%% Theory of Computing
%% File: v021a010.tex
%% Publication date: November 26, 2025
%% Authors: Subhash Khot, Dor Minzer, and Muli Safra
%
%
%
%
%

%
%

%
%
%
%
%
%
%
%
%
%

%

%
%
%
%
%
%
%
%
%
%
%
%
%
%
%
%
%
%
%
%
%
%

%

\documentclass{toc}

%
%
%
\tocdetails{%
title = {On Independent Sets, \texorpdfstring{$2$-to-$2$}{2-to-2} Games and Grassmann Graphs},
%
%
%
  number_of_pages = {53},
  number_of_bibitems = {43},
  number_of_figures = {0},
  conference_version = {STOC17}, %
    %
    %
%
%
%
  author = {Subhash Khot, Dor Minzer, and Muli Safra},
    %
    %
    %
    %
    %
  authorlist = {Subhash Khot, Dor Minzer, Muli Safra},
    %
    %
    %
    %
    %
    %
    %
%
%
  %
  %
  %
  %
%
  acmclassification = {F.2.3},
    %
    %
  amsclassification = {68Q17},
    %
    %
%
%
keywords = {probabilistically checkable proofs, Unique Games Conjecture, Grassmann graph},
    %
    %
    %
    %
    %
}   %

    %
    %
    %
    %
    %
    %
    %
    %
    %
    %
    %
    %
    %
    %
    %
    %
    %
    %
    %
    %
    %
    %
    %
    %
    %
    %
    %
    %
    %

%
%
\tocdetails{%
volume = 21,      %
number = 10,       %
year = 2025,
%
%
%
%
%
%
%
%
%
%
%
%
received = {December 8, 2020}, %
revised = {February 27, 2024},
published = {November 26, 2025},
%
%
%
doi = {10.4086/toc.2025.v021a010},
       %
}   %

%
%
%
%
%
%

%
%
%

\newcommand\gap{ \textsf{gap-} }
\newcommand\gapV{ \textsf{gap$_V$-} }
%
%
   \newtheorem{conj}[theorem]{Hypothesis}
   \newtheorem{rconj}[theorem]{Conjecture}
%
%
%
%
%
%
%
%
\newcommand\comment[1]{\fbox{\tt #1}}
\newcommand\E{\mathbb{E}}
\newcommand\card[1]{\left| {#1} \right|}
\newcommand\tup[1]{\left\langle #1 \right\rangle}
\newcommand\sett[2]{\left\{ \left. #1 \;\right\vert #2 \right\}}
\newcommand\settright[2]{\left\{ #1 \;\left\vert\; #2 \right.\right\}}
\newcommand\set[1]{{\left\{ #1 \right\}}}
\newcommand\Prob[2]{{\Pr_{#1}\left[ {#2} \right]}}
\newcommand\cProb[3]{{\Pr_{#1}\left[ \left. #3 \;\right\vert #2 \right]}}
\newcommand\cProbright[3]{{\Pr_{#1}\left[ #3\;\left\vert  #2\right. \right]}}
\newcommand\cExpect[3]{{\mathbb{E}_{#1}\left[ \left. #3 \;\right\vert #2 \right]}}
\newcommand\Expc[2]{{{\bf E}_{#1}\left[ {#2} \right]}}
\newcommand\ang[1]{< {#1}>}
\newcommand\norm[1]{\| #1 \|}
\newcommand\power[1]{\set{0,1}^{#1}}
\newcommand\ceil[1]{\lceil{#1}\rceil}
\newcommand\round[1]{\left\lfloor{#1}\right\rfloor}
\newcommand\half{{1\over2}}
\newcommand\quarter{{1\over4}}
\newcommand\eigth{{1\over8}}
\newcommand\defeq{\stackrel{def}{=}}
\newcommand\skipi{{\vskip 10pt}}
\newcommand\IP[2]{\left\langle {#1}, {#2} \right\rangle}
\newcommand\TF{ {\set{T, F}} }
\newcommand\pmo{ { \{\pm 1\} } }
\newcommand\Gn{ {\pmo^n} }
\newcommand\BF{ {\colon\Gn\to\pmo } }
\newcommand\Expect[2]{{\mathop{\mathbb{E}}_{#1}\left[ {#2} \right]}}
\newcommand\Set[2]{ \left\{ #1 \vert #2\right\} }
\newcommand\nV{{\cal \#V}}
\newcommand\pV{{ p{\cal V} }}
\newcommand\tV{ {\cal \widetilde{V} } }
\newcommand\sqrtn{ {\sqrt n} }
\newcommand\cube[1]{ { \set{-1,1}^{#1} } }
\newcommand\minimum[2] {  \min\left\{{{#1}, {#2}}\right\}   }
\newcommand\spn{{\sf Span}}
\newcommand\spnn[1]{{\sf Span}({#1})}
\newcommand\inner[2]{\langle{#1},{#2}\rangle}
\newcommand\eps{\varepsilon}
\newcommand\Ext{{\sf s-Ext}}
\newcommand\Extt{{\sf Ext}}
\newcommand\sgn[1]{sign({#1})}
\newcommand*\xor{\mathbin{\oplus}}
\newcommand*\bxor{\mathbin{\bigoplus}}
\newcommand*\Hell[2]{{H\left({#1}, \ {#2}\right)}}
\newcommand*\cHell[4]{{H\left({#1}\vert({#2}),\ {#3}\vert({#4}) \right)}}
\newcommand*\QED{}
\newcommand*\Uni{{\cal U}}
\newcommand*\aconj[1]{\alpha{#1}}
\newcommand*\dimm{{\sf dim}}
\renewcommand\geq{\geqslant}
\renewcommand\leq{\leqslant}

\begin{document}

\begin{frontmatter}%

%
%
%

%
%
\title{On Independent Sets, $2$-to-$2$ Games and Grassmann
  Graphs\titlefootnote{A preliminary version of this paper appeared
    in the \href{https://doi.org/10.1145/3055399.3055432}{Proceedings
      of the 49th ACM Symposium on Theory of Computing
    (STOC'17)}~\cite{conf-version}.}
}  %

%
%

%
%
%
%
%
%

%
%
%
%
%
%
%
%
%
%
%
%


%
%
%
%

%

\author[khot]{Subhash Khot\thanks{Research
supported by NSF CCF-1422159, Simons Collaboration on Algorithms and Geometry, and Simons Investigator Award.}}
\author[minzer]{Dor Minzer\thanks{This work was completed while the author was in the School of Computer Science, Tel Aviv University
and was supported by Clore scholarship.}}
\author[safra]{Muli Safra\thanks{Supported by the European Research Council (ERC) under the European Union’s Horizon 2020 research and innovation programme (Grant agreement No. 835152),
by an ISF grant 2017/13 and by a BSF grant 2016414.}}

%
%
%
%

%
%
%
%
%
%
%
%
\begin{abstract}
  We present a candidate reduction from the  $3$-Lin problem to the $2$-to-$2$ Games problem and present a combinatorial hypothesis about
Grassmann graphs which, if correct, is sufficient to show the soundness of the reduction in
a certain non-standard sense. A reduction that is sound in this non-standard sense
implies that it is NP-hard to distinguish whether an $n$-vertex graph has an independent
set of size $\left( 1- \frac{1}{\sqrt{2}} \right) n - o(n) $ or
%
every  independent set has size $o(n)$, and consequently, that it is NP-hard
to approximate the  Vertex Cover problem within a factor $\sqrt{2}-o(1)$.

This %
article %
initiates and serves as the first installment in a line of work
%
%
%
%
by various subsets of the authors together with Dinur and Kindler
(with additional contributions by Barak, Kothari, and Steurer (ITCS'19)) %
which led to a proof of the $2$-to-$2$ Games Conjecture (albeit with imperfect
completeness), which in %
particular  %
implies the NP-hardness results for independent set
and vertex cover mentioned above.
\end{abstract}

\iffalse %
%
%
%
%
%
%
%
%
%
%
%
%
%
%
%
%
%
%
%
%
%

\tocarxivcategory{!!!}

\fi %

\end{frontmatter}

%
%
%
\section{Introduction}

This paper %
focuses  %
on hardness-of-approximation results for the Vertex Cover and the Independent Set problems, which are closely related to
the hardness of approximating the $2$-to-$2$ Games problem with a certain non-standard notion of soundness.
This %
article  %
represents the starting point of a successful line of attack on the $2$-to-$2$ Games Conjecture %
\cite{conf-version,DKKMS1,DKKMS2,KMS2} %
(with additional contributions from \cite{KS,KMMS,BKS}),
which culminated in a proof of the conjecture (albeit with imperfect completeness) and its implications. These implications in particular include the improved hardness results
for independent set and vertex cover claimed in this paper. We elaborate on these subsequent developments in \expref{Section}{sec:subsequent}.


\subsection{Vertex Cover and Independent Set}
Given an $n$-vertex  graph
$G=(V,E)$, the Vertex Cover problem asks for a
vertex cover of minimum size, namely, a subset $C \subseteq V$ of minimum size that includes
  at least one  endpoint of each edge $e \in E$. This is a classic NP-hard problem and
 has a greedy $2$-approximation algorithm. The algorithm starts with the graph $G$, initializes $C = \emptyset$, and
 until the graph has at least one edge remaining, picks an edge,
adds both its endpoints to $C$, removes all edges incident on either of these two endpoints, and repeats. It is easily seen
that the final set $C$ is a vertex cover of $G$ and has size at most twice that of the minimum vertex cover.  A somewhat better
approximation algorithm achieving factor $2- \Omega\left(\frac{1}{\sqrt{\log n}}\right)$  is known via SDP relaxation \cite{Hal02,Kara09}. However it is a
major open question whether there is a $2-\delta$ approximation algorithm for some fixed positive
constant $\delta$. Surprisingly, as discussed below, there is some evidence to the
contrary: Vertex Cover might actually be hard to approximate
 within a factor $2-\eps$ for every positive constant $\eps$.


The complement $V \setminus C$  of a vertex-cover $C$ is an independent set,
namely, a set of vertices $I \subseteq V$ that has no edge inside it. For constants $0 < \beta < \alpha <1$,
let ${\sf GapIS}(\alpha, \beta)$ denote
the promise gap-problem where the task is to distinguish whether a given $n$-vertex graph has an independent set
of size at least $\alpha n$ or whether every independent set is of size at most $\beta n$. Clearly, if
${\sf GapIS}(\alpha, \beta)$ is hard,\footnote{The statement that ${\sf GapIS}(\alpha, \beta)$ is NP-hard is
equivalent to the statement that NP has a Probabilistically Checkable Proof (PCP) that has zero-free bit complexity,
has completeness at least $\alpha$ and has soundness at most $\beta$, see \cite[Proposition 5.6, Theorem 8.2]{BGS}. We avoid
 the terminology of free bit complexity in this paper.}
 then it would be hard to approximate Vertex Cover  within a factor
strictly less than $\frac{1-\beta}{1-\alpha}$.

Let $\eps$ denote a positive and arbitrarily small constant. We summarize the known
 NP-hardness results for approximating the Vertex Cover problem, obtained in a sequence of
 highly influential papers.
 Building on the
PCP Theorem \cite{FGLSS96, AS98, ALMSS98}, the Parallel Repetition Theorem \cite{Raz98}, and the
Long Code based PCP framework in \cite{BGS}, H\aa stad \cite{Has99} showed that ${\sf GapIS}(\frac{1}{4}-\eps,\frac{1}{8}+\eps)$
is NP-hard, implying $\frac{7}{6}-\eps \approx 1.16$ hardness factor for Vertex Cover.
Dinur and Safra
\cite{DS05}  showed that ${\sf GapIS}(p-\eps, 4p^3-3p^4+\eps)$ is NP-hard for $p = \frac{3-\sqrt{5}}{2}$, implying
$10 \sqrt{5} - 21 - \eps \approx 1.36$ hardness factor for Vertex Cover. Their paper introduced several techniques, \eg, the
Biased Long Code, application of Fourier analytic theorems on Boolean hypercube, and implicitly, the notion of  $2$-to-$2$ Games,
all of which are indispensable in the authors' opinion, for further progress on Vertex Cover.


\subsection{\texorpdfstring{$d$-to-$d$}{s-to-d} games}


In this section, we discuss the $d$-to-$d$ Games and their connection to the Independent Set and
Vertex Cover problems.


\begin{definition}\label{def:game}
A $2$-Prover-$1$-Round Game $G=(V,E,\Sigma,\Phi)$ consists of a set of variables $V$, a set of directed edges
$E\subseteq V\times V$, a set of colors $\Sigma$, and
a constraint $\Phi(u,v)\subseteq \Sigma\times \Sigma$ for every (directed) edge $(u,v) \in E$. The goal is to assign colors to variables, say
 $A: V \to \Sigma$, so as to satisfy the maximum number of the constraints. A constraint $\Phi(u,v)$
is satisfied if $ (A(u), A(v)) \in \Phi(u,v)$, where by abuse of notation, $\Phi(u,v) \subseteq \Sigma
\times \Sigma$ denotes the subset of color-pairs that are deemed satisfactory. The subset  $\Phi(u,v) \subseteq \Sigma \times \Sigma$
may in general depend on the edge $(u,v)$.

Let $d \geqslant 1$ be an integer. A constraint $\Phi(u,v) \subseteq \Sigma \times \Sigma$ is said to be a $d$-to-$d$ constraint if
there are partitions
$A_1,...,A_r$ and  $B_1,...,B_r$ of $\Sigma$ into sets of size $d$  such that ($|\Sigma| = rd$)
\[
\Phi(u,v)=\bigcup\limits_{i=1}^r{(A_i\times B_i)}\,.%
\]
A $2$-Prover-$1$-Round Game $G=(V,E,\Phi,\Sigma)$ is said to be a $d$-to-$d$ Game if every constraint $\Phi(u,v)$ is a $d$-to-$d$ constraint.
A $1$-to-$1$ Game is also called a Unique Game. In this case, $\Phi(u,v)$ is simply a perfect matching on $\Sigma \times \Sigma$.
%
%
\end{definition}

In the above definitions, the number of colors $|\Sigma|$ is thought of as a constant,
possibly large, and the
size of the constraint graph as the growing input size.
Motivated by a hardness-of-approximation result for the $2$-SAT problem,
Khot~\cite{Khot02} formulated the %
following conjecture. %

%
%
%
\begin{rconj}[Unique Games Conjecture]\label{conj:ugc}
For every  %
$\delta>0$, for sufficiently large  %
$|\Sigma|$,
given an instance $G=(V,E,\Phi,\Sigma)$ of a Unique Game, it is NP-hard
to distinguish between
\begin{itemize}
\item YES case: there is a coloring satisfying
a %
$1-\delta$ fraction of the constraints of $G$.
\item NO case: no coloring satisfies more than
a %
$\delta$ fraction of the constraints of $G$.
\end{itemize}
\end{rconj}

The reduction in \cite{DS05} implicitly suggests the idea of $2$-to-$2$ Games (though
therein, the game is a $\alpha$-game in the sense
of \cite{DinurMosselRegev} instead of a $2$-to-$2$ game and
the notion of soundness is non-standard). Motivated by hardness-of-approximation result for the Vertex Cover
problem,  Khot~\cite{Khot02} also formulated the $d$-to-$d$ Conjecture.\footnote{Note that the Unique Games Conjecture is, necessarily,
made with imperfect completeness whereas the $d$-to-$d$ Conjecture is made with perfect completeness. Strictly speaking, the conjecture
in~\cite{Khot02} is a $d$-to-$1$ Conjecture. It implies (and in the authors' opinion,
is morally equivalent to) the $d$-to-$d$ Conjecture stated here.}

%
\begin{rconj}[$d$-to-$d$ Conjecture] \label{conj:d-to-d}
Fix any integer $d \geqslant 2$.
For every %
$\delta>0$, for sufficiently large
%
$|\Sigma|$,  given an instance $G=(V,E,\Phi,\Sigma)$ of a $d$-to-$d$ Game, it is NP-hard
to distinguish between
\begin{itemize}
\item YES case: there is a coloring satisfying all of the constraints of $G$.
\item NO case: no coloring satisfies more than
a %
$\delta$ fraction of the constraints of $G$.
\end{itemize}
\end{rconj}

Let $\eps$ denote a positive and arbitrarily small constant.
It was shown that the $d$-to-$d$ Conjecture implies that ${\sf GapIS} \left(1-\frac{1}{2^{1/d}}-\eps, \eps\right)$ is NP-hard for $d \geqslant 2$ \cite{Khot02}. The result did not
apply in  case of Unique Games,
because the imperfect completeness of Unique Games presented a difficulty.
This difficulty was circumvented in \cite{KR08} where the
authors showed that the Unique Games Conjecture implies that ${\sf GapIS} (\frac{1}{2}-\eps, \eps)$ is NP-hard and therefore, implies that Vertex Cover is NP-hard to approximate
 within a factor $2-\eps$.
The Unique Games Conjecture, and to a lesser extent the $d$-to-$d$ Conjecture, is now a prominent open question in theoretical computer science.
It implies hardness-of-approximation results, often optimal results, for numerous problems and has connections to several areas in algorithms, computational complexity,
and geometry, see \cite{Trevisan-survey, KhotUGC-survey, KhotICM-survey} for surveys on the topic.  It is thus worthwhile to investigate possible lines of attack towards
proving (or disproving) the Unique Games Conjecture, the $d$-to-$d$ Conjectures,
and their variants. In this paper, we
present a line of attack towards proving a variant of the
 $2$-to-$2$  Conjecture with a certain non-standard notion of soundness,
and towards making progress
on the Independent Set and Vertex Cover problems.



Unfortunately we have to consider games where the constraints are a mix of $2$-to-$2$ constraints and $1$-to-$1$ constraints and
the
game satisfies an additional transitivity property. This feature might not be necessary, but we are unable to circumvent it for now
(see \expref{Remark}{remark:transitivity_why}).

\begin{definition}\label{def:trans_game}
A Transitive $2$-to-$2$ Game is a game $G=(V,E,\Phi,\Sigma)$ where
\begin{itemize}
\item Each constraint $\Phi(u,v)$ is a $2$-to-$2$ constraint or a $1$-to-$1$ constraint.
\item Transitivity: If there is a $1$-to-$1$ constraint $\Phi(u,v)$ and a constraint $\Phi(v,w)$, then there is also a constraint $\Phi(u,w)$. The constraint
      $\Phi(u,w)$ is either
  $1$-to-$1$ or $2$-to-$2$ depending on whether $\Phi(v,w)$ is $1$-to-$1$ or
  $2$-to-$2$. %
%

  Moreover, the constraint $\Phi(u,w)$
is a composition of constraints $\Phi(u,v)$ and $\Phi(v,w)$, \ie, for every $a,b,c \in \Sigma$,
$$ (a,b) \in \Phi(u,v), \ (b,c) \in  \Phi(v,w) \implies (a,c) \in \Phi(u,w)\,. $$%
\end{itemize}
\end{definition}

The notion of soundness (NO case) in \expref{Conjectures}{conj:ugc} \expref{and}{conj:d-to-d} states that no coloring satisfies more than a tiny fraction of
the constraints. This notion will be referred to as the standard notion of soundness.
It has been a folklore among the experts (see \cite[Theorem 3.1]{KR08},
where this is stated for the $1$-to-$1$ case) that as far as the Independent Set and Vertex Cover
hardness results are concerned, a non-standard notion of soundness for the $d$-to-$d$ Games suffices. The non-standard notion concerns
``$(j,\delta)$-colorings'' that we define next. \footnote{Our definition has to take into account the transitivity feature that we unfortunately
have to deal with. Also, we restrict ourselves to the case $d=2$ which is our primary concern.}

\begin{definition}\label{def:j,delta}
Let $G(V,E,\Phi,\Sigma)$ be a Transitive $2$-to-$2$ Game, $\delta>0$, $j$ be a positive integer, and $X\subseteq V$.
A coloring $A \colon X\to {\Sigma \choose j}$ is called a $(j,\delta)$-coloring if the following holds
(note that one is allowed to assign a set of $j$ colors to every variable in $X$ and the rest of the variables are
unassigned):
\begin{itemize}
\item $\card{X}\geqslant \delta \card{V}$.
\item For every $u,v\in X$ such that $\Phi(u,v)$ is a $2$-to-$2$ constraint, there are colors $a \in A(u)$, $b \in A(v)$ such that
$(a,b)\in \Phi(u,v)$.
\item For every $u,v\in X$ such that $\Phi(u,v)$ is a $1$-to-$1$ constraint, the color sets $A(u)$ and $A(v)$ are identical up to
the matching $\Phi(u,v)$. More precisely, for every $(a,b) \in \Phi(u,v)$, $a \in A(u)$ if and only if $b \in A(v)$.
\end{itemize}
\end{definition}

Now we state a variant of \expref{Conjecture}{conj:d-to-d} (for $d=2$) with a non-standard notion of soundness, imperfect
completeness, and for transitive $2$-to-$2$ games. This variant is to be thought of as weaker than \expref{Conjecture}{conj:d-to-d}
in the sense that \expref{Conjecture}{conj:d-to-d} (for $d=2$) implies it (up to an insignificant caveat).
\footnote{
The insignificant caveats are twofold. First, the
%
``YES case'' %
in \expref{Conjecture}{conj:j-delta} is a bit stronger
than the yes case in \expref{Conjecture}{conj:d-to-d} with imperfect completeness, since it requires $G$ to have a subgraph
of density $1-\delta$ wherein all constraints can be satisfied. Second, in the ``NO case'' we would need the graph $G$
to satisfy mild connectivity properties, in the sense that a set of density $\delta$ would contain at least
$\Omega(\delta^2)$ fraction of edges. In that case if an instance $G$ has a $(j,\delta)$-coloring, then there is an assignment satisfying at least
%
an $\Omega(\delta^2/j^2)$ fraction  %
of constraints which is the result of choosing a color
randomly for each vertex from its $j$-list.}
\begin{rconj}\label{conj:j-delta}
%
%
%
%
%
For every $\delta>0$ and every positive integer $j$,
for all sufficiently large $|\Sigma|$,  %
given an instance $G=(V,E,\Phi,\Sigma)$ of a Transitive $2$-to-$2$ Game,
it is NP-hard to distinguish between
\begin{itemize}
\item YES case: there is a  $(1,1-\delta)$-coloring to the graph $G$.
\item NO case: there is no $(j,\delta)$-coloring to the graph $G$.
\end{itemize}
\end{rconj}

Finally, we note that the result below follows directly from prior
%
work  %
\cite{DS05, Khot02, KR08}.
A proof is presented in \expref{Section}{sec:vc-is} for the sake of completeness.
%
The ingredients include the Biased Long Code
and analytic theorems of Russo, Margulis and Friedgut on the Boolean hypercube.
Some care is required to handle the transitivity feature.

\begin{theorem}\label{thm:j-delta-vc}
If \expref{Conjecture}{conj:j-delta} holds, then
${\sf GapIS}(1-\frac{1}{\sqrt{2}}-\eps, \eps)$ is NP-hard for every positive
constant~$\eps$.  %
%
\end{theorem}



\subsection{Our results}

Roughly speaking, we give a reduction from an NP-hard problem known as $3$-Lin to
(Transitive) $2$-to-$2$ Game such that
the reduction is sound in the sense of \expref{Conjecture}{conj:j-delta} assuming a
combinatorial hypothesis. Therefore, correctness of the combinatorial hypothesis would
imply \expref{Conjecture}{conj:j-delta}
and the corresponding results for ${\sf GapIS}$ and Vertex Cover via
\expref{Theorem}{thm:j-delta-vc}.


We now state the results more formally. Let $3$-Lin be the following problem. The instance of the problem is
$(X, {\sf Eq})$ where $X$ is a set of variables taking values over ${\mathbb F}_2$ and ${\sf Eq}$ is a set of linear equations over
${\mathbb F}_2$ such that every equation depends on three variables in $X$. The goal is to find an assignment to the variables
so as to maximize the number of equations satisfied. Let ${\sf Gap3Lin}(c, s)$ denote the promise gap-problem
where the task is to distinguish whether a given $3$-Lin instance has an assignment satisfying at least $c$ fraction of
the equations or whether every assignment satisfies at most $s$ fraction of the equations.
A celebrated result of H\aa stad \cite{Has01} shows that for every  $\eps>0$, %
 ${\sf Gap3Lin}(1-\eps, \frac{1}{2}+\eps)$ is NP-hard. For our purposes, it is convenient to work with a $3$-Lin
 instance that is regular, \ie, every equation contains three distinct variables,  every variable appears in exactly, say $5$,
 equations, and two distinct equations share at most one variable.
 Starting with H\aa stad's result, it is a routine exercise to show that ${\sf Gap3Lin}(1-\eps, s^*)$ is NP-hard on regular instances for every positive
 constant $\eps$ and for some absolute constant $s^* < 1$.
 Our main result is this. %

\begin{theorem}\label{thm:main} For every $s^{*}\in (1/2,1)$, every positive integer $j$ and every
$\delta>0$, for sufficiently small $\eps>0$,
 %
there is a
 polynomial time reduction mapping a regular instance
$(X, {\sf Eq})$ of ${\sf Gap3Lin}(1-\eps, s^*)$
to an instance $G=(V,E,\Phi,\Sigma)$ of Transitive $2$-to-$2$ Game
such that:
\begin{itemize}
\item
YES case: If there is an assignment satisfying at least
a  %
$1-\eps$ fraction of the equations
in $(X, {\sf Eq})$, then
there is a $(1,1-\delta)$-coloring to $G$.

\item NO case: Assuming the combinatorial \expref{Hypothesis}{conj:main2},
if no assignment satisfies more than $s^*$ fraction of the equations in $(X, {\sf Eq})$,
then there is no $(j,\delta)$-coloring to $G$.
\end{itemize}
\end{theorem}

We present the combinatorial hypothesis later, after discussing the Grassmann graph and the motivation behind the hypothesis.
The following corollary follows via \expref{Theorem}{thm:j-delta-vc}.
\begin{corollary}\label{thm:main2} Assuming the combinatorial \expref{Hypothesis}{conj:main2},
\begin{itemize}
\item \expref{Conjecture}{conj:j-delta} is correct.
\item ${\sf GapIS}\left(1-\frac{1}{\sqrt{2}}-\eps, \eps\right)$ is NP-hard for every $\eps>0$.  %
%
\item Vertex Cover is NP-hard to approximate to within a factor $\sqrt{2}-\eps$ for every $\eps>0$.%
\end{itemize}
\end{corollary}



\begin{remark} Our reduction, depending on the correctness of the combinatorial hypothesis, would give $\sqrt{2} - o(1)$ hardness
for Vertex Cover, improving on the $1.36$ hardness of Dinur and Safra. While the numerical improvement would be interesting, in the
authors' opinion, a much more interesting feature would be the ``gap-location'' for the Independent Set problem. Our
reduction would show that ${\sf GapIS}(\alpha^*, \beta)$ is NP-hard where $\alpha^*$ is a fixed, absolute constant and $\beta \to 0$ is
an arbitrarily small constant. Such a result would be remarkable, in authors' opinion, irrespective of whether it gives an
improvement in the Vertex Cover hardness factor. The best known result in this direction is that ${\sf GapIS}(2^{-k}- o(1),
 ~2^{-2^k+1}+o(1))$ is NP-hard for every integer $k \geqslant 2$, by
Siu On Chan \cite{Chan13}. Hardness of ~${\sf GapIS}(\alpha, \beta)$ corresponds to Vertex Cover hardness of $\frac{1-\beta}{1-\alpha}$.
An improvement in Vertex Cover hardness would not necessarily yield $\beta \to 0$ while keeping $\alpha$ fixed, which in authors'
opinion, is a more fundamental and challenging question.
\end{remark}



\begin{remark}
$3$-Lin is known to have a %
Lasserre %
integrality gap %
on random instances  with perfect completeness \cite{Grigoriev, Schoen08, Tul09}.
Our reduction from $3$-Lin to $2$-to-$2$ Games and then the reduction from $2$-to-$2$ Games to the Independent Set and
Vertex Cover problems could yield similar Lasserre integrality gap for the latter problems.
%
%
%
\end{remark}



\subsection{Overview of the reduction}

%
The vast majority %
of hardness-of-approximation results are proved by constructing special purpose Probabilistically Checkable Proof Systems (PCPs) (\eg,
\cite{ALMSS98,BGS,Has99,Has01,GHS00,Khot-quasiPCP,Chan13}). Sometimes it is more convenient, and certainly helpful
to a reader not familiar with PCP terminology, to take a combinatorial view and present a PCP construction, equivalently, as a combinatorial reduction (\eg,
\cite{DS05, Dinur07}). In this paper, we adopt the latter view as far as possible, using PCP terminology wherever helpful or necessary.

%
A generic and extremely successful framework to construct PCPs, developed in \cite{ABSS,BGS,Raz98, Has99, Has01}, is as follows.
It consists of two modules, known as an \emph{Inner PCP} and an \emph{Outer PCP}, which are then composed together.  %
%
\begin{itemize}
\item The Inner PCP is best thought of as a combinatorial gadget
combined with an analysis of its structural properties. The gadget is often coding-theoretic and
 amounts to a specific encoding scheme and a probabilistic procedure to test whether a given
word is (close to) a codeword and if so, to decode (or ``list-decode'') that codeword. The choice of the encoding scheme
as well as the nature of the
tester (\eg, number of queries and the acceptance predicate) are dictated by the target problem for which one desires a hardness-of-approximation result.
\item The Outer PCP is a canonical NP-hard problem known variously as $2$-Prover-$1$-Round Game, $2$-CSP, or Label Cover. The problem is known to be
      very hard to approximate \cite{ABSS,Raz98}, via Raz's Parallel Repetition Theorem.
\item The composition amounts to taking several (local) copies of the Inner PCP gadget and
combining %
them via the (global)
      Outer PCP.\footnote{Perhaps a useful analogy here is the text-book reduction from 3SAT to the Traveling Salesperson. Therein, for every variable
      of the 3SAT instance, there is a copy of a fixed graph (= TSP-gadget) and then edges are added between these copies using the clauses of the 3SAT instance.}
\end{itemize}
With this framework in mind, we give a short and informal overview of our reduction, leaving
out several intricate details.
We recall that the reduction is
intended to construct a $2$-to-$2$ Game, where the goal is to assign ``colors'' to vertices
of a graph and once a color %
has been  %
assigned to
a vertex, each one of its neighbors has exactly $2$ colors
 %
%
%
%
that are deemed acceptable.
This consideration dictates our choice of Inner PCP and
specifically,
what we might call \emph{Grassmann encoding}, %
%
wherein the encoded object is a linear function on
an %
%
${\mathbb F}_2$-vector space.
As such, linearity is inherent to
our reduction which dictates, in turn, our choice of Outer PCP as a $2$-Prover-$1$-Round
Game arising from $3$-Lin instances, making
%
%
${\sf Gap3Lin}$ a natural starting point for the reduction.

%
\subsubsection{Inner PCP: Grassmann graph, Grassmann encoding, and our hypotheses}

Let $1 \ll \ell \ll n$ be integers. The vertex set ${\cal L}, ~|{\cal L}|=N$ 
 of the Grassmann graph $G(\power{n}, \ell)$ consists of all $\ell$-dimensional subspaces $L$ of $\power{n}$, the $n$-dimensional vector space over ${\mathbb F}_2$.
 A pair of vertices, $L, L' \in {\cal L}$, %
%
are %
adjacent  %
if and only if ${\dim}(L \cap L')= \ell-1$. Given a linear function $f: \power{n}\to\power{}$ (or equivalently
an $n$-bit string $s_f$ that defines the linear function $x \mapsto \left< s_f, x \right>$), the Grassmann graph leads to a natural encoding of $f$  by a string of
length $N$ over the alphabet $\Sigma =
\{1,2,\ldots, 2^\ell\}$.  The encoding writes down, for every $\ell$-dimensional
subspace $L$, the linear function $f|_L$, namely, the restriction of $f$ to the subspace $L$. There are exactly $2^\ell$ distinct linear functions on an $\ell$-dimensional
space which can be numbered with $\Sigma = \{1,2,\ldots,2^\ell\}$.

Now suppose that the (global) linear function $f$ is unknown, but for an edge $(L, L')$ in the Grassmann graph, $f|_L = \sigma \in \Sigma$ is known. What do we know
about $f|_{L'} = \sigma'$?  We note that ${\sf dim}(L \cap L')=\ell-1$ and since $\sigma, \sigma'$ are (supposed to be)
restrictions of the same global function, it must be the case that they are consistent on $L\cap L'$, \ie,
$\sigma|_{L\cap L'} = {\sigma'}|_{L\cap L'}$. Clearly, for a given $\sigma$, there
are exactly two possible choices for $\sigma'$.
In fact, the
%
consistent pairs %
%
of functions $(\sigma, \sigma')$ on $(L, L')$
%
%
%
%
%
%
%
%
%
%
%
%
%
are in $2$-to-$2$ correspondence with each other.
Let $\Phi(L,L') \subseteq \Sigma \times \Sigma$ denote this set of consistent pairs.

We are naturally led to the following $2$-to-$2$ Game: assign ``colors'' from $\Sigma$ (interpreted as linear functions on $\ell$-spaces) to the vertices
 of the Grassmann graph, and be consistent on a significant fraction of the
 constraints $\Phi(L,L')$.
 %
 Of course, one option is to pick a global linear
 function $f$ and assign $L \mapsto f|_L$; such strategy yields consistency on all edges.
 Is this essentially the only strategy? Before proceeding, let us mention that
 there are two notions of consistency that are natural and relevant:
 \begin{itemize}
 \item (Standard Consistency): An assignment $A: {\cal L} \to \Sigma$ is said to be $\delta$-consistent if it is consistent on
 a %
 $\delta$ fraction of the
       edges, \ie, for
       a %
       $\delta$ fraction of the edges $(L, L')$, $(A[L], A[L']) \in \Phi(L, L')$.
 \item (Non-Standard Consistency): An assignment $F: {\cal L} \to {\Sigma \choose j} \cup \{\emptyset\}$ (\ie, every vertex either gets $j$ colors or
        does not get any color) is said to be $(j,\delta)$-consistent if there is a subset ${\cal S} \subseteq {\cal L}, ~|{\cal S}| \geq \delta |{\cal L}|$
        such that (a) for all $ L \in {\cal S}, ~F[L] \neq  \emptyset$ (b) for all edges $(L, L')$ inside ${\cal S}$, there are colors $\sigma\in F[L], \sigma'\in F[L']$
        with $(\sigma, \sigma') \in \Phi(L,L')$.
\end{itemize}

One is tempted to speculate %
%
%
as follows. (The parameters $\ell, n$ are thought of as arbitrarily
large with $\ell \ll n$.)  %
\medskip
\noindent {\bf Speculation (1):} For every $\delta$, there is $\delta'$ such that given any $\delta$-consistent (in the standard sense)
   assignment $A$, there is a global linear
function $f:\power{n} \to \power{}$ such that for $\delta'$ fraction of the vertices $L$, it holds that $A[L] = f|_L$.

\medskip
\noindent {\bf Speculation (2):} For every $j \geq 1, \delta > 0$, there is $\delta' = \delta'(\delta,j)>0$, such that given any $(j,\delta)$-consistent (in the non-standard sense)
assignment $F$, there is a global linear
function $f:\power{n} \to \power{}$ such that for $\delta'$ fraction of the vertices $L$, it holds that $f|_L \in F[L]$.

\medskip In coding-theoretic terms,
in both the speculations, the assignments $A$ or $F$ are regarded as ``received words.''
If the desired global linear function $f$ exists, it then serves as a ``decoding'' of the received word.
We however know that Speculation (2) is false in the case $j \geq 3$. A counterexample appears in \expref{Section}{section:grassmann}. Since
Speculation (1) implies Speculation (2),\footnote{The Grassmann graph is
dense in the sense that a set %
%
consisting of a delta fraction of the vertices %
contains at least
a   %
$\delta^2$ fraction of the edges. Given a $(j,\delta)$-assignment,
one can pick a random assignment from its $j$-list for every vertex that has been assigned
and satisfy $\frac{\delta^2}{j}$ fraction of the edges in
expectation. Hence
the %
existence of
a $(j,\delta)$-consistent assignment implies
the %
existence of a $\frac{\delta^2}{j}$-consistent assignment.} Speculation (1) is also false.
We believe that Speculation (2) is correct in the case $j=1$ and
seems to present interesting challenges. We show, in \expref{Section}{section:j_1}, that
 it follows from our \expref{Hypothesis}{conj:gras_expand}, via
%
an  %
 ``$\ell$-space vs $b$-space''
linearity test.\footnote{This test is in the spirit of ``line vs point'' and
``plane vs plane'' low-degree test in \cite{RS96,AS98,RS97}.
 However our analysis is Fourier-based instead of algebraic and combinatorial.}
The linearity test and its analysis using Fourier-analytic methods are presented in \expref{Section}{appendix:linearity_testing}.
\expref{Hypothesis}{conj:gras_expand} states that in the Grassmann graph, a set of
constant density contains a connected component of constant density inside it.
In addition, we propose \expref{Hypothesis}{conj:v_expansion} stating
that the Grassmann graph is a small-set vertex expander
(see \expref{Def.}{def:small-set}).%
%
\footnote{We remark that the Grassmann graph $G(\power{n},\ell)$
is not a small-set edge expander.
Indeed, it has sets of subconstant size with edge expansion $\leq \frac{1}{2}$.
For example, one can fix a non-zero point $x \in \power{n}$ and take
the set of all $\ell$-spaces containing $x$.}
These hypotheses might be a good starting point for further investigation.

We now state our main hypothesis informally. We would like to somehow salvage Speculation (2) for $j>1$. We hypothesize that (see \expref{Hypothesis}{conj:main2} for a formal statement)
given a $(j,\delta)$-assignment to the Grassmann graph $G(\power{n}, \ell)$, there exists a $q$-dimensional subspace $Q$ such that if one
``zooms into'' the subgraph induced on $\ell$-spaces $L$ that contain $Q$ (this subgraph is isomorphic to $G(\power{n-q}, \ell-q)$), then indeed
there is a global linear function that is consistent with the given assignment on $\delta'$ fraction of vertices in the induced subgraph.
Here $q, \delta'$ depend on $j,\delta$. We in fact hypothesize that the zoom-in is successful in this sense for $\alpha$ fraction of $q$-dimensional
subspaces $Q$ where $\alpha$ depends on $j,\delta, \ell$. This hypothesis is sufficient to prove \expref{Conjecture}{conj:j-delta}.
The ``zoom-in'' is a new feature in the context of Inner PCPs and our
Outer PCP needs to have an appropriate mechanism to handle it.


%
\subsubsection{Outer PCP: \texorpdfstring{$2$-Prover-$1$-Round}{2-Prover-1-Round} Game}

We present the Outer PCP as a $2$-Prover-$1$-Round Game. Usually, this game is constructed from a hard instance of 3SAT, \eg, \cite{Has01, Has99},
in which case it is compatible with a Long Code based Inner PCP.
%
However in our case, to be compatible with the Grassmann Code
%
 based Inner PCP, the game needs to be constructed from a hard instance
 of a linear constraint satisfaction problem,
${\sf Gap3Lin}$ being the natural choice (this has been done previously, \eg,
\cite{KhotClique,KS}, with a Hadamard Code
%
based Inner PCP).

Let $(X, {\sf Eq})$ be an instance of ${\sf Gap3Lin}(1-\eps, s^*)$ where $\eps$ can be chosen to be arbitrarily small and $s^* < 1$ is an
absolute constant.
The $2$-Prover-$1$-Round Game is a game between a verifier and two non-communicating provers, where
the provers wish to convince the verifier that the instance $(X, {\sf Eq})$ has a $(1-\eps)$-satisfying assignment. Fix a parameter $k$, thought of as a
large integer, and a ``smoothness'' parameter $\beta$, say $\beta = k^{-\frac{3}{4}}$ for the sake of concreteness. The game proceeds as follows. %
%
\begin{itemize}
\item The verifier picks at random $k$ equations $\{e_1,\ldots, e_k \}$, lets $U$
%
be the set of $3k$ variables that appear in these equations, and
      sends $U$ to the first prover as a question.
\item The verifier picks a subset
%
$V \subseteq U$ of variables  %
by including in $V$, independently for $1 \leq i \leq k$,
     (a) all three variables from the equation $e_i$  with probability $1-\beta$ and (b)
      one of the three
      variables chosen at random from the equation $e_i$ with probability $\beta$. Note that the size of $V$ is, w.h.p., close to its
      expected size $3k-2\beta k$, so $V$ is nearly the same as $U$. The verifier sends $V$ to the second prover as a question.
\item %
%
The first and the second prover  %
      answer with bit-strings $s_U, |s_U|=|U|$ and  $s_V, |s_V|=|V|$, respectively, supposedly giving the assignment
      to the set of variables they received.
\item The verifier accepts if and only if $s_U|_V = s_V$ (\ie, if the two provers agree on the shared variables $V$) and  $s_U$ satisfies
      the $k$ equations (this is known as a \emph{side condition}). %
%
\end{itemize}
The parameter $k$ is a constant, so the size of the game
is polynomial in the size of the ${\sf Gap3Lin}$ instance.
Instead of viewing the game as  ``active'' verification, one can write down the description
of the game as a graph, with possible questions as its vertices and possible question-pairs
asked to the provers as its edges.
The game is then viewed as a ``passive'' optimization problem:
assigning colors (= bit-strings of appropriate length) to the vertices, so as to satisfy constraints on the edges.
The following statements show that approximating the provers' optimal strategy (which, in the passive view, is same as a coloring that maximizes
the fraction of the edge-constraints satisfied)
is a very hard problem, and hence can be
used as a canonical hard problem for further reduction.

\smallskip
\noindent {\bf (Completeness):} It is clear that if the instance $(X, {\sf Eq})$ has a $(1-\eps)$-satisfying assignment, the provers can answer according to this
(global) assignment. The $k$ equations chosen by the verifier are all satisfied with probability $\geq 1-k\eps$, in which case the
verifier accepts.

\smallskip
\noindent {\bf (Soundness):} On the other hand, it follows from the Parallel Repetition Theorem \cite{Raz98, Holenstein, Rao, DS14} that if
every assignment to the instance $(X, {\sf Eq})$ is at most $s^*$-satisfying, then any strategy of the provers can make the verifier
accept with probability at most $2^{-\Omega(\beta k)}$.

%
\subsubsection{Composition of Inner and Outer PCP}

We compose %
the Inner PCP and the Outer PCP, constructing an instance $G_{2:2}$ of a Transitive $2$-to-$2$ Game as in \expref{Definition}{def:trans_game} and \expref{Theorem}{thm:main}.
Only the questions $U$ to the first prover in the Outer PCP appear explicitly in the construction whereas the questions $V$ to the second prover
are only implicitly used. The composition, at a high level, is rather straightforward. However, incorporating the %
%
side conditions  %
from the Outer PCP
and ensuring the transitivity %
of the $2$-to-$2$ Game $G_{2:2}$
present serious difficulties. Both of these issues are skipped altogether from this overview. Also, in the actual reduction,
there are more constraints in $G_{2:2}$ than described here.

In the $2$-to-$2$ Game $G_{2:2}$,
for every question $U$ to the first prover, there is a copy
 of the Grassmann graph $G(\power{U},\ell)$. A vertex $L$
in this graph is to be assigned a color (or a $j$-list of colors)
from the alphabet $\Sigma, |\Sigma|=2^\ell$, the colors being interpreted as linear functions on $L$. The intention is as follows. Suppose that
in the Outer PCP, the prover intends to answer with a bit string $s, |s| =|U|=3k$, a supposed assignment to the variables in $U$. The string $s$ is
thought of as the linear ``inner product'' function $f_s: \power{U} \to \power{}, ~f_s(x) = \left< s, x\right>$. The assignment of colors to the vertices of
the Grassmann graph is then precisely the encoding of the linear function $f_s$, \ie, a vertex $L$ is assigned the color ${f_s}|_L$.
We add the $2$-to-$2$ constraints for all edges $(L, L')$ of this Grassmann graph as in the Inner PCP.

To summarize, the $2$-to-$2$ Game $G_{2:2}$ has a block of vertices for every question $U$ to the first prover  and along with the
edge-constraints inside it, the block is exactly a copy of the Grassmann graph/encoding/game. Now we describe the edges across
two different blocks. Let $U, U'$ be two distinct questions to the first prover and $V$ be a question to the second prover such that
$V \subseteq U$ as well as $V \subseteq U'$, \ie, the verifier can potentially ask the question-pair $(U, V)$ as well as the question-pair $(U', V)$.
Obviously the space $\power{V}$ is contained in both $\power{U}$ and $\power{U'}$. There are two types of edges\footnote{In the actual reduction,
potentially, there are edges between blocks of $U, U'$, even when there is no question $V$ that appears along with both $U, U'$.}
 between the block of
$U$ and  the block of $U'$.
\begin{itemize}
\item For any $L \subseteq \power{V}, {\sf dim}(L)=\ell$, $L$ is contained in
both $U$ and $U'$, and hence there are vertices $u, u'$ in their
blocks corresponding to $L$. We add a $1$-to-$1$ constraint between $u, u'$.
Note that (a) the colorings to the blocks of $U, U'$ are supposed to be the
encodings of the
%
functions $f_{s_U}$ and $f_{s_{U'}}$, respectively
%
(b) the assignments $s_U$ and $s_{U'}$ are supposed to be restrictions of some
global assignment to $U$ and $U'$, respectively %
and hence are supposed to
agree with an assignment $s_V$ on $V$. Therefore, the linear function $f_{s_U}|_L$,
\ie, the intended color of $u$, and the linear function $f_{s_{U'}}|_L$, \ie, the
intended color of $u'$, must be the same, \ie, $f_{s_U}|_L = f_{s_V}|_L = f_{s_{U'}}|_L$.
      This defines the $1$-to-$1$ constraint between $u, u'$.
\item Similarly, for any $L, L' \subseteq \power{V}, {\sf dim}(L)={\sf dim}(L')=\ell, {\sf dim}(L\cap L')=\ell-1$,
 $L$ is contained in $\power{U}$ and $L'$ is contained in $\power{U'}$, and hence there are vertices $u, u'$ in the blocks of $U, U'$ %
%
corresponding to $L, L'$, respectively.  We add a $2$-to-$2$ constraint between $u, u'$. As before,
       the intended color of $u$ is $f_{s_U}|_L = f_{s_V}|_L$ and the intended color of $u'$ is $f_{s_{U'}}|_{L'} = f_{s_V}|_{L'}$. Since
       ${\sf dim}(L \cap L')=\ell-1$, there is a $2$-to-$2$ correspondence between the functions $f_{s_V}|_L$ and $f_{s_V}|_{L'}$,
       which defines the $2$-to-$2$ constraint between $u, u'$.
\end{itemize}
This completes the informal description of the reduction. The actual reduction, with several additional details, is
presented in \expref{Section}{section:construction}. The transitivity of the game $G_{2:2}$ is proved in~\expref{Section}{appendix:transitive}.

%
\subsubsection{Advice, Covering Property, and soundness analysis}

Let us attempt the soundness analysis at a
high level, showing the need for an additional
%
feature in the Outer PCP, which we call advice, as well as a certain
covering property.
%
Given a $(j,\delta)$-coloring to
the game $G_{2:2}$, the soundness analysis derives prover strategies in the Outer PCP with a good success probability.
This implies conversely that if the Outer PCP is chosen beforehand to have low enough soundness ($= 2^{-\Omega(\beta k)}$), then
the game $G_{2:2}$ has no $(j,\delta)$-coloring, completing the proof of \expref{Theorem}{thm:main}.

Accordingly, suppose there is a $(j,\delta)$-coloring to the game $G_{2:2}$ and assume for simplicity
that a $\delta$ fraction of vertices in every block have been ($j$-list-)colored. Fix a question $U$ to the first prover. Towards
deriving her answer, she looks at the $(j,\delta)$-coloring of her copy of the Grassmann graph $G(\power{U},\ell)$.
Our main hypothesis
(\expref{Hypothesis}{conj:main2}) implies that for a good fraction of $q$-dimensional subspaces $Q \subseteq \power{U}$, after zooming into the
subgraph induced on vertices (= $\ell$-spaces) that contain $Q$, there is a global linear function $f = f_s:\power{U} \to \power{}$
that has good agreement with the given coloring. The prover returns the string $s$ as her answer. Though
she does not know which zoom-in space $Q$ works, she can choose $Q$ randomly, and it is hypothesized  to work with a good
probability. We emphasize that the decoded
%
global linear function $f_s$, and hence her answer, in general depends on the choice of the zoom-in space $Q$.

Now let $V$ be the question to the second prover. Since $\power{V}$ is contained in $\power{U}$, the Grassmann graph $G(\power{V},\ell)$
is an induced subgraph of the Grassmann graph $G(\power{U},\ell)$. The second prover wishes to derive his answer from the coloring to his
 graph $G(\power{V},\ell)$. Let $S$ be the subset of vertices in the graph $G(\power{U},\ell)$ that are colored, with density of $S$ being
 $\delta$. The prover can only use the coloring to the set $S \cap G(\power{V},\ell)$, which might potentially have negligible density
  in $G(\power{V},\ell)$. If so, he has no information  to derive his answer from, and in the worst case, this could happen for almost every question
  $V$ asked to him, for a fixed question $U$ to the first prover. The purpose of the
  smoothness
  %
  parameter in the Outer PCP is to
  precisely avoid this issue. Provided that $\beta \sqrt{k} \cdot 2^\ell \to 0$ (which happens with $\beta = k^{-\frac{3}{4}}$ and
  $k$ large enough), $V$ has expected size   $|U| - 2\beta k$ that is close enough to $|U| = 3k$ that
  one has the following guarantee: for a fixed question $U$ to the first prover and for any subset $S$ of density $\delta$ in
  $G(\power{U},\ell)$, for almost every question $V$ to the second prover, the density of $S \cap G(\power{V}, \ell)$ in $G(\power{V},\ell)$
  is $\approx \delta$. This guarantee is referred to as the
  covering property,
  %
  a special case of which is defined and used in \cite{KS}.
  The property is stated formally as \expref{Lemmas}{lemma:covering_normal} \expref{and}{lemma:covering_cond_markov} %
and proved in \expref{Section}{appendix:covering}.

 We assume therefore that for the question $V$ to the second prover,
  a %
 $\delta$ fraction of vertices of his Grassmann graph
 $G(\power{V},\ell)$ are colored. In a similar manner as the first prover, he wishes
to zoom into a $q$-dimensional subspace
 $Q' \subseteq \power{V}$, decode a global linear function $f_{s'}: \power{V} \to \power{}$, return $s'$ as the
 answer, and hope that $s' = s|_V$, \ie, that his answer is consistent with the first prover's answer. Strictly speaking,
 he outputs a short list of all $f_{s'}$ that have agreement with coloring to his Grassmann graph and hopes that one of them
 is consistent with the first prover's function $f_s$. It is reasonable to expect this consistency because both provers are using the same coloring:
 if $f_s$ is consistent with coloring to $S \subseteq G(\power{U},\ell)$, then its restriction $f_s|_V$
  is consistent with $S \cap G(\power{V},\ell)$ for a good fraction of questions $V$ to the second prover, and then
  $f_s|_V = f_{s'}$ appears in the decoded list of the second prover. There is one catch however.
  The decoded global functions depend in general on the zoom-in space, so the two provers
 must ``agree'' on the zoom-in space, \ie, manage to choose $Q=Q'$, without communication. We resolve this issue by letting the verifier
 in the Outer PCP choose a random $q$-dimensional subspace $Q \subseteq \power{V}$ and send it to both provers as extra
 advice
 %
 along with their questions. We make sure that this advice does not compromise or hurt the soundness of the Outer PCP.




%
\subsection{Subsequent work}\label{sec:subsequent}  %
Following %
the conference publication of our work~\cite{conf-version}, %
much progress has been %
made %
in understanding the NP-hardness of $2$-to-$2$ Games and Unique Games.
In \cite{DKKMS1}, together with Dinur and Kindler, we propose a candidate NP-hardness reduction for $2$-to-$2$ Games under the
\emph{standard} notion of soundness. The reduction therein builds on the ideas presented in the current paper,
and its soundness was only proved under a combinatorial hypothesis (in the spirit of \expref{Hypothesis}{conj:main2}).
%
In a subsequent article~\cite{DKKMS2},  %
this combinatorial question is linked to the study of
%
edge expansion  %
in the Grassmann graph;
more concretely, to the structure of small sets of vertices whose edge expansion is
bounded away from $1$.
The authors of~\cite{DKKMS2} pose a concrete hypothesis for the structure of such sets,
validate it in a special regime of parameters,
and argue that it is necessary in order to resolve the combinatorial hypothesis
of~\cite{DKKMS2} (and thus prove that $2$-to-$2$ Games Conjecture holds).
Barak, Kothari and Steurer~\cite{BKS} have %
subsequently  %
%
shown  %
that the edge-expansion hypothesis also implies, in a black-box manner,
the combinatorial hypothesis of \cite{DKKMS2}. Finally, the edge-expansion hypothesis
is proved in its full generality in \cite{KMS2}
(with an additional insight gained from~\cite{KMMS}), thus completing the line of
research initiated by the current paper.





\section{The Grassmann graph and related hypotheses}\label{section:grassmann}
In this section we introduce the Grassmann graph and
related hypotheses that are relevant towards the soundness analysis
of our PCP construction (at the Inner PCP level).
The Grassmann graph leads to an encoding of a linear function
on a high-dimensional $\mathbb{F}_2$-vector space and a
$2$-to-$2$ test
%
 to check the encoding. The linear function
is encoded by writing down its restriction to all $\ell$-dimensional subspaces, the restrictions themselves being
linear functions on the respective subspaces.  %
Given a supposed encoding, \ie, an assignment of a linear function to every $\ell$-dimensional
subspace, one can test that the given
linear functions on a pair of $\ell$-dimensional subspaces are consistent
on their intersection. For the test to have the $2$-to-$2$ property, the test is performed only on a pair of $\ell$-dimensional
subspaces that intersect on a $(\ell-1)$-dimensional subspace.
Naturally, the following decoding
%
 question arises: Given an assignment to the $\ell$-spaces that demonstrates some consistency,
is there a global linear function that explains some (or almost all) of the consistency? We hypothesize that the answer
is affirmative, but subtle. Also, towards the analysis of our reduction, we have to deal
with more general notion of $(j,\delta)$-assignments,
%
where a $\delta$ fraction of the $\ell$-spaces are each assigned a list of $j$
linear functions on it.


\subsection{The Grassmann graph}
Let $\mathbb{F}$ be a field of size $p$, $V$ a linear space of dimension $n$ over
$\mathbb{F}$ and $1 \leq \ell \leq n-1$ a positive integer. The Grassmann graph
$G(V,\ell)_p$ is defined as follows. %
\begin{itemize}
\item The vertices are all $\ell$-dimensional subspaces of $V$.
\item The edges are pairs of vertices $L,L'$ such that ${\sf dim}(L\cap L')=\ell - 1$.
\end{itemize}
%

%
The Grassmann graph has been moderately studied in the literature, mainly in the context of distance-regular graphs~\cite{BCN12}.
%
%
Vector-space analogues of classical %
theorems such as
%
%
Erd\H{o}s--Ko--Rado and Kruskal--Katona %
are also known to hold~\cite{CP10}  %
and translate to properties of the Grassmann graphs.  %
Here are some
known facts regarding the Grassmann graph, though we do not necessarily need them.
\begin{fact}
Suppose $1 \leq \ell \leq \frac{n}{2}$.
\begin{itemize}
\item[] 1. The number of vertices in the graph $G(V, \ell)_p$  is the $\ell^{\mbox{th}}$ $p$-nomial coefficient (sometimes referred to as ``Gaussian binomial coefficient'')
\[
{n \brack \ell}_p \defeq \prod\limits_{i=0}^{\ell-1}{\frac{p^{n}-p^{i}}{p^{\ell}-p^{i}}}\,.%
\]
\item[] 2. The graph is regular with degree $d=\frac{p^n-p^\ell}{p^\ell-p^{\ell-1}}{\ell \brack \ell-1}_p$.  For $p=2$,  this is $\Theta(2^n)$.
\item[] 3. The eigenvalues of the (adjacency matrix of the) graph are
\[
\lambda_j = p^{j+1}{\ell-j \brack 1}{n-\ell-j \brack 1}  - {j \brack 1}_p\,.%
\]
with multiplicities ${n \brack j}_p - {n\brack j-1}_p$, for $j=0,...,\ell$ \footnote{${0 \brack 1}$  is defined to be $0$.}.

For $p=2$ the eigenvalues are approximately $(1+2^{j-\ell})2^{n+1-j}$, and so the normalized eigenvalues are $\Theta(2^{-j})$.
\end{itemize}
\end{fact}
We will only be interested in the case $p=2$ and the subscript $p$ will be omitted henceforth. The following observation will be useful.
\begin{fact}\label{fact:isomorphic}
Given the Grassmann graph $G(V,\ell)$,
${\sf dim}(V) = n$ and a $q$-dimensional subspace $Q \subseteq V, ~0 \leq q \leq \ell-1$, let ${\sf Zoom}_Q$ denote the subset
of vertices  $L \in G(V,\ell)$ such that $Q \subseteq L$. Then the induced subgraph on the subset ${\sf Zoom}_Q$ is isomorphic to
the (lower order) Grassmann graph $G(V',\ell')$ with ${\sf dim}(V')=n-q, ~\ell' = \ell- q$. A natural isomorphism is by letting
$V' = V/Q$ to be the quotient space.
\end{fact}

%
\subsection{\texorpdfstring{$(j,\delta)$}{(j,delta)}-assignments and zooming in}
Let $[2^\ell]$ denote the set of linear functions on an $\ell$-dimensional space. We would need to consider
the so-called \emph{$(j,\delta)$-assignments} %
%
to the vertices of the Grassmann graph.

\begin{definition}
The density of a set $S\subseteq G(V,\ell)$ is its fractional size, \ie,
\[
{\sf Density}(S)= \frac{|S|}{\card{\sett{L\subseteq V}{\dim (L) = \ell}}} = \frac{\card{S}}{{\dim(V) \brack \ell}}\,.%
\]
\end{definition}

\begin{definition}\label{def:edge_consistent} Let $S\subseteq G(V,\ell)$ and let $F: S \to {[2^\ell] \choose j}$
 assign, to each $\ell$-space  $L \in S$, a set  $F[L]$ of $j$ linear functions on $L$. The assignment $F$ is
 said to be $(j,\delta)$-edge-consistent, or simply a $(j,\delta)$-assignment, if
 \begin{itemize}
 \item $S$ has density at least $\delta$.
 \item If $L_1,L_2\in S$ are %
adjacent  %
then
there is a pair $a_1\in F[L_1],a_2\in F[L_2]$ such that $a_1,a_2$ agree on $L_1\cap L_2$.
\end{itemize}
\end{definition}

We consider the scenario where $j$ and $\delta$ are fixed,
%
 then $\ell$ is allowed to be
sufficiently large, and finally
the global dimension $n$ is allowed to be sufficiently large compared to $\ell$.
One would hope that a $(j,\delta)$-assignment implies the existence of
a global linear function $g\colon V\to\power{}$ that ``explains'' some of the consistency.
Specifically, is there a global linear function $g$ such that $g|_L \in F[L]$ for
 a %
$\delta' = \delta'(j,
\delta)>0$ fraction of the $\ell$-dimensional subspaces $L\subseteq V$? The answer turns out to be negative as seen from the following example.

\subsubsection{Subspace example}
Fix $Z\subseteq V$ to be a subspace of dimension $n-\ell$ and pick a set of vertices $S \subseteq G(V,\ell)$ as
\[
S=\sett{L\in G(V,\ell)}{\dimm(L\cap Z) =2}\,.%
\]
It is not difficult to see that $S$ has constant density which can be computed to be $\approx 0.20$.
We will exhibit a $(3, \approx 0.20)$-assignment $F[\cdot]$ which
has no non-trivial consistency with any global linear function.
For each $z \in Z \setminus \set{0}$, choose $f_z: V \to \{0,1\}$ to be a global linear function arbitrarily.
The assignment $F: S \to {[2^\ell] \choose 3}$ is now defined as
%
\[
F[L]=\{ f_z|_L \mid z \in (L\cap Z)\setminus\set{0}\}\,.%
\]
In words, $L$ is assigned three linear functions that are restrictions to $L$ of the three global linear
functions $f_z$ for $z \in  (L\cap Z)\setminus\set{0}$. Note that ${\sf dim}( L \cap Z) =2 $ and hence
$|F[L]|=3$. Now we show that
 if $L_1,L_2\in S$ have an edge connecting them, then the assignments $F[L_1], F[L_2]$ are consistent.
 Indeed, when ${\sf dim}(L_1)=\ell$,~${\sf dim}(L_2) = \ell$,~${\sf dim}(L_1 \cap L_2) = \ell-1$, ~${\sf dim}(L_1 \cap Z)=2$,  we have
 $$  1 \leqslant {\sf dim} (L_1 \cap L_2 \cap Z)  \leqslant 2\,.$$%
 In particular, there exists  $z \in (L_1 \cap L_2 \cap Z) \setminus \set{0}$. By design, we
  have $f_z|_{L_1} \in F[L_1]$
 and $f_z|_{L_2} \in F[L_2]$ and moreover that $f_z|_{L_1}, ~f_z|_{L_2}$ agree on  $L_1 \cap L_2$, both
 being restrictions of the same global function $f_z$. Finally, we note that since the choice of
 functions $f_z$ is arbitrary, no global linear function has a non-trivial consistency with the assignment $F[\cdot]$.
 For the sake of concreteness, one can let $f_z$ be the linear
 %
 function $x \to \left< z, x\right>$. It is not
 difficult to see that for any global linear function $f: V \to \{0,1\}$,
   $$
  \Prob{L \in G(V, \ell)}{f|_L \in F[L] } \leqslant 2^{-\Omega(\ell)}\,.%
   $$




\subsubsection{Zooming in}
We now consider a plausible way to circumvent the above example. For a Grassmann graph $G(V,\ell)$
and a subspace $Q \subseteq V$, let 
$$ {\sf Zoom}_Q =  \{ L \mid L \in G(V, \ell), ~Q \subseteq L \}$$
be the subset of vertices that contain $Q$ (we intend to take ${\sf dim}(Q) \ll \ell$).
 Similarly, for a subset $S \subseteq G(V,\ell)$, let 
$$ {\sf Zoom}_Q [S] =  \{ L \mid L \in S, ~Q \subseteq L \}$$
be the subset of vertices of $S$ that contain $Q$.
Returning to the subspace example above, let $F[\cdot]$ be the $(3, \approx 0.20)$-assignment to the Grassmann graph $G(V,\ell)$
as therein. Let $S$ be the subset of density $\approx 0.20$ to which $F[\cdot]$ actually assigns a list of $3$ linear functions.
We noted that no
global linear function has a non-trivial consistency with $F[\cdot]$. To be specific, for any global linear function $f: V \to \{0,1\}$,
 $$  \Prob{L \in G(V, \ell)}{f|_L \in F[L] }   \leqslant 2^{-\Omega(\ell)}.   $$
We observe however that there exists a one-dimensional subspace $Q$ such that after
zooming into $Q$ (\ie, conditioning on the $\ell$-spaces containing $Q$), there does exist a global linear function with good consistency with $F[\cdot]$.
Indeed, let $z \in Z \setminus \set{0}$ be an arbitrary point, $f_z$ be the global linear function associated with $z$ and let
$Q=\spn\set{z}$. For any subspace $L \in {\sf Zoom}_Q[S]$, we have $z \in Q \subseteq L$ and hence $f_z|_L \in F[L]$. Thus the global linear function $f_z$
is consistent with the assignment $F[\cdot]$ on every $L \in {\sf Zoom}_Q[S]$. Moreover,
$ {\sf Zoom}_Q[S] $ when regarded as a subset of ${\sf Zoom}_Q$ has a constant density, say $C$
(in fact its density is higher than the original density of $S$ in $G(V, \ell)$ which is $\approx 0.20$).
Thus $$  \Prob{L \in G(V, \ell)}{f_z|_L \in F[L] \mid Q \subseteq L }   \geqslant   C  $$
where the probability is conditional on the $\ell$-spaces containing $Q$. Further,
if the point $z$ were chosen at random from the global space $V$, with probability $\approx 2^{-\ell}$, we have $z \in Z \setminus \set{0}$ and
      then zooming into $Q = \spn\set{z}$ gives a global linear function with good consistency.
To summarize our specific example,

\medskip
\emph{Given a $(3,\approx 0.20)$-assignment $F[\cdot]$ to $G(V,\ell)$, for $\approx 2^{-\ell}$ fraction of one-dimensional subspaces
$Q \subseteq V$, zooming into $Q$ gives a global linear function that is $\Omega(1)$-consistent with $F[\cdot]$.}
\medskip


%
%
%
%
We now hypothesize that something to this effect always holds for any $(j,\delta)$-assignment to a Grassmann graph $G(V,\ell)$
when one is allowed to zoom into a $q$-dimensional subspace $Q$ with constant $q$ and the zoom-in succeeds for a non-negligible
fraction (that may depend arbitrarily on $\ell$)
of $q$-dimensional subspaces $Q$. Our main hypothesis appears below, followed by its variants and special cases.

\subsection{Our hypotheses}

%
\subsubsection{The main hypothesis}

\begin{conj}\label{conj:main2}
For every integer $j \geqslant 1$ and %
real  %
$\delta>0$, there exist an integer $q \geqslant 0$, %
a real  %
$C \in (0,1)$,
  and a function $\aconj\colon\mathbb{N}\to (0,1]$ such that
  for all sufficiently large integers $\ell$, for all sufficiently large integers $n$,
  the following holds:
  Let $F[\cdot]$ be a $(j,\delta)$-assignment to the Grassmann graph $G(V,\ell)$ with ${\sf dim}(V) = n$.
  Then for at least $\aconj{(\ell)}$ fraction of the $q$-dimensional subspaces $Q \subseteq V$,
  there exists a global linear function $g_Q: V \to \set{0,1}$ such that (note the conditional probability)
  \begin{equation} \label{eqn:global-fn} 
  \Prob{L \in G(V, \ell)}{g_Q|_L \in F[L] \mid Q \subseteq L}\geqslant C\,.%
  \end{equation}
\end{conj}


%
\subsubsection{Upper bound on list-decoding size}
In \expref{Hypothesis}{conj:main2}, given a $(j,\delta)$-assignment $F[\cdot]$,
a global linear function $g_Q$ satisfying
\expeqref{Equation}{eqn:global-fn} is viewed as a decoded
%
global linear function.
Naturally, such a decoded function is only useful when there are not many functions satisfying
\expeqref{Equation}{eqn:global-fn}, hence one would like to obtain an upper bound on the number of such
functions.
%
%
%
Indeed, a reasonable upper bound, stated below, follows from
%
a result    %
of Blinovsky~\cite{Bli1}.
We include a proof in \expref{Section}{appendix:list_dec_upper} for the sake of completeness.

\begin{theorem}\label{thm:list_dec_upper} Let $F[\cdot]$ assign to every $L \in G(V,\ell)$, a list $F[L]$ of
at most $j$ linear functions on $L$.  Let $Q$ be any $q$-dimensional subspace of $V$.
Then there are at most $\frac{jC}{C^2-j\cdot 2^{-(\ell-q)}}$ global functions $g$ for which
\[
\Prob{L\in G(V,\ell)}{g|_L\in F[L] \mid Q \subseteq L}\geqslant C\,.%
\]
\end{theorem}





%
\subsubsection{Hypotheses about connectivity of Grassmann graph}

We are far from understanding the case of general $j$ in the
\expref{Hypothesis}{conj:main2} and even the case $j=1$
presents interesting challenges. We show in \expref{Section}{section:j_1} that
\expref{Hypothesis}{conj:main2} in the case $j=1$ follows from the
\expref{Hypothesis}{conj:gras_expand} below without the need for zoom-in.
A linearity test is developed therein that could be of independent interest.

\begin{conj}\label{conj:gras_expand}
For every %
$\delta >0$,  there exists %
%
$\eps > 0$
  such that
  for all sufficiently large integers $\ell$, for all sufficiently large integers $n$,
 the following holds:
 Let $S$ be any set
of vertices in the Grassmann graph $G(V,\ell), ~{\sf dim}(V)=n$ with density at least $\delta$. Then the
induced subgraph on $S$ contains a connected component of density at least $\eps$.
\end{conj}

%
%
%
%
The hypothesis below  seems
like a natural related question, stating that the Grassmann graph is a
small-set vertex expander (defined below), and it could be a good starting point
for further investigations.
%
\begin{definition}
For a graph $G = (V,E)$, the \emph{neighbourhood} of a set %
%
$S\subseteq V$   %
of vertices  %
is defined as
  $\Gamma(S) = \sett{v\in V}{\exists u\in V, (u,v)\in E}$.
\end{definition}
\begin{definition}[Small-set vertex expander] %
\label{def:small-set}  %
A family of graphs, $\{G_{n,\ell} = (V_{n,\ell},E_{n,\ell})\}_{n,\ell\in\mathbb{N}}$,
%
%
is called a \emph{small-set vertex expander} if there %
exists a function $\lambda\colon (0,1)\to [0,\infty)$ such that
\begin{enumerate}
  \item $\lim_{\eps\rightarrow 0^{+}}\lambda(\eps)=\infty$.
  \item For all $\eps>0$ there are $N$ and $L$, such that for all $n\geq N$, $\ell\geq L$
and sets $S\subseteq V_{n,\ell}$ with size at least $\eps|V_{n,\ell}|$,
  it holds that $|\Gamma(S)|\geq \lambda(\eps)|S|$.
\end{enumerate}
\end{definition}

\begin{conj}\label{conj:v_expansion}
The family of Grassmann graphs $G_{n,\ell} = G(\{0,1\}^n,\ell)$ is a small-set expander.
\end{conj}

%
%
%
%
%
%
%
%
%
%
%
%
%
%
%
%
%
%
%
%
%
%

%


%
\subsubsection{Hypothesis with side condition}
As is standard, while composing the
%
Inner PCP with the Outer PCP,  %
we require that the decoded global linear
function $g_Q$ in \expeqref{Equation}{eqn:global-fn}, \expref{Hypothesis}{conj:main2} itself respects certain linear %
side condition.  %
We state a
variant of \expref{Hypothesis}{conj:main2} that takes into account the side condition and show
that this variant follows easily from \expref{Hypothesis}{conj:main2}.


\begin{definition}
A pair $(\set{h_1,...,h_r},(b_1,...,b_r))$, where $\{ h_i\in \power{n} \}_{i=1}^r$  are linearly
independent and $b_i\in \power{}$,
is called a \emph{side condition} for a function $g\colon \power{n}\to\power{}$. We say that $g$ respects the side
condition if $g(h_i)=b_i$, for every $i$.
\end{definition}



Note that when $g$ is a linear function respecting the side condition $(\set{h_1,\ldots,h_r},(b_1,\ldots,b_r))$,
the value of $g$ on the space $H=\spn\set{h_1,...,h_r}$ is fixed. We will often simplify notation and say
$g$ respects the side condition $H$, when $ (b_1,\ldots,b_r)$ is clear from the context. Note that the
vertices of the Grassmann graph $G(V,\ell)$ are
%
%
the  %
$\ell$-dimensional subspaces of $V$. Now we instead think of the
vertex set as
 $$  \sett{L\xor H}{L\in G(V,\ell)},$$
restricted to only those $L$ such that $L \cap H = \{0\}$ and moreover, if $L \xor H = L' \xor H$, then the two
vertices are identified together\footnote{
In our application, we will have $r = \frac{n}{3} \gg \ell$, so almost all $\ell$-dimensional spaces $L \subseteq V$ satisfy
$L \cap H = \{ 0 \}$. }. Note that ${\sf dim}(L\xor H) = \ell+r$. There is an edge between $L \xor H$ and $L' \xor H$ if and only
their intersection has dimension $\ell+r-1$. It can be easily seen that the resulting graph is isomorphic to a %
lower-order  %
Grassmann graph $G(V', \ell)$, where $V'\subseteq V$ is a
%
complement  %
to $H$ (\ie, $V'$ is a subspace and $V' \xor H = V, ~V'\cap H = \{0\}, ~{\sf dim}(V')=n-r$).


A \emph{$(j,\delta)$-assignment respecting the side condition} %
%
 is an assignment $F[\cdot]: S \to {[2^\ell] \choose j}$ to a set of
vertices $S$ such that
\begin{itemize}
\item $S$ has density at least $\delta$.
\item For each vertex $L\xor H \in S$, $F[L\xor H]$ is a list of $j$ linear functions
on $L \xor H$ that respect the side condition. Note that since the side condition already specifies the values
of a linear function on $H$, the number of linear functions on $L\xor H$ that respect the side condition is
exactly $2^\ell$.
\item For any $L\xor H, L' \xor H \in S$ that are %
adjacent  %
there are linear functions
      $a \in F[L \xor H], ~a' \in F[L' \xor H]$ that agree on the intersection $L\xor H \cap L' \xor H$.
\end{itemize}

We now state the variant of \expref{Hypothesis}{conj:main2} that takes into account the side condition. A (rather self-evident) proof that
\expref{Hypothesis}{conj:main} follows from \expref{Hypothesis}{conj:main2} appears in \expref{Section}{appendix:hyp_side_cond}.
%
%
%

\begin{conj}\label{conj:main}
  For every integer $j \geqslant 1$ and %
real  %
$\delta>0$, there exist an integer $q \geqslant 0$, a %
real  %
$C \in (0,1)$, and a
  function $\aconj\colon \mathbb{N}\to(0,1)$
  such that
  for all sufficiently large integers $\ell$, for all sufficiently large integers $n$,  the following holds:
  Let $F[\cdot]$ be a $(j,\delta)$-assignment respecting the side condition $\left( \{ h_i \}_{i=1}^r, ~\{b_i\}_{i=1}^r \right)$
  to the Grassmann graph $G(V,\ell)$ with ${\sf dim}(V) = n$ and $r \leq \frac{n}{3}$.
  Then for at least $\aconj{(\ell)}$ fraction of the $q$-dimensional subspaces $Q \subseteq V$,
  there exists a global linear function $g_Q: V \to \set{0,1}$ that respects the side condition
  such that (note the conditional probability)
  \begin{equation} \label{eqn:global-fn-side} 
  \Prob{L \in G(V, \ell)}{g_Q|_{L\xor H}  \in F[L\xor H] \mid Q \subseteq L}\geqslant C.
  \end{equation}
\end{conj}




\section{The outer PCP}\label{section:outer_pcp_game}

Our Outer PCP is a carefully constructed $2$-Prover-$1$-Round Game from a regular instance of the $3$-Lin problem. Recall (see
the paragraph before \expref{Theorem}{thm:main}) that an instance $(X, {\sf Eq})$ of the $3$-Lin problem consists of a set of
${\mathbb F}_2$-valued variables $ X$ and a set of equations ${\sf Eq}$, each equation containing three (distinct) variables. The
instance is regular if  every variable appears in exactly, say $5$,
 equations, and two distinct equations share at most one variable.
 Starting with a $3$-Lin instance given by H\aa stad's reduction \cite{Has01}, a standard sequence of transformations can turn the
 instance into a regular one, while preserving the near-perfect completeness and keeping the soundness bounded away from $1$.
 To summarize:

 \begin{theorem}\label{thm:regular_3lin}
 There exists an absolute constant $\frac{1}{2} < s^* < 1 $ such that for every constant $\eps > 0$, the ${\sf Gap3Lin}(1-\eps, ~s^*)$
 problem on regular instances is NP-hard.
 \end{theorem}



Let $(X, {\sf Eq})$ be an instance of ${\sf Gap3Lin}(1-\eps, s^*)$ as in \expref{Theorem}{thm:regular_3lin}. We intend to construct a $2$-Prover-$1$-Round Game
that is used as our Outer PCP.  Instead
of taking a passive view of $2$-Prover-$1$-Round Game as a constraint satisfaction problem as in \expref{Definition}{def:game},
it is more intuitive to take
an equivalent active view in terms of two provers and a probabilistic verifier. The two provers wish to convince the verifier that the $3$-Lin instance
is near-satisfiable. Since our construction has multiple subtle features, we present it incrementally, adding one feature at a time.
The construction is along the lines of \cite{KS}, smoothness and covering features are as therein and there is an additional
advice feature.



\subsection{Equation vs variable game}



We start with a standard ``equation vs variable'' game that the reader might be already familiar with. In this game,
the verifier chooses an equation $e\in {\sf Eq}$ uniformly at random, sends it to
the first prover, chooses a variable $x$ randomly from the three variables occurring in the equation $e$ and sends it to the second prover.
The provers are expected to provide an ${\mathbb F}_2$-value for each of the variables they receive.
The verifier accepts if and only if the first prover provides a satisfying assignment to $e$ and if
both provers give $x$ the same value.

\smallskip
\noindent {\bf Completeness:} Suppose there is an assignment to $(X, {\sf Eq})$ that satisfies
a %
$1-\eps$ fraction of the equations.
The provers can answer according to this assignment and the verifier accepts with probability at least $1-\eps$.

\smallskip
\noindent {\bf Soundness:} Suppose no assignment to  $(X, {\sf Eq})$
satisfies more than $s^*$ fraction of the equations. The strategy of the second prover is simply an assignment to all the variables.
This assignment fails to satisfy $1-s^*$ fraction of the equations. For every equation that fails, the second prover either has to
give inconsistent answer to at least one of its variables or answer with an unsatisfying assignment to the equation. Thus the provers
cannot make the verifier accept with probability more than $1 - \frac{1-s^*}{3}$ (\ie, bounded away from $1$).



\subsection{Smooth equation vs variable game}

We modify the equation vs variable game slightly and call it a smooth game.\footnote{Smoothness refers to the following
property of a game: consider a fixed question for the first prover and two distinct answers $a,b$ to the first prover,
sample the question for the second prover, and let $a',b'$ be the answers it must answer for the verifier to accept if
the first prover answered $a,b$, respectively.  %
The game is called smooth if with high probability over the choice of the
question to the second prover, it must be the case that $a'\neq b'$. The game described is smooth provided $\beta \ll 1$. }
 Let $\beta \in (0,1)$ be a smoothness parameter.
The verifier sends an equation $e$ to the first prover as before. To the second prover however,
the verifier sends a random variable $x$ occurring in $e$ with probability $\beta$, and sends the equation $e$
with probability $1-\beta$ (hence asking the same question to both the provers).

\smallskip
\noindent {\bf Completeness:} As before, the completeness is at least $1-\eps$.

\smallskip
\noindent {\bf Soundness:} The new game is effectively a trivial game with probability $1- \beta$ and is same
as the equation vs variable game  with probability $\beta$. Hence the soundness is at most $1- \Omega\left(\beta\right)$,
where the $\Omega$-notation hides the dependence on $s^*$ (which is an absolute constant anyways).





\subsection{Smooth equation vs variable game with advice}


Our application requires a further modification of the smooth game.
Roughly speaking,
the provers are also provided extra advice that acts like publicly shared randomness.
Nevertheless, this advice cannot considerably help the provers.\footnote{As far as we know, this is the first
instance of a PCP construction that incorporates the feature of advice.}

As before, the verifier picks an equation $e$ at random, say $x_{i_1} + x_{i_2} + x_{i_3} = b_{i}$,
and sends it to the first prover. With probability $1-\beta$, the second prover receives the
equation $e$ as well, and otherwise a single variable from the equation $e$ chosen at random.
Let $V\subseteq \set{x_{i_1},x_{i_2},x_{i_3}}$ be the set of variables sent to the second prover (so
$|V|$ is $1$ or $3$).
The verifier chooses an advice vector $a \in\power{V}$ at random. If $|V|=3$, define $a^* = a$, and
if $|V|=1$, let $a^*$ be obtained from $a$ by padding with $0$ in place of $\set{x_{i_1},x_{i_2},x_{i_3}}\setminus V$.
The verifier sends the first prover the vector $a^*$ and the second prover the vector $a$. As before, the
provers are expected to provide a value for each of the variable they receive.


Call this game $G_{\beta, 1}$. The extra advice could give the first prover a hint as to which variables the second prover receives.
For example, if the first prover's advice vector is $a^* = (0,0,1)$, she knows that the second prover has received
either all three variables or (just) the variable $x_{i_3}$.
However, when the first prover receives the vector $(0,0,0)$, she does not know whether the second prover has received
all three variables along with advice $a = (0,0,0)$ or a single variable, whose identity she does not know,
along with advice $a = (0)$. It is clear from this discussion that:

\smallskip
\noindent {\bf Completeness:} The completeness of game $G_{\beta, 1}$ is at least $1-\eps$.

\smallskip
\noindent {\bf Soundness:} The soundness of game $G_{\beta, 1}$ is at most $1- \Omega\left(\beta\right)$.

\smallskip
We further generalize to the game $G_{\beta, q}$ for any integer $q \geq 0$
where instead of sampling and sending the provers one pair,
%
%
$(a^*, a)$,
%
%
the verifier samples, %
independently, $q$ pairs $(a_1^*, a_1), \ldots, (a_q^*, a_q)$, and sends the list $[a_1^*, \ldots, a_q^* ]$ to the first prover and
the list $[a_1, \ldots, a_q]$ to the second prover. It is not difficult to see that:

\smallskip
\noindent {\bf Completeness:} The completeness of game $G_{\beta, q}$ is at least $1-\eps$.

\smallskip
\noindent {\bf Soundness:} The soundness of game $G_{\beta, q}$ is at most $1- \Omega\left(\frac{\beta}{2^q}\right)$. Intuitively,
the verifier rejects with constant probability when the second prover is sent a single variable (which happens with probability $\beta$)
along with the advice-list $[(0), \ldots, (0)]$ (which happens with probability $\frac{1}{2^q}$).

\begin{remark}\label{remark:constant_rej}
The soundness of the $\left(\frac{2^q}{\beta} \right)$-fold parallel
repetition game $G_{\beta, q}^{\otimes 2^q/\beta}$ is
%
less than %
an absolute constant less than $1$. Intuitively, in $\frac{2^q}{\beta}$ ``trials,''
with constant probability, there is a  ``coordinate'' on which
the second prover receives a single variable
along with the advice-list $[(0), \ldots, (0)]$,
and then the verifier rejects with a constant probability.
\end{remark}

\subsection{The final game (outer PCP) }\label{final_game}\label{sec:final-outer-pcp}

Finally, our Outer PCP is a $k$-fold parallel repetition of the game $G_{\beta, q}$, \ie, the game $G_{\beta, q}^{\otimes k}$.

\smallskip
\noindent {\bf Completeness:} The completeness of game $G_{\beta, q}^{\otimes k}$ is at least $1-k\eps$.

\smallskip
\noindent {\bf Soundness:} The soundness of game $G_{\beta, q}^{\otimes k} $ is at most $2^{-\Omega(\beta k/2^q)}$. The game can be considered
as $\frac{\beta k}{2^q}$-fold parallel repetition of the game $G_{\beta, q}^{\otimes 2^q/\beta}$ which has constant soundness as per
\expref{Remark}{remark:constant_rej}. One can then apply the parallel repetition theorem for
projection games with no dependency on the answer size as in \cite{Rao,DS14}.

\begin{remark} Let $U, V$ be the questions sent to the first and the second prover in the game $G_{\beta, q}^{\otimes k}$, not taking into
account the advice yet. Thus $U$ is a set of $3k$ variables and $V \subseteq U$ with expected size $\Expect{}{|V|} = 3k - 2\beta k$.
With a careful look, it can be seen that the advice-list for the first prover is a list $[x_1,\ldots, x_q]$ with $\forall ~1\leq i \leq q,
~ x_i \in \power{U}$. Similarly, the advice-list for the second prover is a list $[y_1,\ldots, y_q]$ with $\forall ~1\leq i \leq q,
~ y_i \in \power{V}$. Moreover, if one regards the space $\power{V}$ as a subspace of $\power{U}$ in a natural manner, then
$\forall ~1\leq i \leq q, ~x_i = y_i$. Thus the advice is to be interpreted as a list of $q$ points in $\power{V}$ that is sent to both provers.
\end{remark}





\section{The main reduction}\label{section:construction}
In this section we present our reduction towards proving \expref{Theorem}{thm:main}. The soundness analysis of the reduction is
presented in \expref{Section}{section:soundness}.

\subsection{Setting of the parameters}

Let $(X, {\sf Eq})$ be an instance of regular  ${\sf Gap3Lin}(1-\eps, s^*)$ as in \expref{Theorem}{thm:regular_3lin}. We will use the
game $G_{\beta, q}^{\otimes k}$ in \expref{Section}{sec:final-outer-pcp} as the Outer PCP. Since there are several parameters involved, we specify the (tedious) order in which the parameters are chosen.

Let $(j,\delta)$ be the parameters required of the reduction in \expref{Theorem}{thm:main}. Depending on $(j,\delta)$, let $q, C, \alpha(\cdot)$ be as given in
\expref{Hypothesis}{conj:main} and let $\ell$ be an integer large enough so that \expref{Hypothesis}{conj:main} holds for all sufficiently large integers $k$ ($= n$ therein). The soundness
analysis of the reduction shows (modulo \expref{Hypothesis}{conj:main}) that a $(j,\delta)$-coloring to the $2$-to-$2$ Game yields a prover strategy in the
Outer PCP with success probability roughly $\frac{\delta\cdot \alpha(\ell) \cdot C^2}{j}$. Conversely, by setting the parameters $\beta, k$ appropriately, the
 soundness of the Outer PCP, ~$2^{-\Omega(\beta k/2^q)}$, is ensured to be small enough beforehand so that the $2$-to-$2$ Game does not have a $(j,\delta)$-coloring.
In addition, we will want the construction to have the ``covering property'' (which we elaborate on
later), which requires $\beta  \sqrt{k}\cdot 2^\ell$ to be sufficiently small. One can choose $\frac{1}{k} \ll \beta \ll
\frac{1}{\sqrt{k}}$ and $k$ large enough so that both the soundness of the Outer PCP is small enough and the ``covering property'' holds.
 Finally the completeness parameter $1-\eps$ for the ${\sf Gap3Lin}$
instance is chosen to be close enough to $1$ so that the Outer PCP as well as the $2$-to-$2$ Game have completeness $1-k\eps \geq 1-\delta$.




\subsection{The reduction}

Consider the game $G_{\beta, q}^{\otimes k}$ and ignore the advice for now.
Let ${\cal U}$ and ${\cal V}$ denote the sets of questions asked to the first and the second prover, respectively.  %
Specifically, ${\cal U}$
is the set of all $k$-tuples of equations, $U = (e_1, \ldots, e_k)$ from the regular ${\sf Gap3Lin}$
instance $(X, {\sf Eq})$.
For our purposes, it will be convenient to retain only those ``legitimate'' $U = (e_1, \ldots, e_k)$ such that (a) the equations
 $e_1, \ldots, e_k$ are distinct and do not share variables and (b) for any pair of variables $x \in e_i$ and
 $y \in e_j, ~i\neq  j$, $x, y$ do not appear together in any equation in the instance $(X, {\sf Eq})$. Due to
 regularity of the instance $(X, {\sf Eq})$, every variable appears in a constant number of equations, and hence
 the fraction of $U$ that are not legitimate is negligible, \ie, $O(\frac{k^2}{|X|})$, and dropping these does
 not affect our analysis. We assume henceforth that ${\cal U}$ consists of only the legitimate  tuples $U$.

 The verifier in the game $G_{\beta, q}^{\otimes k}$ picks a $k$-tuple $U = (e_1,\ldots, e_k) \in {\cal U}$ uniformly at random
 and then constructs a $k$-tuple $V$ such that independently for $1 \leq i \leq k$, the $i^{th}$ element of
 $V$ is the equation $e_i$ with probability $1-\beta$ and is a variable in the equation $e_i$ with probability $\beta$.
 Thus the set of questions ${\cal V}$ to the second prover consists of ``mixed'' tuples. In the following, we will work
 only with the set ${\cal U}$ and the role of the set ${\cal V}$ will be implicit.





We are now ready to describe the Transitive $2$-to-$2$ Game $G_{2:2} ( V(G_{2:2}), E(G_{2:2}), \Sigma, \Phi)$ that our
 reduction constructs. For any $U \in {\cal U}$, we regard $U$ as the tuple of $k$ equations $(e_1,\ldots, e_k)$ as well
 as the set of $3k$ variables appearing in these equations, say $(x_{11}, x_{12}, x_{13}, \ldots, x_{k1}, x_{k2}, x_{k3})$.
For each equation $e_i$, define a vector $v_{e_i}\in \power{U}$ that has $1$ on coordinates corresponding
to variables in $e_i$ and $0$ on the rest.
Denote $H_U={\sf Span}{\set{v_{e_1},...,v_{e_k}}}$ referred to as the space of side conditions. Let $b_1, \ldots, b_k \in \power{}$
be the ``right hand sides'' of the equations, \ie, the equation $e_i$ is $x_{i1} + x_{i2} + x_{i3} = b_i$. Define
\[
{\cal L}_U\defeq \sett{L\subseteq \power{U}}{{\sf dim}(L)=\ell,  ~L \cap H_U = \{0\}}\,.%
\]
 Note that for $L \in {\cal L}_U$, its intersection with $H_U$ is trivial and hence ${\sf dim}(L \xor H_U) = \ell+k$.
 Also, $|U|=3k$, ${\sf dim}(H_U) = k$ and ${\sf dim}(L)=\ell$. The fraction of $\ell$-spaces $L \subseteq \power{U}$ such that
$L \cap H_U \neq  \{0\}, ~L \not\in {\cal L}_U$ is negligible ($\approx 2^{\ell-2k}$, see \expref{Fact}{fact:aux_2}).

\medskip
\noindent \textbf{Vertices of $G_{2:2}$}: The game $G_{2:2}$ has a block of vertices ${\sf Block}[U]$ for every $U \in {\cal U}$ defined as
  $${\sf Block}[U] = \sett{L\xor H_U}{ L\in{\cal L}_U}\,.$$%
The vertex set of $G_{2:2}$ is the (disjoint) union of all blocks:
\[
V(G_{2:2})=  \bigcup_{U \in {\cal U}}  ~{\sf Block}[U]\,.%
\]

\medskip
\noindent \textbf{Colors of  $G_{2:2}$}: The set of colors $\Sigma$ has size $|\Sigma|=2^\ell$. For a vertex $L\xor H_U$, its color
set $\Sigma$ is identified with  
 $$ \{ \psi: L \xor H_U \to \power{} \mid \psi ~\mbox{is linear}, \ \forall ~1\leq i \leq k, ~\psi(v_{e_i})=b_i  \}\,.$$%
 In words, the
 vertex $L\xor H_U$ is  to be assigned a linear function $\psi: L \xor H_U \to \power{}$
that respects the side conditions, meaning $\psi(v_{e_i}) = b_i$ for $1 \leq i \leq k$. Since the values of $\psi$ are already determined
on $H_U$, there are exactly $2^\ell$ eligible linear functions $\psi$.




\medskip
\noindent \textbf{Edges and Constraints of $G_{2:2}$}: Towards defining the edges and constraints of the game $G_{2:2}$, we stress
a notational (and perhaps conceptual) point.  $X$ is the set of all variables in the ${\sf Gap3Lin}$ instance, so $U \subseteq X$ and
 $\power{U}$ is a subspace of
$\power{X}$ in a natural manner. Every subspace under consideration can be considered as a subspace of $\power{X}$ and we can freely
take the intersections or direct sums of subspaces. For instance
 if $U_1,U_2$ are two sets of variables and $L_1\subseteq \power{U_1}$,
 $L_2\subseteq \power{U_2}$ are subspaces, we can consider both $L_1, L_2$ as subspaces of $\power{U_1 \cup U_2}$ (which in turn is
 a subspace of $\power{X}$) and then the subspaces $L_1\cap L_2, L_1+ L_2$ make sense.

\smallskip We are ready to define the edges and the constraints of $G_{2:2}$. For $U, U' \in {\cal U}$ (allowing the possibility that $U=U'$), we describe the edges
between their respective blocks.\footnote{Here $U, U'$ are thought of as sets of variables that are nearly identical. There might be
edges between the blocks of $U, U'$ that differ significantly, but those edges are merely ``accidental'' and do not have much relevance
towards the soundness of the reduction.}
There is an edge between vertices $L\xor H_U, ~L'\xor H_{U'}$ if either of the two conditions holds. Either
\begin{equation}\label{eqn:1-to-1}
{\sf dim} (L + H_U + H_{U'}) = {\sf dim} (L' + H_U + H_{U'}) = {\sf dim} (L + L' + H_U + H_{U'})\,,%
\end{equation}
in which case, the constraint is $1$-to-$1$, or
\begin{equation}
 {\sf dim} (L + H_U + H_{U'}) ={\sf dim} (L' + H_U + H_{U'}) = {\sf dim} (L + L' + H_U + H_{U'}) - 1\,,%
\end{equation}
in which case, the constraint is $2$-to-$2$. This definition is, admittedly, rather mysterious and we try to clarify it somewhat. We recommend reading the proofs of \expref{Lemmas}{lem:intuition_aux1} \expref{and}{lem:intuition_aux3} %
to start having some intuition. We first consider the $1$-to-$1$ constraints.

\medskip \noindent {\bf $1$-to-$1$ Constraints:} From \expref{Lemma}{lem:intuition_aux1}, we always have
 $${\sf dim} (L + H_U + H_{U'}) = {\sf dim} (L' + H_U + H_{U'})\,.$$%
If, in addition, this dimension is same as that of $L + L' + H_U + H_{U'}$ which contains both the spaces above, then all the three spaces   must be identical, \ie,
$L + H_U + H_{U'} = L' + H_U + H_{U'} = L + L' + H_U + H_{U'} = Z$, say.
From \expref{Lemma}{lem:intuition_aux3}, there is a $1$-to-$1$ correspondence between linear functions on $L\xor H_U$
(that respect the side condition on $H_U$) and linear functions on $L+ H_U + H_{U'}=Z$  (that respect the side condition on both $H_U, H_{U'}$), and the same holds
between $L' \xor  H_{U'}$ and $L' + H_U + H_{U'} = Z$.
This gives a $1$-to-$1$ correspondence between linear functions on $L\xor H_U$ and $L'\xor H_{U'}$ (respecting the relevant side conditions)
which is regarded as the $1$-to-$1$ constraint on the
%
coloring  %
of $L\xor H_U$ and $L'\xor H_{U'}$.


\medskip \noindent {\bf $2$-to-$2$ Constraints:} As before,
from \expref{Lemma}{lem:intuition_aux1}, we always have
 $${\sf dim} (A = L + H_U + H_{U'}) = {\sf dim} (A' = L' + H_U + H_{U'}) = d ~\mbox{(say)}.$$
Now suppose that $Z = L + L' + H_U + H_{U'}, ~{\sf dim}(Z) = d+1$. Since $Z = A \xor A'$, it follows
that ${\sf dim}(A \cap A') = d-1$. Thus, it is possible to choose a basis $I$ for $A \cap A'$ and $v \in L, ~v' \in L'$
so that $I \cup \{ v \}$ is a basis for $A$,  $I \cup \{ v' \}$ is a  basis for $A'$, and $I \cup \{v, v'\}$ is a basis for $Z$. In the
following, all linear functions considered are supposed to respect the side condition on $H_U$ or $H_{U'}$ or both, depending on
whether the relevant space contains $H_U$, $H_{U'}$ or both.

Every linear function $f$ on $A \cap A' = {\sf Span}(I)$ has exactly two extensions $f_1, f_2$ to $A = {\sf Span}(I \cup \{v\})$,
depending on their value on $v$, and has exactly two extensions $f'_1, f'_2$ to $A' = {\sf Span}(I \cup \{v'\})$,
depending on their value on $v'$. Moreover by \expref{Lemma}{lem:intuition_aux3}, linear functions on $A$ are in one-to-one correspondence with
those on $L\xor H_U$.
Denote by $\tilde{f}_1, \tilde{f}_2$ the mates of $f_1, f_2$, respectively, via this correspondence.  %
Similarly, linear functions on $A'$ are in one-to-one correspondence with
those on $L'\xor H_{U'}$ and let $\tilde{f}'_1, \tilde{f}'_2$ be the mates of $f'_1, f'_2$. This gives a $2$-to-$2$ constraint
between $L\xor H_U$ and $L' \xor H_{U'}$ that matches the pair $(\tilde{f}_1, \tilde{f}_2)$ with the pair $(\tilde{f}'_1, \tilde{f}'_2)$.

\paragraph{Transitivity:} The transitivity of $G_{2:2}$ is proved in \expref{Section}{appendix:transitive}.

\begin{remark}\label{remark:useful_222}
Another useful way to describe the constraint, both in the $1$-to-$1$ and $2$-to-$2$
%
cases,  %
is as follows: if there is a space $Z$ that includes both
$L\xor H_U$ and $L' \xor H_{U'}$ and has an assignment $\beta$ that respects the side conditions on $H_{U}, H_{U'}$, then
$\beta|_{L\xor H_U}, \beta|_{L' \xor H_{U'}}$ are colorings to $L\xor H_U$ and  $L'\xor H_{U'}$, respectively that satisfy the constraint.  %
In both cases above $Z$ happens to be $L+ L' + H_U + H_{U'}$, but we will have occasion to use an even larger space $Z$ in certain proofs.
\end{remark}

\begin{remark}\label{remark:transitivity_why}
  The reason we do not identify vertices $L\oplus H_U$ and $L'\oplus H_{U'}$ that have a
  $1$-to-$1$ constraint between them, and instead resort to the transitivity property is due
  to the completeness of the reduction. A priori, if $(X,{\sf Eq})$ has a good solution $A$,
  then it may potentially be the case that while it satisfies all of the equations of $U$,
  it may fail to satisfy some of the equations of $U'$, in which case it would not be clear
  how to assign the vertex resulting from the identification of $L\oplus H_U$ and $L'\oplus H_{U'}$.
  Transitivity may be seen as a softer way of doing this identification which enjoys all of the benefits
  proper identification would have.
\end{remark}


%
\subsubsection{Auxiliary lemmas}

\begin{lemma}\label{lem:intuition_aux1} Let $U, U' \in {\cal U}$ and
let  %
${\sf Eq}[U], {\sf Eq}[U']$ denote the sets of equations ($k$ in number)
in $U, U'$, respectively.  %
%
Then, for
$L \in {\cal L}_U$, $$ {\sf dim} (L + H_U + H_{U'}) = \ell + 2k  - |{\sf Eq}[U] \cap {\sf Eq}[U']|\,.$$%
\end{lemma}
\begin{proof}Let ${\sf Eq}[U'] = \{ e'_1, \ldots, e'_k \}$ and recall that $H_{U'} = {\sf Span}(v_{e'_1},\ldots,v_{e'_k} )$.
Let $C$ denote the %
\emph{current space}  %
that is initialized to  $C = L \xor H_U$ and has
dimension $\ell+k$. We consider
equations $e'_1, \ldots, e'_k  \in {\sf Eq}[U']$ one by one,
and check whether %
%
%
%
%
%
%
%
%
%
$v_{e'_i}$  
belongs to the   %
%
%
%
to the current space.
%
If %
$e'_i \in {\sf Eq}[U]$, then $v_{e'_i} \in H_U$ already,
and hence ${\sf dim} (C + {\sf Span}(v_{e'_i}))= {\sf dim}(C)$.
Otherwise $e'_i \not\in {\sf Eq}[U]$ and shares at most one variable with $U \cup (U' \setminus e'_i)$. This is where we use the fact that $U, U'$ are ``legitimate''
tuples in the sense described in the first paragraph of current section.
Thus $v_{e'_i}$ is linearly independent of
%
$L + H_U +\sum_{j \neq  i} {\sf Span}( v_{e'_j} )\,.$ %
%
Hence $C \xor  {\sf Span}(v_{e'_i})$ has dimension $1$ larger than that of $C$. Carrying the
argument for $i=1,\ldots,k$ shows that in the end $C = L + H_U +H_{U'}$ and ${\sf dim}(C)$ is as desired.
\end{proof}

%
%
%
%
%


\begin{lemma}\label{lem:intuition_aux3} Let $U, U' \in {\cal U}$ and $L \in {\cal L}_U$.  Then any linear function on $L\xor H_U$ that respects the side condition on $H_U$,
has a unique extension to $L + H_U + H_{U'}$ that respects the side condition on both $H_U$ and $H_{U'}$.
\end{lemma}
\begin{proof} Let $f$ be a linear function on $L\xor H_U$ that respects the side condition on $H_U$. Clearly, it has at most one extension
to $L + H_U + H_{U'}$ that respects the side condition on both $H_U$ and $H_{U'}$, so the main point is to show that  there indeed is
such an extension. Similar to the proof of \expref{Lemma}{lem:intuition_aux1}, let $C$ denote the current space, $g$ denote the current linear function on $C$,
 so that initially $C = L \xor H_U, g = f$ and at each step,
$g$ respects the side condition on $H_U$ and the side condition due to equations $e'_1, \ldots, e'_{i-1}$ considered so far. Consider
the equation $e'_i$. If $e'_i \in {\sf Eq}[U]$ then $C + {\sf Span} (v_{e'_i}) = C$ and we keep $g$ unchanged and proceed next. If $e'_i \not\in {\sf Eq}[U]$, then
as in the proof of \expref{Lemma}{lem:intuition_aux1}, $v_{e'_i}$ is linearly independent of
$L + H_U +_{j \neq  i} {\sf Span}( v_{e'_j} ) . $ Hence $C \xor  {\sf Span}(v_{e'_i})$ has dimension $1$ larger than that of $C$ and
the function $g$ can be safely extended to
the  %
vector $v_{e'_i}$ as required. To be precise, one sets $g(v_{e'_i}) = b'_i$ where $b'_i$ is the
`right hand side'' of the equation $e'_i$ and then extends $g$ linearly to $C \xor  {\sf Span}(v_{e'_i})$. Carrying the argument for $i=1,\ldots,k$,
completes the proof.
\end{proof}

\subsection{Completeness}\label{section:completeness}
It is easily observed that the reduction satisfies the completeness condition as in \expref{Theorem}{thm:main}. Let
 $\sigma \colon X \to \power{}$ be an assignment to the ${\sf Gap3Lin}$ instance $(X, {\sf Eq})$  that satisfies
 a %
 $1-\eps$ fraction of
 the equations. Let ${\sf Eq}'$ denote the set of the equations satisfied so that $|{\sf Eq}'| \geq (1-\eps) |{\sf Eq}|$.
Let $~{\cal U}' \subseteq {\cal U}$ be the subset of $k$-tuples of equations $U$ such that all its $k$ equations are satisfied, \ie,
$U \subseteq {\sf Eq}'$. Clearly, $|{\cal U}'| \geq (1-k\eps)|{\cal U}| \geq (1-\delta) |{\cal U}|$ by choosing $\eps$ sufficiently
small.
%

%
For every $U \in {\cal U}'$, let $\sigma[U]$ denote the linear function on $\power{U}$ that maps $x \in \power{U}$ to $\left< \sigma|_U, x \right>$.
Since $\sigma$ satisfies all equations inside $U$, the linear function $\sigma[U]$ respects the side condition $H_U$. Now assign to every vertex
$L\xor H_U$ in ${\sf Block}[U]$, the linear function $\sigma[U]|_{L\xor H_U}$. We show that this assignment satisfies all constraints whose
both endpoints have been assigned.
Indeed if $(L\xor H_U, ~L'\xor H_{U'})$ is a constraint such that both endpoints are assigned,
then the constraint  is satisfied since all spaces are assigned using the same global assignment $\sigma$. Thus the $2$-to-$2$ Game has a
$(1, 1-\delta)$-assignment.


\subsection{Covering property}\label{section:covering}

We need a certain covering property towards the soundness analysis. While this property was introduced in \cite{KS}, we need a more general notion.
The covering property, the zoom-in required in \expref{Hypothesis}{conj:main}, and the ``advice'' in the 2-Prover-1-Round game in \expref{Section}{sec:final-outer-pcp} (the Outer PCP) are all used in a coordinated  manner in the soundness analysis.

\smallskip
Let $U$ be the set of $3k$ variables in a fixed set of $k$ equations.
We recall that in the Outer PCP game, the verifier chooses
$V\subseteq U$ randomly by choosing from each equation independently
(a) with probability $\beta$, one of the variables from the equation and
(b) with probability $1-\beta$, all three  variables from the equation.
We consider $\power{V}$ as a subspace of $\power{U}$ in a natural manner.
Slightly rephrasing a result from \cite{KS}, the statistical distance
between the following two distributions over one-dimensional subspaces of $\power{U}$ is small, \ie, at most $O(\beta \sqrt{k})$.\footnote{The intuition is as follows.
 A one-dimensional subspace is same as a non-zero point. A random point in $\power{U}$ (and in $\power{V}$) has negligible chance of being zero, so we
 might as well consider the distribution of (a) a random point in $\power{U}$ and (b) a random point in $\power{V}$ after choosing $V$ and then ``lifting it up''
 by appending $0$ in the coordinates $U \setminus V$. We note that $|U\setminus V| \approx 2\beta k$.
 A point chosen from the second distribution has  $\approx 2\beta k$ more zeroes
than that from the first distribution. However the imbalance between the number of zeroes and ones in a typical point in
$\power{U}$ is $\approx \sqrt{k}$, so when  $\beta k  \ll \sqrt{k}$, a deviation of $2\beta k$ zeroes is nearly imperceptible.}
\begin{itemize}
\item Choose a random one-dimensional subspace $P \subseteq \power{U}$.
\item Choose $V \subseteq U$ as described above, choose a random one-dimensional subspace $P' \subseteq \power{V}$ and regard it as a
      subspace of $\power{U}$.
\end{itemize}
We will need an analogous statement regarding two distributions over $\ell$-dimensional subspaces of $\power{U}$. We define the two distributions
below and prove the subsequent lemmas in \expref{Section}{appendix:covering}.
\begin{definition}\label{def:covering_dist} Let $U$ be a fixed set of $k$ equations and $V \subseteq U$ be chosen as above with
parameter $\beta$. Let $\ell \geq 1$ be an integer. Let ${\cal L}, ~{\cal L}'$ be distributions over $\ell$-dimensional subspaces of $\power{U}$ sampled as follows.
\begin{itemize}
\item ${\cal L}$: Choose  a uniformly random $\ell$-dimensional subspace of $\power{U}$.
\item ${\cal L'}$: Choose $V \subseteq U$ as above,  choose  a uniformly random $\ell$-dimensional subspace of $\power{V}$ and regard it as a
      subspace of $\power{U}$.
\end{itemize}
\end{definition}


The covering property, stated below, asserts that the distributions $\mathcal{L},\mathcal{L}'$ (and some conditioning of theirs)
are close in statistical distance. Here and throughout, if $P, Q$ are probability distributions over a universe $\mathcal{U}$, then
the statistical distance between them is defined as ${\sf SD}(P,Q) = \frac{1}{2}\sum\limits_{u\in\mathcal{U}}\card{P(u) - Q(u)}$.
\begin{lemma}\label{lemma:covering_normal}
Suppose $2^{\ell}\beta \leqslant \eigth$. Let ${\cal L}, ~{\cal L}'$ be distributions over $\ell$-dimensional subspaces over $\power{U}$
 sampled as in \expref{Definition}{def:covering_dist}. Then the statistical
distance between ${\cal L}, ~{\cal L}'$ is bounded as
$$
{\sf SD}({\cal L}, ~{\cal L}')\leqslant  \beta \sqrt{k}\cdot 2^{\ell+4}\,.$$%
\end{lemma}


\begin{lemma}\label{lemma:covering_cond_markov} Let $0 \leq q \leq \ell-1$ be an integer. Let $Q$ be
a $q$-dimensional subspace of $\power{U}$. Let ${\cal L}_Q$ and ${\cal L}'_Q$ be
distributions ${\cal L}$ and ${\cal L}'$ conditioned on the event that a sampled $\ell$-subspace $L$ contains $Q$.
Suppose $2^{\ell}\beta \leqslant \eigth$. Then for at least $1 - \sqrt{\beta}~k^{\frac{1}{4}} $
fraction of $Q$,
\begin{equation}\label{eq:SD_cond_Q}
{\sf SD}  ({\cal L}_Q, ~{\cal L}'_Q ) \leqslant \sqrt{\beta} ~k^{\frac{1}{4}} \cdot 2^{\ell+5}\,.%
\end{equation}
\end{lemma}


\section{Soundness analysis}\label{section:soundness}

In this section, given a $(j,\delta)$-assignment to the game $G_{2:2}$ constructed in \expref{Section}{section:construction},
we show how to extract a
provers' strategy in the Outer PCP game ($G_{\beta, q}^{\otimes k}$ as in \expref{Section}{sec:final-outer-pcp})
that succeeds with probability
$p = p(j, \delta,\ell)$. If the soundness of the Outer PCP
game is chosen to be smaller than $p$ to begin with, it implies that the game $G_{2:2}$ has no $(j,\delta)$-assignment,
proving \expref{Theorem}{thm:main}. Formally:
\begin{lemma}\label{lem:main_soundness}
  Suppose that there is a $(j,\delta)$-assignment to the game $G_{2:2}$ constructed in \expref{Section}{section:construction}.
  Then there is a strategy for the provers in the game $G_{\beta, q}^{\otimes k}$ that achieves winning probability
  $p = p(j,\delta,\ell) > 0$.
\end{lemma}

The rest of this section is devoted for the proof of \expref{Lemma}{lem:main_soundness}.
We recall that the first prover (the ``larger'' prover) receives as a question, a set $U$ of $3k$ variables (in $k$ equations) and an advice-list
$[x_1, \ldots, x_q]$ of $q$ points in $\power{U}$ (in fact in $\power{V}$ as stated next). The second prover (the ``smaller'' prover) receives as a question a subset
$V \subseteq U$ of variables and an advice-list $[x_1, \ldots, x_q]$ of the same $q$ points in $\power{V}$. We will extract the provers' strategies in
the next two subsections, and show in the last subsection that these strategies succeed with good probability.





\subsection{Strategy for the first (larger) prover }
We recall that a typical vertex in the game $G_{2:2}$ is denoted as $L\xor H_U$. Specifically, for $U \in {\cal U}$, the block of
vertices corresponding to $U$ is\footnote{We emphasize that ${\sf Block}[U]$ contains the vertex $L\xor H_U$ for essentially all
$\ell$-dimensional subspaces $L \subseteq \power{U}$.} 
 $$ {\sf Block}[U] = \{ L \xor H_U \mid L \subseteq \power{U}, ~{\sf dim}(L)=\ell, ~L \cap H_U = \{0\} \}\,,$$%
 and the set of vertices of $G_{2:2}$ is the union of all blocks of vertices over $U \in {\cal U}$. Let $F[\cdot]$ be the
 given $(j,\delta)$-assignment to the game $G_{2:2}$. Let us emphasize that this means:
 \begin{itemize}
 \item For at least
 a %
 $\delta$ fraction of the vertices $L \xor H_U$, a
 list $F[L\xor H_U]$ of $j$ linear functions on $L\xor H_U$ (that respect the side condition on $H_U$) is given. The remaining vertices
 are unassigned and do not play any role in the analysis.
 \item If there is a $2$-to-$2$ constraint between vertices $L\xor H_U$ and $L'\xor H_{U'}$, both of which are assigned, then there are linear functions
       $f \in F[L\xor H_U], ~f' \in F[L' \xor H_{U'}]$ that satisfy the constraint.
\item If there is a $1$-to-$1$ constraint between vertices $L\xor H_U$ and $L'\xor H_{U'}$, both of which are assigned, there is a one-to-one correspondence
      between the lists $F[L\xor H_U], ~F[L'\xor H_{U'}]$, via the same one-to-one correspondence that defines the $1$-to-$1$ constraint.
\end{itemize}

By an averaging argument, for at least $\frac{\delta}{2}$ fraction of the tuples $U$,
at least $\frac{\delta}{2}$ fraction of vertices in ${\sf Block}[U]$ are assigned. Call such a tuple $U$ good and let ${\cal U}_{\sf good}$ be
the set of good tuples with $|{\cal U}|_{\sf good} \geq \frac{\delta}{2} \cdot |{\cal U}|$.

Let the question to the first prover be $U \in {\cal U}$ along with the advice-list  $[x_1, \ldots, x_q]$ of points in $\power{U}$. If
$U \not\in {\cal U}_{\sf good}$, the prover gives up, so let us assume $U \in {\cal U}_{\sf good}$, and
let ${\sf Assigned}[U] \subseteq {\sf Block}[U]$ denote the set of vertices in its block that have been assigned,
$|{\sf Assigned}[U]| \geq \frac{\delta}{2} \cdot |{\sf Block}[U]|$. Since this is a $(j,\frac{\delta}{2})$-assignment respecting the
side condition on $H_U$, \expref{Hypothesis}{conj:main} states that for some $q, \alpha(\cdot), C$ that depend on $(j,\frac{\delta}{2})$, for
at least $\alpha(\ell)$ fraction of the $q$-dimensional subspaces $Q \subseteq \power{U}$,
there exists a global linear
function $g_Q : \power{U} \to \power{}$ that respects the side condition on $H_U$ and
\begin{equation} \label{eqn:soundness-1}  
  \Prob{L \subseteq \power{U}, ~{\sf dim}(L)=\ell}{g_Q|_{L\xor H_U}  \in F[L\xor H_U] \mid Q \subseteq L}\geqslant C\,.%
\end{equation}
We call such a choice of $Q$ ``lucky'' and let ${\cal Q}_{\sf lucky}$ be the set of all lucky $q$-dimensional subspaces of
$\power{U}$.  We note that the parameter $q$ was chosen beforehand to exactly match with that arising in \expref{Hypothesis}{conj:main}.
Moreover, call a $q$-dimensional subspace $Q$ ``smooth'' if it satisfies \expeqref{Condition}{eq:SD_cond_Q} in \expref{Lemma}{lemma:covering_cond_markov}\footnote{To recall,
the condition is that the distributions ${\cal L}_Q$ and ${\cal L}'_Q$ are close in statistical distance; the former distribution
chooses a random $\ell$-subspace of $\power{U}$ containing $Q$ and the latter distribution chooses a random question $V \subseteq U$ to the
second prover and then a random $\ell$-subspace of $\power{V}$ containing $Q$.}, and let ${\cal Q}_{\sf smooth}$ be the set of all smooth
$q$-dimensional subspaces of
$\power{U}$.

The prover looks at the advice-list and lets $Q = {\sf Span}(x_1,\ldots, x_q)$. If $Q \not\in {\cal Q}_{\sf lucky}$ or
$Q \not\in {\cal Q}_{\sf smooth}$,  the prover
gives up. Otherwise the prover picks a global linear function $g_Q:\power{U} \to \power{}$ respecting the side condition and satisfying \expeqref{Equation}{eqn:soundness-1}
(if there is more than one, one of them is picked arbitrarily), and outputs $g_Q$ as the answer. Strictly speaking,
 the linear functions $g_Q$  amounts to a function
$x \to \left< \sigma_Q, x\right> $ on $\power{U}$ for some $\sigma_Q \in \power{U}$ and the prover answers that $\sigma_Q$ is the
assignment to the $3k$ received variables ($\sigma_Q$ satisfies the $k$ equations as $g_Q$ respects the side condition). However, it
is more convenient to view the function $g_Q$ itself as the answer.

We note that a uniformly random $q$-subspace of $\power{U}$ is lucky with probability $\geq \alpha(\ell)$ (by \expref{Hypothesis}{conj:main})
and is smooth with probability $\geq 1-\beta k^{\frac{1}{4}} \geq 1- \frac{\alpha(\ell)}{2}$ when the parameters $\beta, k$ are
chosen appropriately.
Thus with probability at least $\frac{\alpha(\ell)}{2}$, the space $Q$ dictated by the advice-list is both lucky and smooth.


\subsection{A Strategy for the second (smaller) prover}
Let the question to the second prover be $V$ along with the advice-list  $[x_1, \ldots, x_q]$ of points in $\power{V}$. Let
$Q = {\sf Span}(x_1,\ldots,x_q)$ be the $q$-dimensional subspace of $\power{V}$. Let ${\cal L}_V$ be the set of all $\ell$-dimensional
subspaces of $\power{V}$ (though when the prover decides on an answer, only the subspaces containing $Q$ are relevant):
\[
{\cal L}_V =\sett{L}{L\subseteq \power{V}, ~{\sf dim}(L)=\ell}\,.%
\]
The prover first obtains an assignment $F_V[\cdot]$ to ${\cal L}_V$. Fix
 $L\in {\cal L}_V$. The prover examines every $k$-tuple of equations $U$ such that $V \subseteq U$, \ie, every question that could
 have been asked to the
 first prover, when the question of the second prover is $V$. Note that $L \subseteq \power{V} \subseteq \power{U}$ and hence there is a vertex $L\xor H_U$
 of the game $G_{2:2}$ in ${\sf Block}[U]$. If the vertex $L\xor H_U$ has been assigned, then the prover defines
\[
F_V[L]  = \{ f|_L  \mid f \in F[L\xor H_U] \}\,,%
\]
\ie, restrictions of all functions in $F[L\xor H_U]$ to $L$. In general there are several $U$ that contain $V$, so a priori, there is ambiguity
in the definition of $F_V[\cdot]$. The claim below shows however that the definition is unambiguous.
\begin{claim}
$F_{V}[\cdot]$ is well defined. That is, if $V \subseteq U, ~V \subseteq U'$ and if $L\xor H_{U}$, ~$L\xor H_{U'}$ are both assigned, then the
restrictions of $F[L\xor H_{U}]$ and $F[L\xor H_{U'}]$ to $L$ are identical.
\end{claim}
\begin{proof}
Notice that $L\xor H_{U}, L\xor H_{U'}$ have a $1$-to-$1$ constraint between them in the game $G_{2:2}$. By \expref{Definition}{def:j,delta}, a $(j,\delta)$-assignment
must assign identical sets of $j$ assignments to the vertices that have a $1$-to-$1$ constraint between them.
\end{proof}

Once $F_{V}$ is defined, the prover zooms into $Q$,
 and chooses at random any linear function
$h_Q\colon \power{V}\to\power{}$ that satisfies (if one exists and if so, we show that the list-size is bounded)
\begin{equation} \label{eqn:prover2-strat}
\Prob{L \in {\cal L}_V}{h_Q|_{L} \in F_V[L]] \mid Q \subseteq L}\geqslant \frac{C}{4}\,.%
\end{equation}
The prover outputs $h_Q: \power{V} \to \power{}$ as the answer. Again, strictly speaking,
 the linear function $h_Q$  amounts to a function
$y \to \left< \tau_Q, y\right> $ on $\power{V}$ for some $\tau_Q \in \power{V}$ and the prover answers that $\tau_Q$ is the
assignment to the  variables he received. However, it
is more convenient to view the function $h_Q$ itself as the answer.

\subsection{The success probability of the provers}
We now show that the provers' strategy succeeds with good probability.  Let
%
%
 $U, V, [x_1,\ldots, x_q],$ $x_i \in \power{V}$  %
be the provers' questions and $Q = {\sf Span}(x_1,\ldots, x_q)$.
We already observed that with probability
at least $\frac{\delta}{2}$,
 $U \in {\cal U}_{\sf good} $ and with probability $\frac{\alpha(\ell)}{2}$, $Q$ is both lucky and smooth (from the first
prover's perspective). Assume that all these properties hold. Then the answer of the first prover is a global function
$g_Q : \power{U} \to \power{}$ that satisfies the side condition on $H_U$, and
\[
  \Prob{L \subseteq \power{U}, ~{\sf dim}(L)=\ell}{g_Q|_{L\xor H_U}  \in F[L\xor H_U]  ~|  ~Q \subseteq L}\geqslant C\,.%
\]
Since $Q$ is smooth, by \expref{Lemma}{lemma:covering_cond_markov}, the uniform distribution on $\ell$-spaces in $\power{U}$ containing $Q$ is $(\sqrt{\beta} k^{\frac{1}{4}}\cdot 2^{\ell+5})$-close in statistical distance
to the distribution that chooses a question $V \subseteq U$ to the second prover and then chooses uniformly an $\ell$-space in $\power{V}$ containing $Q$. By setting the
parameters $\beta, k$ appropriately, we can assume that this statistical distance is at most $\frac{C}{2}$ and
conclude from the above inequality:
\[
  \Prob{V, ~L \in {\cal L}_V}{g_Q|_{L\xor H_U}  \in F[L\xor H_U]  ~|  ~Q \subseteq L}\geqslant \frac{C}{2}\,.%
\]
By an averaging argument, with probability at least $\frac{C}{4}$ over the choice of question $V$ to the second prover,
\[
  \Prob{L \in {\cal L}_V}{g_Q|_{L\xor H_U}  \in F[L\xor H_U]  ~|  ~Q \subseteq L}\geqslant \frac{C}{4}\,.%
\]
Fix any such good choice of question $V$.
Note that the assignment $F_V[L]$ to the $\ell$-spaces $L$ of the second prover is precisely the restriction of the assignment $F[L\xor H_U]$ to
the $\ell$-spaces of the first prover. Letting $h^*_Q: \power{V}\to\power{}$ to be the restriction of $g_Q:\power{U}\to \power{}$ to $\power{V}$,
we can rewrite the inequality above as:
\[
  \Prob{L \in {\cal L}_V}{h^*_Q|_L  \in F_V[L]  ~|  ~Q \subseteq L}\geqslant \frac{C}{4}\,.%
\]
Thus the function $h^*_Q$ satisfies \expeqref{Condition}{eqn:prover2-strat} and is a legitimate candidate for the second prover's answer. By \expref{Theorem}{thm:list_dec_upper}, the number of functions $h_Q$ satisfying \expeqref{Condition}{eqn:prover2-strat} is at most $\frac{jC/4}{(C/4)^2-j\cdot 2^{q-\ell}} \leq
 \frac{8 j}{C}$ for a large enough choice of $\ell$. When the second prover does pick $h^*_Q$ as the answer, both the provers' answers are
 consistent ($h^*_Q$ being a restriction of $g_Q$) and the provers succeed. Their overall success probability is at least
\[
\frac{\delta}{2} \cdot \frac{\alpha(\ell)}{2} \cdot \frac{C}{4} \cdot \frac{C}{8j} =
\frac{\delta ~\alpha(\ell) ~C^2}{128 ~j}\,,%
\]
thus completing the proof of \expref{Lemma}{lem:main_soundness} for $p(j,\delta,\ell) = \frac{\delta ~\alpha(\ell) ~C(j,\delta)^2}{128 ~j}$.


\section{The case \texorpdfstring{$j=1$}{j=1} of \expref{Hypothesis}{conj:main2}}
\label{section:j_1}
In this section it is shown that \expref{Hypothesis}{conj:main2},
in the case $j=1$, follows from \expref{Hypothesis}{conj:gras_expand} without the need of zoom-ins.

\smallskip
Let $f:S \to [2^\ell], ~{\sf Density}(S)=\delta$ be the given $(1,\delta)$-assignment to
the Grassmann graph $G(V=\{0,1\}^n, \ell)$, meaning, for any two $\ell$-spaces
$L_1, L_2 \in S$ such that ${\sf dim} (L_1 \cap L_2) = \ell-1$, we have the consistency $f[L_1]_{L_1 \cap L_2} = f[L_2]_{L_1 \cap L_2}$.
We intend to show, using \expref{Hypothesis}{conj:gras_expand}, that there is a global
linear function $g: \{0,1\}^n \to \{0,1\}$ such that $g|_L = f[L]$ for a fraction $C = C(\delta)$ of the $\ell$-spaces $L$.

\smallskip
Here is the idea. Fix an integer $b = \frac{\ell}{10}$.
Using \expref{Hypothesis}{conj:gras_expand}, we conclude rather easily, that for a 
constant fraction of
pairs $L_1, L_2 \in S$ such that  ${\sf dim} (L_1 \cap L_2) = b$, we still have  
$f[L_1]_{L_1 \cap L_2} = f[L_2]_{L_1 \cap L_2}$.
This enables us to assign linear functions to
all %
$b$-dimensional
spaces that have a good agreement with the given assignment to
the %
$\ell$-dimensional spaces.
In other words, this assignment passes
 the ``$\ell$-space vs $b$-space linearity test''\footnote{Analogous to the ``line vs point low-degree test.''}
with good probability. Using a Fourier analytic approach,\footnote{As opposed to
 the rather involved algebraic (and/or combinatorial) analysis of the ``line vs point'' and ``plane vs plane'' low-degree test in \cite{AS98, RS97}. }
we are able to show a soundness guarantee for the ``$\ell$-space vs $b$-space linearity test,''
implying the existence of a desired global linear function. A formal proof appears below.
The analysis of the linearity test is presented in \expref{Section}{appendix:linearity_testing} which might be of independent interest.


\smallskip

Let $b=\frac{\ell}{10}$. For a $b$-dimensional space $B\subseteq V$, let ${\sf Zoom}_B$ be the (lower order) Grassmann graph
induced on the set of vertices $\{L | L \in G(V,\ell), B \subseteq L\}$ (see \expref{Fact}{fact:isomorphic}).
Let ${\sf Zoom}_B[S] = S \cap {\sf Zoom}_B$
be the set of vertices in $S$ that contain $B$,
and let ${\sf Density} ( {\sf Zoom}_B[S] )$ be its density inside ${\sf Zoom}_B$. Clearly,
 $$ \Expect{B\subseteq V,\ \dimm(B)=b}{ {\sf Density} ( {\sf Zoom}_B[S] )  } =
   {\sf Density} (S) =  \delta\,.$$%
By an averaging argument, ${\sf Density} ( {\sf Zoom}_B[S] ) \geq \frac{\delta}{2}$ for at least $\frac{\delta}{2}$ fraction of $b$-spaces $B$; denote by
${\cal B}$ the set of all such ``good'' $b$-spaces.
Fix any $B \in {\cal B}$. Note that there are $2^b$ different $\mathbb{F}_2$-valued linear functions on $B$. Partition ${\sf Zoom}_B[S]$ into classes ${\cal C}_1,\ldots,
{\cal C}_{2^b}$ according to the restriction of $f[L]|_B$ for $L \in {\sf Zoom}_B[S]$. We observe that for any edge $(L, L')$ of the Grassmann graph inside
${\sf Zoom}_B[S]$, the linear functions $f[L], f[L']$ agree on $L_1 \cap L_2 \supseteq B$ and hence the edge is inside one of the partitions
${\cal C}_i$. Since ${\sf Density} ( {\sf Zoom}_B[S] ) \geq \frac{\delta}{2}$, \expref{Hypothesis}{conj:gras_expand} implies that there is a connected component ${\cal C}$ of density $\geq \eps$ in  ${\sf Zoom}_B[S]$ and as observed, ${\cal C} \subseteq {\cal C}_{i_0}$ for some $1 \leq i_0 \leq 2^b$. Let $h[B]$ denote the linear function on $B$ that equals the common
function $f[L]|_B$ over $L \in {\cal C}_{i_0}$. This gives an assignment $h: {\cal B} \to [2^b]$ of linear functions to $b$-spaces. From
the discussion, if a pair $(B, L), B \subseteq L$ of a $b$-space and a $\ell$-space is chosen at random from $V = \{0,1\}^n$, then
 $$ \Prob{B \subseteq L \subseteq V}{ f[L]|_B = h[B] } \geq \frac{\delta}{2} \cdot \eps\,,$$%
where $B \in {\cal B}$ with probability at least $\frac{\delta}{2}$ and then $L$ is in the ``large'' connected component of ${\sf Zoom}_B[S]$
 with probability at least $\eps$.
That is, the ``$\ell$-space vs $b$-space test'' succeeds with probability $\geq \frac{\delta}{2} \cdot \eps$. \expref{Theorem}{thm:linearity_dec} now implies that there
is a global linear function $g: V \to \{0,1\}$ that agrees with at least $C$ fraction of the $L$-spaces in $G(V,\ell)$.


\section{Transitivity of \texorpdfstring{$G_{2:2}$}{G\_{2:2}}}
\label{appendix:transitive}
In this section,  we show that the game  $G_{2:2}$ constructed in \expref{Section}{section:construction} is transitive as per
\expref{Definition}{def:trans_game}.

\begin{lemma}\label{lemma:transitive}
Suppose $L_1\xor H_{U_1}, L_2\xor H_{U_2}$ have a $1$-to-$1$ constraint between them in $G_{2:2}$, and
$L_2\xor H_{U_2}, L_3\xor H_{U_3}$ have a $1$-to-$1$ or a $2$-to-$2$ constraint between them.
Then there is a constraint
between $L_1\xor H_{U_1}, L_3\xor H_{U_3}$, and it is  $1$-to-$1$ or a $2$-to-$2$ depending on whether
the constraint between $L_2\xor H_{U_2}, ~L_3\xor H_{U_3}$ is $1$-to-$1$ or $2$-to-$2$.
%
%
\end{lemma}
\begin{proof}
Since there is $1$-to-$1$ constraint between $L_1\xor H_{U_1}, L_2\xor H_{U_2}$, we have
\begin{equation} \label{eqn:tr-1}
L_1+ H_{U_1}+ H_{U_2} = L_2+ H_{U_1}+ H_{U_2}\,.%
\end{equation}
We first consider the case when the constraint between $L_2\xor H_{U_2}, ~L_3\xor H_{U_3}$ is also $1$-to-$1$. This gives
\begin{equation}\label{eqn:tr-2}
L_2+ H_{U_2}+ H_{U_3} = L_3+ H_{U_2}+ H_{U_3}\,.%
\end{equation}
Combining the above equations gives (``add'' $H_{U_1}$ to \expeqref{Equation}{eqn:tr-2} and do a ``substitution'' using \expeqref{Equation}{eqn:tr-1})
%
\begin{equation}\label{eqn:tr-4}
L_1 + H_{U_1} +  H_{U_2} + H_{U_3} = L_3 + H_{U_1} +  H_{U_2} + H_{U_3}\,.%
\end{equation}
Now we would like to ``remove'' $H_{U_2}$ from both the sides so as to obtain $L_1 + H_{U_1} +  H_{U_3} = L_3 + H_{U_1} +  H_{U_3}$ and
implying that there is a $1$-to-$1$ constraint between $L_1\xor H_{U_1}$ and $L_3\xor H_{U_3}$. This ``removal'' can be done for the following
reason. Write $H_{U_2} = A \xor B, ~A \cap B = \set{0}$ where (a) $A$ is the span of all vectors $v_e$ such that the equation $e$ occurs in $U_2$, but also occurs
 in $U_1$ or  $U_3$, and hence $A \subseteq H_{U_1} + H_{U_3}$. (b) $B$ is the span of all vectors $v_e$ such that the equation $e$ occurs
in $U_2$, but not in $U_1$ nor in $U_3$. Any such equation $e$ shares at most one variable with $U_1$ and at most one variable with $U_3$ and no
variable with $U_2 \setminus e$. Hence there is a variable that is ``private'' to $e$, meaning it does not occur in $U_1 \cup U_3 \cup (U_2 \setminus e)$.
In particular, the ``private'' variables of the equations contributing to $B$
are distinct. Thus the intersection of
$B$ and $L_1+ L_3+ H_{U_1}+ H_{U_3} \subseteq \power{U_1}\xor \power{U_3}$
is $\set{0}$. To summarize, we can write \expeqref{Equation}{eqn:tr-4} as
\[
L_1 + \left( H_{U_1} + H_{U_3} +  A\right) + B  = L_3 + \left( H_{U_1} + H_{U_3} + A\right)  + B\,,%
\]
which simplifies to
\[
\left( L_1 + H_{U_1} + H_{U_3}\right)  + B  = \left( L_3 +  H_{U_1} + H_{U_3}\right)  + B\,,%
\]
and we can now safely ``remove'' $B$, using \expref{Fact}{fact:aux_3}.

\medskip
We now consider the case when the constraint between $L_2\xor H_{U_2}, ~L_3\xor H_{U_3}$ is $2$-to-$2$. We have \expeqref{Equation}{eqn:tr-1} as before,
but instead of \expeqref{Equation}{eqn:tr-2}, we now have
\begin{equation} \label{eqn:tr-3}
{\sf dim}(L_2+ H_{U_2}+ H_{U_3}) = {\sf dim}(L_3+ H_{U_2}+ H_{U_3}) = {\sf dim}(L_2+ L_3+ H_{U_2}+ H_{U_3})-1\,.%
\end{equation}
We claim that one can ``add'' $H_{U_1}$ to all three ``sums'' in \expeqref{Equation}{eqn:tr-3}. Arguing as earlier, one can write $H_{U_1} = A' \xor B', ~A' \cap B' = \{0\}$ where
$A' \subseteq H_{U_2}+ H_{U_3}$ and $B'$ is linearly independent of $\power{U_2}+ \power{U_3}$. Thus ``adding'' $H_{U_1}$ to all three ``sums,'' increases
the dimension of each ``sum'' by precisely ${\sf dim}(B')$. Thus
\[
{\sf dim}(L_2+ H_{U_1} + H_{U_2}+ H_{U_3}) = {\sf dim}(L_3+ H_{U_1} + H_{U_2}+ H_{U_3}) = {\sf dim}(L_2+ L_3+ H_{U_1} + H_{U_2}+ H_{U_3})-1\,.%
\]
Using \expeqref{Equation}{eqn:tr-1} and ``substituting,'' we get
\[
{\sf dim}(L_1+ H_{U_1} + H_{U_2}+ H_{U_3}) = {\sf dim}(L_3+ H_{U_1} + H_{U_2}+ H_{U_3}) = {\sf dim}(L_1+ L_3+ H_{U_1} + H_{U_2}+ H_{U_3})-1\,.%
\]
Now, arguing as earlier again, we ``remove'' $H_{U_2}$ from all three ``sums.'' One can write $H_{U_2} =  A \xor B, ~A \cap B = \{0\}$ where
$A \subseteq H_{U_1}+ H_{U_3}$ and $B$ intersects $L_1+ L_3+ H_{U_1} + H_{U_3}$ only at $\set{0}$. Thus ``removing'' $H_{U_2}$ from  all three ``sums,'' decreases
the dimension of each ``sum'' by precisely ${\sf dim}(B)$. Thus
\[
{\sf dim}(L_1+ H_{U_1} + H_{U_3}) = {\sf dim}(L_3+ H_{U_1} + H_{U_3}) = {\sf dim}(L_1+ L_3+ H_{U_1} + H_{U_3})-1\,,%
\]
implying that there is a $2$-to-$2$ constraint between $L_1 \xor H_{U_1}$ and $L_3 \xor H_{U_3}$.
\end{proof}




\begin{lemma}\label{lemma:constraint_composition}
Let $s_1 = L_1\xor H_{U_1},  s_2 = L_2\xor H_{U_2}, s_3 = L_3\xor H_{U_3}$
be vertices in $G_{2:2}$ such that there is $1$-to-$1$ constraint between $(s_1, s_2)$ and a constraint between $(s_2, s_3)$.
Then the  constraint between $(s_1, s_3)$ (as guaranteed by \expref{Lemma}{lemma:transitive}) is a composition of the constraints
between $(s_1, s_2)$ and $(s_2, s_3)$.

Specifically, if linear functions (respecting relevant side
conditions) $f$ on $L_1 \xor H_{U_1}$, $g$ on $L_2 \xor H_{U_2}$, and $h$ on $L_3\xor H_{U_3}$ are such that
  $(f,g)$ satisfy $(s_1, s_2)$ and
  $(g,h)$ satisfy $(s_2, s_3)$,
then $(f,h)$ satisfy $(s_1, s_3)$.
\end{lemma}

\begin{proof} In the following, whenever we construct a linear function on a certain space, it will always respect the side condition contained in that
space. Since $(g,h)$ satisfy the constraint $(s_2, s_3)$, there is a linear function $\beta$ on $W = L_2 + L_3 + H_{U_2} + H_{U_3}$ that respects
side conditions on $H_{U_2}$ and $H_{U_3}$ and
  $$g = \beta|_{L_2 \xor H_{U_2}},  \quad \quad h = \beta|_{L_3 \xor H_{U_3}}\,.$$%
Let $Z = W + H_{U_1}$ and extend the linear function $\beta$ on $W$ uniquely to a linear function $\gamma$ on $Z$ so as to respect the side condition $H_{U_1}$.  This is possible because
every equation in $U_1$ that does not appear in $U_2$ or
$U_3$ has a ``private variable,'' as in the proof of the previous lemma. We note that
 $$\gamma|_{L_2 \xor H_{U_2}} = (\gamma|_W)|_{L_2 \xor H_{U_2}} = \beta|_{L_2 \xor H_{U_2}} = g\,,$$%
  $$\gamma|_{L_3 \xor H_{U_3}} = (\gamma|_W)|_{L_3 \xor H_{U_3}} = \beta|_{L_3 \xor H_{U_3}} = h\,.$$%
Since there is a $1$-to-$1$ constraint $(s_1, s_2)$, we have
  $$L_1 +  H_{U_1} + H_{U_2} =  L_2 + H_{U_1} + H_{U_2}\,.$$%
It therefore holds  that
 $$ Z =  L_2 + L_3 + H_{U_1} + H_{U_2} + H_{U_3}  \supseteq L_2 + H_{U_1} + H_{U_2} = L_1 + H_{U_1} + H_{U_2} \supseteq L_1 \xor H_{U_1}\,.$$%
Since $\gamma$ is an assignment on $Z$, by \expref{Remark}{remark:useful_222}, $\gamma|_{L_1 \xor H_{U_1}}$ and $\gamma|_{L_2 \xor H_{U_2}} = g$
satisfy the constraint $(s_1, s_2)$. However $(f,g)$ is supposed to satisfy this $1$-to-$1$ constraint, and hence we must have $f = \gamma|_{L_1 \xor H_{U_1}}$.
Now we have
  $f = \gamma|_{L_1 \xor H_{U_1}}$ and $h = \gamma|_{L_3 \xor H_{U_3}}$ and by \expref{Remark}{remark:useful_222}, $(f,h)$ satisfy the constraint $(s_1, s_3)$.
\end{proof}

\section{Covering property}\label{appendix:covering}
The goal of this section is to prove the covering property, namely
\expref{Lemmas}{lemma:covering_normal} \expref{and}{lemma:covering_cond_markov}.
\subsection*{Proof of \expref{Lemma}{lemma:covering_cond_markov} from \expref{Lemma}{lemma:covering_normal}}

Let $U, V, k, \beta, \ell, {\cal L}, {\cal L}', q, Q, {\cal L}_Q, {\cal L}'_Q$ be as in \expref{Definition}{def:covering_dist} and \expref{Lemmas}{lemma:covering_normal}
\expref{as~in}{lemma:covering_cond_markov}.  %
Let ${\cal Q}, {\cal Q}'$ be
distributions over $q$-dimensional subspaces of $\power{U}$ that are analogous to ${\cal L}, {\cal L}'$, respectively %
(\ie, as in \expref{Definition}{def:covering_dist}, with
parameter $q$ instead of $\ell$). It is easily observed that an equivalent way to sample
%
%
from ${\cal Q}$ and ${\cal Q}'$) is to sample an $\ell$-space $L$
from ${\cal L}$ and ${\cal L}'$), respectively,  %
and then sample a uniformly random $q$-dimensional subspace of $L$.
We stress that ${\cal Q}$ and ${\cal L}$ are
uniform distributions on $q$-dimensional and $\ell$-dimensional subspaces
of $\power{U}$, respectively.  %
We have the sequence of arguments
\begin{eqnarray*}
\Expect{Q \sim {\cal Q}}{{\sf SD}({\cal L}_Q, {\cal L}'_Q) } & = & \sum_{Q} \Prob{}{{\cal Q}=Q} \sum_{L \supseteq Q} \left| \Prob{}{{\cal L}_Q = L} - \Prob{}{{\cal L}'_Q = L} \right| \\
 & = & \sum_{Q} \sum_{L \supseteq Q} \left|  \Prob{}{{\cal Q}=Q} \cdot \Prob{}{{\cal L}_Q = L} - \Prob{}{{\cal Q}=Q} \cdot \Prob{}{{\cal L}'_Q = L}       \right| \\
 & \leq &  \sum_{Q} \sum_{L \supseteq Q} \left|  \Prob{}{{\cal Q}=Q} \cdot \Prob{}{{\cal L}_Q = L} - \Prob{}{{\cal Q'}=Q} \cdot \Prob{}{{\cal L}'_Q = L}       \right| + \\
 &  &  \sum_{Q} \sum_{L \supseteq Q} \left|  \Prob{}{{\cal Q}'=Q} \cdot \Prob{}{{\cal L}'_Q = L} - \Prob{}{{\cal Q}=Q} \cdot \Prob{}{{\cal L}'_Q = L}       \right|  \\
 & = &  \sum_{L} \left| \Prob{}{{\cal L}=L} - \Prob{}{{\cal L}' = L} \right|  + \\
 & & \sum_Q \left| \Prob{}{{\cal Q}'=Q} - \Prob{}{{\cal Q}=Q} \right| \sum_{L \supseteq Q} \Prob{}{{\cal L}'_Q = L} \\
 & = & {\sf SD}({\cal L}, {\cal L}') + {\sf SD}( {\cal Q}, {\cal Q}')\,,%
\end{eqnarray*}
where we used triangle inequality and the fact that sampling $Q \in {\cal Q}$ and then $L \sim {\cal L}_Q$ is equivalent to sampling $L \sim {\cal L}$ and
similarly, sampling $Q \in {\cal Q}'$ and then $L \sim {\cal L}'_Q$ is equivalent to sampling $L \sim {\cal L}'$.  Now \expref{Lemma}{lemma:covering_normal} upper bounds both
${\sf SD}({\cal L}, {\cal L}')$ and $ {\sf SD}( {\cal Q}, {\cal Q}')$ by $\beta \sqrt{k} \cdot 2^{\ell+4}$ and hence
  $$ \Expect{Q \sim {\cal Q}}{{\sf SD}({\cal L}_Q, {\cal L}'_Q )} \leq  \beta \sqrt{k} \cdot 2^{\ell+5}\,.$$%
\expref{Lemma}{lemma:covering_cond_markov}  now follows by
Markov's inequality.%
%
%



\subsection*{Proof of \expref{Lemma}{lemma:covering_normal} }

We recall that $U, |U|=3k$ is a set of $3k$ variables in $k$ equations. A subset $V \subseteq U$ is chosen by choosing independently
for each equation, one of the variables in the equation with probability $\beta$ and all three variables in the equation with
probability $1-\beta$. The expected size of $V$ is $3k-2\beta k$ and except with probability $2^{-\Omega(k)}$, we have $|V| \geq 2k$.

We note that choosing a uniformly random $\ell$-subspace $L$ of $\power{U}$ (resp. $\power{V}$) is equivalent to
choosing uniformly a sequence of points $x_1, \ldots, x_\ell$ in $ \power{U}$ (resp. $\power{V}$) that are
linearly independent and letting $L = {\sf Span}(x_1,\ldots, x_\ell)$.
Since a uniformly random and independent sequence of
points $x_1, \ldots, x_\ell$ in $ \power{U}$ (resp. in $\power{V}$) is linearly independent except with probability
$\leq 2^{\ell - {\sf dim}(U)}$ (resp. $\leq 2^{\ell - {\sf dim}(V)}$, see \expref{Fact}{fact:aux_1}), we might as well focus on such sequences of points. It is thus enough to bound the
statistical distance between distributions ${\cal D}, {\cal D}'$ over $(\power{U})^\ell$ sampled as:
\begin{itemize}
\item ${\cal D}$: Choose uniformly and independently $x_1, \ldots, x_\ell \in \power{U}$.
\item ${\cal D}'$: Choose $V \subseteq U$, choose uniformly and independently $x'_1, \ldots, x'_\ell \in \power{V}$ and regard them as
      points in $\power{U}$ (by appending $0$ in coordinates $U \setminus V$).
\end{itemize}
We now observe that since the process of choosing $V \subseteq U$ is independent over the $k$ equations,
${\cal D} = {\cal S}^k$ and ${\cal D}' = {\cal S}'^k$ where ${\cal S}, {\cal S}'$ are the ``basic'' distributions
exactly as above, but with $k=1, |U|=3$. A bound on the statistical distance between ${\cal D}, {\cal D}'$ now follows in the
same manner as in \cite[Lemma 3.1]{KS}, by bounding the Hellinger distance between ${\cal S}, {\cal S}'$, using the
multiplicativity of the Hellinger distance to bound the Hellinger distance between ${\cal D}, {\cal D}'$ and finally,
bounding the Hellinger distance in terms of the statistical distance. We observe how a bound on Hellinger distance between ${\cal S},
{\cal S}'$ also follows already from the proof of \cite[Lemma 3.1]{KS}. Re-writing the sampling process for
${\cal S}, {\cal S}'$ for convenience (this is the special case $k=1, |U|=3$):
\begin{itemize}
\item ${\cal S}$: Choose uniformly at random $x_1,\ldots, x_\ell \in \{0,1\}^3$.
\item ${\cal S}'$:  With probability $1-\beta$, choose uniformly at random $x_1,\ldots, x_\ell \in \{0,1\}^3$. Otherwise:

$\quad$ ~Choose uniformly at random $b_1, \ldots, b_\ell \in \power{}$. Output with probability $\frac{\beta}{3}$ each,
      $$b_100, \ldots, b_\ell00, \quad \mbox{or} \quad  0b_10, \ldots, 0b_\ell0, \quad \mbox{or} \quad 00b_1, \ldots, 00b_\ell
      \,.$$%
\end{itemize}
The distributions ${\cal S}, {\cal S}'$ are over $(\power{3})^\ell$ which is
equivalent to $\Sigma^3$ with $\Sigma=\power{\ell}$, where the
%
first, second, and third coordinate %
in $\Sigma^3$ %
correspond %
to the concatenation of the
%
first, second, and third coordinate, respectively,  %
in the $\ell$ output triplets.
%
%
%
%
%
%
Note that $0^\ell \in \Sigma$.
Denoting the uniform distribution over $\Sigma$ by ${\sf Uniform}(\Sigma)$, it is seen that
  $$ {\cal S} =  \left( {\sf Uniform}(\Sigma), {\sf Uniform}(\Sigma), {\sf Uniform}(\Sigma) \right)\,,$$%
\ie, three independent and uniform copies of $\Sigma$, whereas,
\begin{eqnarray*}
 {\cal S}' & = & (1-\beta) \left( {\sf Uniform}(\Sigma), {\sf Uniform}(\Sigma), {\sf Uniform}(\Sigma) \right)  + \\
         & & \frac{\beta}{3}
        \left(  {\sf Uniform}(\Sigma), 0^\ell, 0^\ell  \right)  + \frac{\beta}{3}
        \left(  {0^\ell, \sf Uniform}(\Sigma), 0^\ell  \right)  + \frac{\beta}{3}
        \left( 0^\ell, 0^\ell, {\sf Uniform}(\Sigma) \right)\,.%
        \end{eqnarray*}
With this viewpoint,  the Hellinger distance between ${\cal S}, {\cal S}' $ is calculated to be at most $4 \beta^2 |\Sigma|^2$ in the
proof of \cite[Lemma 3.1]{KS}.  The statistical distance between ${\cal D}, {\cal D}'$ is then at most $16 \beta \sqrt{k}\cdot |\Sigma|$.

%


%
\appendix
\section*{Appendix}

\section{Reduction from \texorpdfstring{$2$-to-$2$}{2-to-2} Games
to the Independent Set Problem}\label{sec:vc-is}

In this section, we present a reduction from the (Transitive) $2$-to-$2$ Games problem to the Independent Set problem, proving
\expref{Theorem}{thm:j-delta-vc}.
The reduction is along the lines of \cite{DS05, Khot02, KR08}, using the Biased Long Code
and analytic theorems of Russo, Margulis and Friedgut, introduced in \cite{DS05}. Some care is required to handle the
transitivity feature.


\subsection{Biased Long Code}\label{appendix:blc}
While the Biased Long Code can be viewed as an encoding scheme, it is more convenient to take a combinatorial
view and treat it as a weighted Kneser graph. The valid codewords then amount to certain large (in fact the largest) independent sets
in this graph. The analysis of the Biased Long Code amounts to a structural theorem about independent sets of
moderately large (= linear) size.


\begin{definition} For a bias parameter $p \in (0, 1)$ and alphabet $\Sigma$,
the vertex set of weighted Kneser graph  $G_p[\Sigma]$ is ${\cal P}(\Sigma)$, the family of all subsets of $\Sigma$.
The weight of a vertex $A \subseteq \Sigma$ is $\mu_p(A)=p^{\card{A}}(1-p)^{|\Sigma|-\card{A}}$.
The edge set is $\sett{(A,B)}{A, B \subseteq\Sigma, ~A\cap B=\emptyset}$.
\end{definition}

It can be shown easily that the largest independent sets in $G_p[\Sigma]$ have weight $p$. These are precisely
the sets $I_{\sigma_0} = \{ F \mid \sigma_0 \in F \}$ for any fixed $\sigma_0 \in \Sigma$.

\begin{definition} For a set family ${\cal F} \subseteq {\cal P}(\Sigma)$, let $\mu_p({\cal F})$ denote its
weight under $\mu_p$. Let $A \sim \mu_p$ denote the process of picking a set $A \subseteq \Sigma$
according to the distribution $\mu_p$.
For a fixed element $\sigma \in \Sigma$, let ${\sf Infl}_\sigma({\cal F})$ denote its
influence on the family ${\cal F}$ defined as
 $$ {\sf Infl}_\sigma({\cal F})  =  \Prob{A \sim \mu_p}{\mbox{~Exactly one of the pair} ~A ~\mbox{and}  ~A \Delta \{\sigma\} ~\mbox{is in} ~{\cal F}}\,. $$%
 The average sensitivity of a family ${\sf as}_p({\cal F})$ is the sum of all influences, \ie,
 $${\sf as}_p({\cal F}) = \sum_{\sigma \in \Sigma}  {\sf Infl}_\sigma({\cal F})\,.   $$%
\end{definition}



\subsection{The reduction}

Let $G=(V,E,\Phi,\Sigma)$ be the instance of a Transitive $2$-to-$2$ Game as in \expref{Conjecture}{conj:j-delta}. The parameters $j$ and $\delta$ therein will be chosen later.
The Independent Set instance $G'=(V',E')$ is defined as follows. Set the parameter $p=1-\frac{1}{\sqrt{2}}-\delta$.
The vertex set of the instance is
\[
V'= \sett{(x,A)}{ x \in V, ~A\subseteq \Sigma}\,.%
\]
The weight of the vertex $(x,A)$ is  $\frac{1}{|V|} \cdot \mu_p(A)$, so that the total weight of all the vertices is $1$. The edge set is
\[
E' = \sett{((x_1,A_1),(x_2,A_2))}{ (x_1,x_2)\in E \ \land \ \forall  ~\sigma_1 \in A_1,
 \ \sigma_2 \in A_2,  (\sigma_1, \sigma_2) \not\in \Phi(x_1, x_2)}\,.%
\]

In words, there is a cloud of vertices for every $x \in V$.
For every constraint $(x_1, x_2) \in E$, there are cross edges between the
respective clouds.  %
There is an edge between $(x_1,A_1),(x_2,A_2)$ if there is no pair of
colors in the sets $A_1,A_2$ that satisfy the constraint on $(x_1,x_2)$. \footnote{ One could add edges inside each cloud according to the
 Kneser graph. The reduction does not need it though.}


\subsection{Completeness}
Let ${\cal C}: X \to \Sigma$ be a $(1,1-\delta)$-coloring of the game  $G=(V,E,\Phi,\Sigma)$ where $X \subseteq V$, $|X|=(1-\delta)|V|$. The coloring
satisfies all the constraints inside $X$. Consider the set of vertices in $G'(V',E')$,
\[
I\defeq\sett{(x,A)}{x\in X, ~ {\cal C}(x)\in A}\,.%
\]

Clearly, the set $I$ includes a weight $p$ of the vertices inside the cloud for every $x \in X$. Hence the weight of $I$ is
$(1-\delta) p \geq 1-\frac{1}{\sqrt{2}} - 2\delta $. We observe that
$I$ is an independent set. For every pair of vertices
 $(x,A),(x',A')\in I$, we show that there is no edge between them in $G'$. Since the coloring
${\cal C}$ satisfies the constraint $(x,x')$, we have $({\cal C}(x), {\cal C}(x')) \in \Phi(x,x')$.
By definition of the set $I$, we have  ${\cal C}(x) \in A,  ~{\cal C}(x') \in A'$.  Thus $A,A'$ contain a consistent pair of colors,
so there is no edge between $(x,A)$ and $(x',A')$.



\subsection{Soundness}\label{appendix:blc_analysis}

We begin by stating two auxiliary lemmas towards the soundness analysis. The relevance of the $2$-to-$2$-ness of the
constraints and the choice of $p \approx 1-\frac{1}{\sqrt{2}}$ is apparent from the statements of these lemmas.
Let $\Sigma$ and $\Gamma$ be alphabets such that $|\Gamma| = \frac{|\Sigma|}{2}$ and $\pi: \Sigma \to \Gamma$
be a $2$-to-$1$ map. For $F \subseteq \Sigma$, its projection $\pi(F) \subseteq \Gamma$ is defined naturally as
$\{\pi(\sigma) | \sigma \in F\}$. For a family ${\cal F} \subseteq {\cal P}(\Sigma)$, the projected family
$\pi({\cal F}) \subseteq {\cal P}(\Gamma)$ is defined naturally as $\{ \pi(F) | F \in {\cal F}\}$. For a
subset $H \subseteq \Gamma$, the set $\pi^{-1}(H), ~|\pi^{-1}(H)|= 2|H|$ is defined naturally as
$\{ \sigma | \sigma \in \Sigma, ~\pi(\sigma) \in H \}.$

\begin{lemma}\label{lem:aux-1}  For $q \in (0,1)$, ~$\mu_{1-(1-q)^2} ( \pi( {\cal F} ) ) \geq  \mu_q ( {\cal F} ). $
\end{lemma}
\begin{proof} For every $H \subseteq \Gamma$, we define $\pi^{\uparrow}(H) = \{ F \subseteq \Sigma | \pi(F) = H \}$. We observe that
\begin{itemize}
\item $\mu_{1-(1-q)^2} ( H ) =   \mu_q ( \pi^{\uparrow} (H)  ). $
\item The families $\pi^{\uparrow}(H)$ over all $H \subseteq \Gamma$
%
form a  %
%
partition of the family ${\cal P} (\Sigma)$.
\end{itemize}
The lemma follows by noting that
\[
 \mu_{1-(1-q)^2} ( \pi( {\cal F} ) ) = \sum_{H \in \pi( {\cal F} )} \mu_{1-(1-q)^2}(H) =
 \sum_{H \in \pi( {\cal F} )} \mu_q ( \pi^{\uparrow} (H)  ) \geq  \sum_{H \in \pi( {\cal F} )} \mu_q ( \pi^{\uparrow} (H) \cap {\cal F} )  =
 \mu_q ({\cal F})\,.   \qedhere%
\]
\end{proof}

\begin{lemma}\label{lem:aux-2}  Let ${\cal F}\subseteq {\cal P}(\Sigma), ~{\cal F}' \subseteq {\cal P}(\Sigma')$ be two families, each of weight strictly larger than
$\frac{1}{2}$ under the distribution $\mu_q$ with $q = 1-\frac{1}{\sqrt{2}}$. Let $\pi: \Sigma \to \Gamma, ~\pi': \Sigma' \to \Gamma$ be
$2$-to-$1$ maps (so $|\Sigma| = |\Sigma'| = 2|\Gamma|$).
Then there exist $F \in {\cal F},
~F' \in {\cal F}'$ such that $\pi(F) \cap \pi'(F') = \emptyset$.
\end{lemma}
\begin{proof} We note that $1-(1-q)^2 = \frac{1}{2}$ and from \expref{Lemma}{lem:aux-1}, $\mu_{\frac{1}{2}} (\pi({\cal F})) \geq \mu_q ({\cal F}) >
\frac{1}{2}$ and similarly $\mu_{\frac{1}{2}} (\pi'({\cal F}')) > \frac{1}{2}$. Thus $\pi({\cal F})$ and  $\pi'({\cal F}')$ are families,
each containing more than half (in the usual counting sense) of the sets from ${\cal P}(\Gamma)$. Hence there must exist $H \in \pi({\cal F}),
~H' \in \pi'({\cal F}')$ that are complements of each other and in particular $H \cap H' = \emptyset$.
\end{proof}



We now present the soundness analysis. Given
 a maximal independent set $I$ of weight at least $\eps$ in $G'$, we show how to construct
a $(j,\delta)$-coloring for $G=(V,E,\Phi,\Sigma)$ where $j,\delta$ depend only on $\eps$.
For every $x\in G$, consider the part of $I$ inside the cloud of $x$,
\[
{\cal F}_x=\sett{A}{A \subseteq \Sigma, ~(x,A)\in I}\,.%
\]
\begin{claim}
The family ${\cal F}_x$ is monotone.
\end{claim}
\begin{proof}
Otherwise, there are $A\subseteq B$ such that $A\in {\cal F}_x$ and $B\not\in{\cal F}_x$.
Then $I\cup\set{(x,B)}$ is an independent set larger than $I$,  contradicting  the maximality
of $I$.
\end{proof}

Since the independent set $I$ has weight $\eps$, by an averaging argument, there is a set
$X \subseteq V, ~|X| \geqslant \frac{\eps}{2} \cdot |V|$ such that $I$ includes a weight
$\geqslant \frac{\eps}{2}$ of vertices from the cloud of $x$, \ie, $\mu_p({\cal F}_x) \geqslant  \frac{\eps}{2}$
for $x \in X$.



\begin{theorem}[Russo -- Margulis ~\cite{Russo1982,Marg74}]
Suppose ${\cal F}$ is a monotone family. Then $\mu_q({\cal F})$ is an increasing function of $q$ and
\[
\frac{d\mu_q({\cal F})}{dq}={\sf as}_q({\cal F})\,.%
\]
\end{theorem}

\begin{claim}\label{clm:as}
There exists $p'\in (p,p+\delta)$ such that
\[
\Expect{x\in X}{{\sf as}_{p'}({\cal F}_x)} \leqslant \frac{1}{\delta}\,.%
\]
\end{claim}
\begin{proof}
By Lagrange's mean value theorem, it follows that there exists $p'\in (p,p+\delta)$ such that
\[
\Expect{x\in X}{{\sf as}_{p'}({\cal F}_x)}=
\frac{d}{dq}\Big( \Expect{x\in X}{\mu_{q}({\cal F}_x)}     \Big) \Big|_{q=p'} =
\frac
{
\Expect{x\in X}{\mu_{p+\delta}({\cal F}_x)}
-\Expect{x\in X}{\mu_{p}({\cal F}_x)}
}
{\delta}
\leqslant \frac{1}{\delta}\,.   \qedhere%
\]
\end{proof}
From \expref{Claim}{clm:as} and an averaging argument, there is a set $X' \subseteq X, ~|X'| \geq \frac{|X|}{2}$ such that
for all $x \in X'$, $~{\sf as}_{p'}({\cal F}_x)\leqslant \frac{2}{\delta}.$ A theorem of Friedgut states that families
with bounded average sensitivity are well-approximated by ``juntas.''

\begin{definition}[Junta]
A family ${\cal F}\subseteq {\cal P}(\Sigma)$ is called a $j$-junta, if there exists $J\subseteq \Sigma, |J|=j$ such that the membership of a set
$A$ in ${\cal F}$ is determined by only $A\cap J$.
\end{definition}



\begin{theorem}[Friedgut ~\cite{Friedgut98}]
There exists %
$C(q) \geq 1$ such that for every
${\cal F}\subseteq P(\Sigma)$ and %
every  %
accuracy parameter $\eta>0$,
there exists ${\cal F}'\subseteq P(\Sigma)$ that is a $j$-junta and
\begin{itemize}
\item $j = C(q)^{ {\sf as}_q ({\cal F})/\eta}. $
\item $\mu_{q}({\cal F}\Delta{\cal F}')\leqslant \eta$.
\end{itemize}
\end{theorem}

Fix $x\in X'$ and set $\eta\defeq \frac{\eps}{20}$.
Since ${\sf as}_{p'}({\cal F}_x) \leqslant \frac{2}{\delta}$,
it follows from Friedgut's Theorem that ${\cal F}_x$ is $\eta$-close to a $k$-junta with $k = C(p')^{2/(\delta \eta)}$. Let $J_x \subseteq \Sigma$
denote the set of elements  on which the junta depends. Clearly, the set-family that is a junta on $J_x$ and is closest to ${\cal F}_x$
 is the ``majority vote'' on each setting of $J_x$, namely
\[
[{\cal F}_x]_\half\defeq \sett{F\cup F'}{F\subseteq\Sigma\setminus J_x, ~F' \subseteq J_x, ~\Prob{A\subseteq \Sigma\setminus J_x, ~A \sim \mu_{p'}}{A\cup F'\in{\cal F}_x} > \half}\,.%
\]
The following claim shows that the family $[{\cal F}_x]_{\frac{3}{4}}$ is also close to ${\cal F}_x$ (and will be more useful to work with):
\[
[{\cal F}_x]_{\frac{3}{4}}\defeq \sett{F\cup F'}{F\subseteq\Sigma\setminus J_x, ~F'\subseteq J_x ~\Prob{A\subseteq \Sigma\setminus J_x, ~A\sim \mu_{p'}}{A\cup F'\in {\cal F}_x} > \frac{3}{4}}\,.%
\]

\begin{claim}$
\mu_{p'}({\cal F}_x \Delta [ {\cal F}_x ]_{\frac{3}{4}}   )\leqslant 5\eta.
$
\end{claim}
\begin{proof}
Let ${\cal F}^* \subseteq {\cal P}(J_x)$ be the family of subsets $F \subseteq J_x$ such that
\[
\half<\Prob{A \subseteq \Sigma\setminus J_x, ~A\sim \mu_{p'}}{A\cup F \in {\cal F}_x}\leqslant \frac{3}{4}\,.%
\]
Notice that for each such $F$  (a) at least  $\quarter$ (weighted) fraction of its extensions to $\Sigma$ are not in ${\cal F}_x$  (b)
each extension is in $[{\cal F}_x]_\half$ (c) no extension is in $[{\cal F}_x]_{\frac{3}{4}}$. Hence
\[
\quarter\cdot \Prob{F \subseteq J_x, ~F\sim\mu_{p'}}
{F \in {\cal F}^*}\leqslant
\mu_{p'}({\cal F}_x\Delta [{\cal F}_x]_\half)
\leqslant
\eta\,.%
\]
It follows that
\[
\mu([{\cal F}_x]_\frac{1}{2}\Delta [{\cal F}_x]_\frac{3}{4})
=
\Prob{F\subseteq J_x, ~F\sim\mu_{p'}}
{F \in {\cal F}^*} \leqslant 4\eta\,.%
\]
We finish the proof using the triangle inequality,
\[
\mu_{p'}({\cal F}_x\Delta [{\cal F}_x]_{\frac{3}{4}})
\leqslant
\mu_{p'}({\cal F}_x\Delta [{\cal F}_x]_{\frac{1}{2}})
+\mu_{p'}([{\cal F}_x]_\frac{1}{2}\Delta [{\cal F}_x]_\frac{3}{4})
\leqslant 5\eta\,.    \qedhere%
\]
\end{proof}

\begin{claim}\label{claim:inf_non_empty}
 $[{\cal F}_x]_{\frac{3}{4}}\neq \emptyset$.
\end{claim}
\begin{proof}
Using the triangle inequality, the previous claim, and that $\eta = \frac{\eps}{20}$,
\[
\mu_{p'}([{\cal F}_x]_{\frac{3}{4}}) \geqslant
\mu_{p'}({\cal F}_x)-\mu_{p'}({\cal F}_x \Delta [ {\cal F}_x ]_{\frac{3}{4}}   )
 \geqslant \frac{\eps}{2} - 5\eta
> 0\,.      \qedhere%
\]
\end{proof}

\begin{definition}
The extended junta $EJ(x)$ of $x$ is defined by
\[
EJ(x)= J_x ~\cup  ~\sett{\sigma \in \Sigma}{ {\sf Infl}_\sigma({\cal F}_x)\geqslant 2^{-10 k}}\,.%
\]
\end{definition}

We note that for $x \in X'$, ${\sf as}_{p'}({\cal F}_x)\leqslant \frac{2}{\delta}$ and since the average sensitivity is the sum of all influences,
$ |EJ(x)| \leq j = k +  \frac{2 \cdot 2^{10k} }{\delta}$. Our coloring to the game $G=(V,E,\Phi,\Sigma)$ will assign, to
every $x \in X'$, a set of at most $j$ colors $EJ(x)$. We now show that this is indeed a $(j,\delta)$-coloring. Firstly,
 $$|X'| \geq \frac{|X|}{2} \geq \frac{\eps}{4} |V|\geq \delta |V|$$
  as $\delta$ will be chosen accordingly. Secondly, we need
 to show that every constraint $(x_1, x_2)$ inside $X'$ is satisfied in the sense of \expref{Definition}{def:j,delta}.
 Fix any such constraint. It is a $2$-to-$2$ or a $1$-to-$1$ constraint. Our main
soundness lemma below takes care of the $2$-to-$2$ case, and the $1$-to-$1$ case then follows directly from the
transitivity of the game.


\begin{lemma}\label{lemma:there_is_constraint}
Suppose $x_1,x_2\in X'$ are such that $(x_1,x_2)$ is a $2$-to-$2$ constraint.
Then there exist consistent colors
for $x_1,x_2$ in their respective  %
extended juntas. \ie, there exist $\sigma_1  \in EJ(x_1), ~\sigma_2 \in EJ(x_2)$
such that $(\sigma_1, \sigma_2) \in \Phi(x_1, x_2)$.
\end{lemma}
\begin{proof}
It will be
convenient to think of the $2$-to-$2$ constraint in terms of a pair of $2$-to-$1$ maps
$\pi_1: \Sigma_1 \to \Gamma, \pi_2: \Sigma_2 \to \Gamma$. Here $\Sigma_1 = \Sigma_2 = \Sigma$ are
the same alphabet, but it will be convenient to think of them as separate.
A coloring $(a_1, a_2)$ to vertices $(x_1, x_2)$ satisfies the
$2$-to-$2$ constraint if and only if $\pi_1(a_1) = \pi_2(a_2)$. Assume towards a contradiction that
there is no pair of consistent colors in the extended juntas for $x_1$ and $x_2$.
The assumption can be stated as
 $$  \pi_1 ( EJ(x_1) ) \cap \pi_2 (EJ(x_2) ) = \emptyset\,. $$%
Note in particular that $J_{x_2} \subseteq EJ(x_2)$ and hence
\begin{equation} \label{eqn:assumption}
    \pi_1 ( EJ(x_1) ) \cap  \pi_2 ( J_{x_2} ) = \emptyset\,.%
\end{equation}
Our goal is to exhibit $F_1 \in {\cal F}_{x_1}, ~F_2 \in {\cal F}_{x_2}$ such that $\pi_1(F_1)
\cap \pi_2(F_2) = \emptyset$. We will zero in on such $F_1, ~F_2$ in progressive manner.
%
We consider the case of $F_1$, the other case being similar. We zero in on
a sequence of sets,  %
 $$ A_1  \subseteq B_1 \subseteq B_1  \subseteq  F_1\,, $$%
that are contained, respectively, %
in a %
progressively expanding ``universe of focus''
 $$ J_{x_1} \ \ \subseteq \ \  \pi_1^{-1} ( \pi_1( J_{x_1} ) )  \ \  \subseteq \  \
    \pi_1^{-1}( \pi_1 ( J_{x_1}  ) ) \cup  \pi_1^{-1} ( \pi_2  ( J_{x_2} ) )  \ \  \subseteq \ \ \Sigma_1\,.$$%
    We clarify that the set $ \pi_1^{-1} ( \pi_1( J_{x_1} ) ) $ is a superset of $J_{x_1}$
    and can have size up to $2 |J_{x_1}|$ since that map $\pi_1$ is $2$-to-$1$.
The weights (sizes) of set-families are with respect to $\mu_{p'}$, unless stated otherwise.
\begin{itemize}
\item Recalling the definition of $[{\cal F}_{x_1}]_{\frac{3}{4}}$ and using \expref{Claim}{claim:inf_non_empty},
    there is
      $A_1 \subseteq J_{x_1}$ such that at least $\frac{3}{4}$ of its extensions outside $J_{x_1}$
      are in ${\cal F}_{x_1}$.
\item We let $B_1 =  A_1  \cup \left(  \pi_1^{-1} ( \pi_1( J_{x_1} ) )  \setminus J_{x_1}  \right)$. Due to monotonicity of ${\cal F}_{x_1}$,
      at least $\frac{3}{4}$ of extensions of $B_1$ outside $  \pi_1^{-1} ( \pi_1( J_{x_1} ) ) $
      are in ${\cal F}_{x_1}$.
\item We now retain $B_1$ as is, but consider it as subset of enlarged universe $ \pi_1^{-1}( \pi_1 ( J_{x_1}  ) ) \cup  \pi_1^{-1} ( \pi_2  ( J_{x_2} ) ) $.
      The elements added to the enlarged universe, namely  $\pi_1^{-1} ( \pi_2  ( J_{x_2} ) )$
      are outside of $EJ(x_1)$ (using \expeqref{Equation}{eqn:assumption}), hence have influence at most $2^{-10k}$, and are at most $2k$ in number. %
%
%
%
%
      The fraction of extensions of
      $B_1$ outside $ \pi_1^{-1}( \pi_1 ( J_{x_1}  ) ) \cup  \pi_1^{-1} ( \pi_2  ( J_{x_2} ) ) $
      remains at least $$\frac{3}{4}- 2^{-10k} \cdot {p'}^{-2k-1}(1-p')^{-2k-1} \geq \frac{5}{8}\,.$$%
\end{itemize}

Using a similar argument for $x_2$, to summarize, there exist
 $$ B_1 \subseteq  D_1 = \pi_1^{-1}( \pi_1 ( J_{x_1}  ) ) \cup  \pi_1^{-1} ( \pi_2  ( J_{x_2} ) )
     , \ \  B_2 \subseteq D_2  = \pi_2^{-1}( \pi_1 ( J_{x_1}  ) ) \cup  \pi_2^{-1} ( \pi_2  ( J_{x_2} ) )   $$
such that at least $\frac{5}{8}$ of their extensions outside
$D_1$ and $D_2$ %
%
are in ${\cal F}_{x_1}$ and ${\cal F}_{x_2}$, respectively.  %
Note that

 \begin{equation} \label{eq:B}
 \pi_1(B_1) \cap \pi_2(B_2) = \emptyset\,.%
 \end{equation}
 We are almost done.
 Denote
  $$ {\cal F}_1 = \{ S_1 \subseteq  \Sigma_1 \setminus D_1 \mid B_1 \cup S_1 \in {\cal F}_{x_1}\}, \quad
     {\cal F}_2 = \{ S_2 \subseteq  \Sigma_2 \setminus D_2 \mid B_2 \cup S_2 \in {\cal F}_{x_2}\}\,, $$%
so that $\mu_{p'} ({\cal F}_1) \geq \frac{5}{8}$ and due to monotonicity, letting $q = 1-\frac{1}{\sqrt{2}} \geq p'$, $\mu_{q} ({\cal F}_1) \geq \frac{5}{8}$,
and similarly    $\mu_{q} ({\cal F}_2) \geq \frac{5}{8}$.  Applying \expref{Lemma}{lem:aux-2} to $ {\cal F}_1, {\cal F}_2$ along with $2$-to-$1$
 maps $\pi_1: \Sigma_1 \setminus D_1 \to \Gamma
\setminus \pi_1(D_1)$  and  $\pi_2: \Sigma_2 \setminus D_2 \to \Gamma
\setminus \pi_2(D_2)$ (we have $\pi_1 (D_1) = \pi_2(D_2) = J_{x_1} \cup J_{x_2}$),
there exist $F_1^* \subseteq  \Sigma_1 \setminus D_1, ~F_2^* \subseteq  \Sigma_2 \setminus D_2$ such that
\begin{equation} \label{eq:F-star}
 \pi_1 (F_1^*)  \cap \pi_2 (F_2^*) = \emptyset\,.%
 \end{equation}
 Finally, letting $F_1 = B_1 \cup  F_1^*$  and  $F_2 = B_2 \cup  F_2^*$, and using
\expeqref{Equations}{eq:B} \expeqref{and}{eq:F-star}, %
we conclude that
 $F_1 \in {\cal F}_{x_1}, ~F_2 \in {\cal F}_{x_2}, ~\pi_1(F_1) \cap \pi_2(F_2) = \emptyset$ as desired.
 \end{proof}



Finally, we handle the $1$-to-$1$ constraints inside $X'$. Let $G_{X'}$ be the subgraph of
$G(V,E,\Phi, \Sigma)$ induced on $X'$.  The
transitivity of the game, as per \expref{Definition}{def:trans_game},
then implies that $G|_{X'}$ can be partitioned into
cliques ${\cal C}_1,\ldots,{\cal C}_m$ such that
\begin{itemize}
\item All constraints inside a clique ${\cal C}_r$ are $1$-to-$1$, \ie, matchings on $\Sigma \times \Sigma$.
      For any $x,y,z \in {\cal C}_r$,
      the matchings between $(x,y), (y,z)$ can be composed to derive the matching between $(x,z)$.
\item For $r\neq s$, either there is no edge between ${\cal C}_r, {\cal C}_s$ or there is
a complete bipartite graph between ${\cal C}_r, {\cal C}_s$ with all constraints being $2$-to-$2$.
In the latter case, for any $x, y \in {\cal C}_r$ and $z \in {\cal C}_s$, the $2$-to-$2$ constraint $(x,z)$ is
a composition of the $1$-to-$1$ constraint $(x,y)$ and the $2$-to-$2$ constraint $(y,z)$.
\end{itemize}
These considerations show that all vertices inside a clique $C_r$ play an essentially identical role. Therefore,
in a maximal independent set $I$, for any $x, y \in {\cal C}_r$,
the families ${\cal F}_x, {\cal F}_y \subseteq {\cal P}(\Sigma)$ are identical up to the permutation of
$\Sigma$ that defines the $1$-to-$1$ constraint $(x,y)$, and hence the color-sets $EJ(x), EJ(y)$ are
identical up to the (same) permutation. This shows that $1$-to-$1$ constraints are satisfied in the
sense of \expref{Definition}{def:j,delta}.







\section{\texorpdfstring{``$\ell$-space vs $b$-space''}{``l-space vs b-space''}
  Linearity Test}\label{appendix:linearity_testing}
In this section, we present and analyze ``$\ell$-space vs $b$-space'' linearity test. The analysis is Fourier
analytic and, as is standard, it is convenient to think of boolean values as $\cube{}$ and replace addition
over ${\mathbb F}_2$ by product of the signs $\cube{}$.
A function $f: \Omega = \cube{n}
\to \cube{}$ is linear if $f(x) f(y) = f(x\cdot y)$ for all $x, y \in \Omega$ and $x\cdot y$ denotes the coordinatewise
product of $x, y$.

\begin{center} {\bf  The $\ell$-space vs $b$-space Linearity Test }
\end{center}
For  $\Omega = \cube{n}$, let ${\cal B}$  and ${\cal L}$ denote the set of all $b$-dimensional and $\ell$-dimensional
subspaces of $\Omega$. Let %
$A$ and $F$  %
be tables that assign, for
%
$B \in {\cal B}$ and $L \in {\cal L}$, %
%
linear functions %
$A[B]:B \to \cube{}$ and $F[L]:L \to \cube{}$  %
on the respective subspaces.  %
 The test picks a pair
$(B,L)$ uniformly at random with $B \subseteq L, B \in {\cal B}, L \in {\cal L}$
and accepts
if %
%
%
 $$  F[L]|_B \equiv  A[B]\,.$$%
Our result is the following:
\begin{theorem}\label{thm:linearity_dec}
Let $\Omega, {\cal B}, {\cal L}$ and parameters $n, \ell, 1 \leq b \leq \frac{\ell}{4}$
be as in the description of
the test above. Let %
$A$ and $F$  %
be tables that assign linear functions to $B \in {\cal B}$ and
$L \in {\cal L}$, respectively.  %
Suppose the tables pass the linearity test with probability at least
$\frac{1}{2^b}+\eps$ where
$\eps \geq 2^{2-b/4}$, \ie,
\[
\Prob{B\subseteq  L, B \in {\cal B}, L \in {\cal L}}{F[L]|_B\equiv A[B]}\geqslant \frac{1}{2^b}+\eps\,.%
\]
Then there exists a global linear function $g: \Omega \to \cube{} $ that agrees with at least $\frac{\eps^3}{300}$ fraction of the $\ell$-spaces,
that is
\[
\Prob{L\in {\cal L} }{F[L]\equiv g|_L}\geqslant \frac{\eps^3}{300}\,.%
\]
\end{theorem}


The rest of the section is devoted to proving \expref{Theorem}{thm:linearity_dec}. We start by viewing the entire table $A[\cdot]$ as
a function $f: \cube{nb} \to \cube{b}$ as follows. In notation, for $(v_1, \ldots, v_b) \in (\cube{n})^b = \cube{nb}$,
 $$ f(v_1,\ldots,v_b) = ( A[{\sf Span}(v_1,\ldots,v_b)] (v_1), \ldots, A[{\sf Span}(v_1,\ldots,v_b)] (v_b) )\,. $$%
In words, to evaluate  $f(v_1, \ldots, v_b)$, one considers the $b$-space $B= {\sf Span}(v_1,\ldots,v_b)$, and the
linear function $A[B]$ on $B$. The linear function assigns, in particular, $\cube{}$-values to the vectors
$v_1, \ldots, v_b$.  %
%
The list of these $b$ values is defined to be $f(v_1,\ldots,v_b)$. Since the output
of $f$ is a string of length $b$, we can think of $f$ as a collection of $\cube{}$-valued functions,
$f_1,\ldots,f_b$, one for each output coordinate. In notation,  $f_i: \cube{nb} \to \cube{}$ is defined as
  $$  f_i(v_1,\ldots,v_b) =  A[{\sf Span}(v_1,\ldots,v_b)] (v_i)\,. $$%
We must make
a couple of clarifying remarks.
First, when the input vectors $\set{v_1,...,v_b}$ are linearly dependent, then their span $B$ has dimension less than $b$ and $A[B]$
is undefined. However the fraction of such inputs is negligible (at most $2^{b-n}$) and on those inputs $f$ can be defined arbitrarily
without affecting the analysis. Second, since the same $b$-space may have different bases,
$f$ has many symmetries, \eg, $f_1(v_1,...,v_b)=f_2(v_2,v_1,...,v_b)$. We will use these symmetries, but not in any explicit manner.

%
%
%
%
%
\subsection{The Gowers Test}  %
%
The main idea behind the analysis is to use a %
%
``Gowers Test'' as an auxiliary tool. We can relate the acceptance probability
of the $\ell$-space vs $b$-space test to that of the acceptance probability of
the Gowers Test. The Gowers Test allows us
to conveniently switch from local considerations to global considerations.
Let $\vec{1}$ denote a $b$-dimensional vector with all coordinates $1$.

\begin{definition}{\bf [The Gowers Test]} Given $h: \cube{nb} \to \cube{b}$,
pick $x,y,z\in  \cube{nb}$ randomly and check if
$$h(x)h(y)h(z)h(x\cdot y \cdot z)=\vec{1}\,. $$%
\end{definition}

Represent a function $h: \cube{nb} \to \cube{b}$ as $h = (h_1,\ldots,h_b)$ where $h_i$ are the
coordinatewise functions. For $T \subseteq [b]$, let $h_T = \prod_{i\in T} h_i$ be the product functions.
The lemma below expresses the probability of $h$ passing the Gowers Test
in terms of the Fourier coefficients of products of functions $h_T$.




\begin{lemma}\label{lemma:gowers_to_fourier}
 The probability that  $h\colon \cube{nb} \to \cube{b}$ passes the Gowers Test is:
\[
\Prob{x,y,z\in \cube{nb}}{h(x)h(y)h(z)h(x\cdot y \cdot z)=\vec{1}}=
\frac{1}{2^b} + \frac{1}{2^b} \sum\limits_{T\subseteq [b], T\neq \emptyset}{\sum\limits_{S \subseteq [nb]}{{\widehat{h}_T}^4(S)}}\,.%
\]
\end{lemma}
\begin{proof}
For the test to pass, it must pass on every coordinate. Thus,
%
%
%
%
%
%
%
%
%
%
\begin{align*}
 & \Prob{x,y,z \in \cube{nb}}{h(x)h(y)h(z)h(x\cdot y\cdot z)=\vec{1}} \\
& = \Expect{x,y,z \in \cube{nb}}{  \prod_{i=1}^b  \frac{1+h_i(x)h_i(y)h_i(z)h_i(x\cdot y\cdot z)}{2}}  \\
& = \frac{1}{2^b} + \frac{1}{2^b} {\sum\limits_{T\subseteq [b], T \neq \emptyset}{\Expect{x,y,z\in \cube{nb}}{h_T(x)h_T(y)h_T(z)h_T(x\cdot y\cdot z)}}}  \\
& =  \frac{1}{2^b} + \frac{1}{2^b} \sum\limits_{T\subseteq [b], T \neq \emptyset}{\sum\limits_{S \subseteq [nb]}{{\widehat{h}_T}^4(S)}}
\,.    \qedhere %
\end{align*}
\end{proof}

The main trick is that \expref{Lemma}{lemma:gowers_to_fourier} is applied globally as well
as locally and then the information gained from the
two applications is combined.
Globally, the lemma is applied to the
function $f: \cube{nb} \to \cube{b}$ that (essentially) represents the entire assignment $\{A[B] | B \in {\cal B}\}$. Locally, for
a fixed $\ell$-space $L$,
the lemma is applied to the function $g: \cube{\ell b} \to \cube{b}$
that represents, in a similar manner, the assignment  $\{A[B] | B \subseteq L\}$ (\ie, only the assignment to $b$-spaces
that are contained in $L$). We present the local application first.

Fix an $\ell$-space $L$. Locally, $L$ can be identified with $\cube{\ell}$ and the linear function $F[L]$ on it can be
identified with a Fourier character $\chi_S$ for some $S \subseteq [\ell]$. The assignment  $\{A[B] | B \subseteq L\}$
can be represented, in a similar manner as before, by a function $g: \cube{\ell b} \to \cube{b}$, $g = (g_1,\ldots,g_b)$
where for $(w_1, \ldots, w_b) \in (\cube{\ell})^b = \cube{\ell b}$,
 $$ g_i (w_1, \ldots, w_b)  =  A[{\sf Span} (w_1, \ldots, w_b)  ] (w_i)\,.  $$%
We note that $g$ really is the restriction of $f$ to $L^b$.
As before, for $T \subseteq [b]$, let $g_T = \prod_{i\in T} g_i$ be the product functions.
We now relate the probability that $g$ passes the Gowers Test with the probability that the linearity test passes for
the fixed $L$, \ie, the probability that $F[L]|_B = A[B]$  for a random $B \subseteq L$. Let $1-\gamma$ be the probability
that random vectors $w_1, \ldots, w_b \in \cube{\ell}$ are linearly independent, so that $\gamma \leq 2^{b-\ell}$ is
negligible. Thus choosing a random $b$-dimensional subspace of $L$ is essentially same as choosing $b$ random vectors from $L = \cube{\ell}$.
We now have
\begin{eqnarray*}
(1-\gamma) \cdot \Prob{B \subseteq L}{F[L]|_B = A[B]} & \leq &
    \Prob{w_1,\ldots, w_b \in \cube{\ell}}{ \wedge_{i=1}^b   ~F[L](w_i) = A[{\sf Span}(w_1,\ldots, w_b)](w_i) } \\
    & = &   \Prob{w_1,\ldots, w_b \in \cube{\ell}}{ \wedge_{i=1}^b   ~\chi_S(w_i) = g_i(w_1,\ldots, w_b) }\\
    & = &  \Expect{w_1,\ldots, w_b \in \cube{\ell}} {\prod_{i=1}^b \frac{1+\chi_S(w_i) g_i(w_1,\ldots, w_b) }{2} } \\
    & = &  \frac{1}{2^b} + \frac{1}{2^b} \sum_{T \subseteq [b], T \neq  \emptyset } \Expect{w_1,\ldots, w_b \in \cube{\ell}}{ g_T (w_1,\ldots, w_b) \prod_{i \in T}\chi_S(w_i)  }\\
     & = &  \frac{1}{2^b} + \frac{1}{2^b} \sum_{T \subseteq [b], T \neq  \emptyset }  \widehat{g}_T (S_T)\,,%
\end{eqnarray*}
where $S_T \subseteq [\ell b]$ is defined as $(S_T(1), \ldots, S_T(b))$  and  $S_T(i)
\subseteq [\ell]$ equals $S$ if $i \in T$ and equals $\emptyset$ if $i \not\in T$. Hence noting
that $\gamma \leq 2^{b-\ell}$,
 $$ \Prob{B \subseteq L}{F[L]|_B = A[B]}  ~\leq   ~ \frac{1}{2^b} + 2 \cdot 2^{b-\ell}  + \frac{1}{2^b} \sum_{T \subseteq [b], T \neq  \emptyset }  \widehat{g}_T (S_T)\,. $$%
Now we take average of this inequality over the choice of $L \in {\cal L}$ and note that the L.H.S. then equals the probability that the linearity test accepts (which is
~$\geq \frac{1}{2^b}+\eps$). This gives
 $$ \frac{\eps}{2} \leq   ~\eps - 2 \cdot 2^{b-\ell}  ~\leq  \Expect{L \in {\cal L}} { \frac{1}{2^b} \sum_{T \subseteq [b], T \neq  \emptyset }  \widehat{g}_T (S_T) }\,. $$%
We keep in mind that $g$ and $S$ depend on the choice of $L$. Using convexity of the function $x \to x^4$, we get
 $$ \frac{\eps^4}{16} \leq    \Expect{L \in {\cal L}} {  \frac{1}{2^b} \sum_{T \subseteq [b], T \neq  \emptyset }  \widehat{g}^4_T (S_T) }\,.$$%
Applying \expref{Lemma}{lemma:gowers_to_fourier} to $g: \cube{b\ell} \to \cube{b}$, we get
 $$ \frac{\eps^4}{16} \leq     ~\Expect{L \in {\cal L}}{ \Prob{}{g ~\mbox{passes Gowers Test}  }}\,. $$%
Now we relate the R.H.S. to the probability that $f$ passes the Gowers Test, using the fact that $g$ really is the restriction of
$f$ to $L^b$. Let $x = (x_1,\ldots,x_b), y = (y_1,\ldots,y_b), z= (z_1,\ldots,z_b)$ where $x_i, y_i, z_i$ are either in $L$ or in the global space $\cube{n}$,
as understood from the context. We would like to argue as
\begin{eqnarray*}
  \frac{\eps^4}{16} & \leq &    \Expect{L \in {\cal L}}{ \Prob{}{g ~\mbox{passes
the Gowers Test}  }}\\
 & = &   \Prob{L \in {\cal L}, ~x_i,y_i,z_i \in L}{g(x)g(y)g(z)g(x\cdot y\cdot z) = \vec{1} }\\
 & \approx & \Prob{x_i,y_i,z_i \in \cube{n}}{f(x)f(y)f(z)f(x\cdot y\cdot z) = \vec{1}} \\
 & = & \Prob{}{f ~\mbox{passes the Gowers Test}}\,.%
 \end{eqnarray*}
This is an almost correct argument, except that the distribution ${\cal D}$ of $x_i, y_i, z_i \in \cube{n}$ is slightly different from
the distribution ${\cal D}'$ of $L \in {\cal L}, ~x_i, y_i, z_i \in L = \cube{\ell}$  (\ie, first choosing $L$ at random and then
choosing $x_i, y_i, z_i$ from inside $L$). The distributions are identical however if conditioned on the $3b$ vectors $x_i, y_i, z_i$
being linearly independent. The probability of this happening is at least $1-2^{3b-n}$ and $1-2^{3b-\ell}$ depending on the space they
are chosen from. It follows that the statistical distance between the distributions is at most $3\cdot 2^{3b-\ell}$ and the argument
above is correct up to that much error. It follows that (provided $\eps \geq 4\cdot 2^{-b/4}$)
 $$  \frac{\eps^4}{24} \leq \Prob{}{f ~\mbox{passes the Gowers Test}}\,. $$%
Applying \expref{Lemma}{lemma:gowers_to_fourier} to $f$,
 $$ \frac{1}{2^b} + \frac{1}{2^b} \sum\limits_{T\subseteq [b], T \neq \emptyset}{\sum\limits_{S \subseteq [nb]}{{\widehat{f}_T}^4(S)}}
 \geq  \frac{\eps^4}{24}\,. $$%
Noting that the sum of squares of Fourier coefficients of a boolean function equals $1$, we see that there exists $T \subseteq [b],
T \neq \emptyset$ and $S \subseteq [nb]$ such that
${\widehat{f}_T}^2(S)  \geq \frac{\eps^4}{32}$. We are almost done, by inspecting the coefficient  $\widehat{f}_T(S)$. Let $S = (S_1,\ldots, S_b)$,
$S_i \subseteq [n]$, and  denote $B = {\sf Span}(v_1,\ldots,v_b)$  below.
 By definition of Fourier coefficients and of the functions $f, f_T$,
\begin{eqnarray*}
   \widehat{f}_T(S) & = & \Expect{v_1,\ldots,v_b \in \cube{n}}{ f_T(v_1,\ldots,v_b) \cdot \prod_{i=1}^b \chi_{S_i}(v_i) } \\
    & = &  \Expect{v_1,\ldots,v_b \in \cube{n}}{ \prod_{i\in T} A[B](v_i)  \cdot \prod_{i=1}^b \chi_{S_i}(v_i) } \\
    & = & \Expect{B, ~{\sf dim}(B)=b, ~v_1,\ldots, v_b \in B,  \atop{{\sf Rank}(v_1,\ldots,v_b)=b}  }{ \prod_{i\in T} A[B](v_i)  \cdot \prod_{i=1}^b \chi_{S_i}(v_i)}
     \pm 2^{b-n}\,,%
\end{eqnarray*}
where while choosing $v_1,\ldots,v_b \in \cube{n}$, they are assumed to be linearly independent (introducing the negligible error term $2^{b-n}$) and then
their choice is same as first choosing a random $b$-space $B$ and then letting  $v_1,\ldots,v_b$ be a random basis of $B$. Regard $B = \cube{b}$ and $A[B]$ as
the linear function $\chi_{S[B]}$ for $S[B] \subseteq [b]$. The global function $\chi_{S_i} (v_i)$ where $S_i \subseteq [n], v_i \in \cube{n}$, after restricting
to $v_i \in B$, amounts to a linear function on $B$, say $\chi_{S_i \downarrow B}$ with  $S_i \downarrow B \subseteq [b]$. Thus
$$
   \widehat{f}_T(S)  =  \Expect{B, ~{\sf dim}(B)=b, ~v_1,\ldots, v_b \in B,  \atop{{\sf Rank}(v_1,\ldots,v_b)=b}  }
                                 { \prod_{i\in T}  \chi_{S[B] \Delta S_i\downarrow B} (v_i)   \cdot \prod_{i\in [b]\setminus T}  \chi_{S_i\downarrow B}(v_i)}
     \pm 2^{b-n}\,.%
$$
Let us look at the expectation for a fixed $B$. Call $B$ good if
\begin{equation}  \label{eqn:good-B}
 \forall  ~i\in T,  \  S_i\downarrow B = S[B], \quad \quad \forall ~i\in [b]\setminus T, \  S_i\downarrow B=\emptyset\,,%
\end{equation}
and let ${\cal B}'$ be the set of such good $B$.
For a good $B$, the expectation equals $1$ and from \expref{Lemma}{lemma:basis_exp},
the expectation is bounded by
$2^{-b+2}$
%
in magnitude for a bad $B$. Thus
$$  \widehat{f}_T(S)  =  \Prob{B}{B \in {\cal B}'}  \pm 2^{-b+2} \pm 2^{b-n}\,. $$
 %
Since $|  \widehat{f}_T(S)| \geq \frac{\eps^2}{6}$, it follows that $\Prob{B}{B \in {\cal B}'} \geq \frac{\eps^2}{10}$ (since $\eps \geq 2^{2-b/4}\geqslant 2^{3-b/2}$). %
%
Now we show that in fact for some $S^* \subseteq [n]$, for all $i \in T$, $S_i = S^*$ and for all $i \in [b]\setminus T$, $S_i = \emptyset$. This is
because if this were not the case, for a random $b$-space $B$, \expeqref{Condition}{eqn:good-B} holds with probability at most $2^{-b}$, upper
bounding  $\Prob{B}{B \in {\cal B}' }$  by $2^{-b}$, a contradiction. It follows that $\chi_{S^*}: \cube{n}\to \cube{}$ is a global linear function
that agrees with the given linear function $A[B]$ on $\geq \frac{\eps^2}{10}$ fraction of $B, ~{\sf dim}(B)=b$.

%
\subsection{Agreement with \texorpdfstring{$\ell$-spaces}{l-spaces}}
What we have concluded so far is that if tables $F, A$ pass
the $\ell$-space vs $b$-space linearity test with probability $\geq \frac{1}{2^b} + \eps$, then there is a
global linear function $g: \cube{n} \to \cube{n}$ that agrees with $A[B]$ for $\geq  \frac{\eps^2}{10}$ fraction of $b$-spaces $B$. \expref{Theorem}{thm:linearity_dec} however demands a good agreement with $F[L]$ for $\ell$-spaces $L$. This is easy to fix. Let

 $$ {\cal B}^* =
 \left\{ B \mid \Prob{L: B \subseteq L}{F[L]|_B = A[B]}  \geq \frac{\eps}{2} \right\}\,. $$%

Since the linearity test succeeds with probability $\geq \frac{1}{2^b} + \eps$, by an averaging argument, $|{\cal B}^*| \geq \frac{\eps}{2} \cdot |{\cal B}|$.
Now modify the table $A[\cdot]$ to table $A'[\cdot]$ so that $A'[B] = A[B]$ for $B \in {\cal B}^*$ and $A'[B]$ is a random linear function on $B$ otherwise.
Clearly, the tables $F, A'$ still pass the linearity test with probability $\geq 2^{-b} + \frac{\eps}{2}$ and by the analysis so far, there is a global linear
function $g$ that agrees with $A'[B]$ for $\geqslant \frac{\eps^2}{40}$ fraction of $B \in {\cal B}$. Since $A'[B]$ for $B \not\in {\cal B}^*$ was defined
at random, their contribution to consistency with $g$ is negligible, \ie, at most $2^{-b}$. Thus we have
 $$ \Prob{B}{ g|_B = A[B]  ~\wedge ~B \in B^*    }  \geq  \frac{\eps^2}{80}\,. $$%
For every $B \in {\cal B}^*$, by definition, $A[B]$ is consistent with $F[L]$ for at least $\frac{\eps}{2}$ fraction of $L$ containing $B$.
Hence,
   $$ \Prob{B \subseteq L}{ g|_B = A[B]  ~\wedge ~B \in B^*  ~\wedge F[L]_B = A[B]  }  \geq  \frac{\eps^3}{160}\,. $$%
In particular,
   $$ \Expect{L} {  \Prob{B \subseteq L} {  F[L]|_B = g|B   }          } \geq  \frac{\eps^3}{160}\,. $$%
This implies immediately that $g|_L = F[L]$ for at least  $ \frac{\eps^3}{160} - 2^b \geq  \frac{\eps^3}{300} $ fraction of $L$, since for $L$ not satisfying this,
the inside probability is at most $2^{-b}$.



\section{Missing proofs}
\subsection{ \expref{Hypothesis}{conj:main2} implies \expref{Hypothesis}{conj:main}}\label{appendix:hyp_side_cond}

Let $G(V,\ell)$, ${\sf dim}(V)=n$ and the side condition $\set{h_i}_{i=1}^r, \set{b_i}_{i=1}^r, ~r \leq \frac{n}{3}$ be as in \expref{Hypothesis}{conj:main}. Let
$H = {\sf Span}(h_1,\ldots,h_r), ~{\sf dim}(H)=r$. Let $W[H]$ be any ``complementing space'' to $H$ so that
 $V = H \xor W[H], \ H \cap W[H] = \{0\}$.
We identify linear
functions on subspaces of $V$ that respect the side condition $H$  with  their restrictions to $W[H]$, move to the ``lower order'' Grassmann graph $G(W[H], \ell)$,
apply \expref{Hypothesis}{conj:main2}, and then ``pull-back'' the function on $W[H]$ guaranteed by \expref{Hypothesis}{conj:main2} to a function on $V$
that respects the side condition.

Formally, let $q,\alpha(\cdot),C$ be as in \expref{Hypothesis}{conj:main2} given $(j, \frac{\delta}{2})$ and $n$ sufficiently large.
Let $Q \subseteq V$ be a random $q$-dimensional space. With probability at least $1-2^{r+q-n}$ (see \expref{Fact}{fact:aux_2}), we have $Q \cap H = \{0\}$ and we condition
on this event henceforth.
\begin{claim}
$W[Q] \defeq (Q\xor H)\cap W[H]$ is a random $q$-dimensional subspace of $W[H]$.
\end{claim}
\begin{proof} Firstly, the dimension consideration shows that
\[
\dimm(W[Q]) = \dimm(W[H])+\dimm(Q\xor H) - \dimm(W[H]\xor Q\xor H) = (n-r) + (q+r) - n = q\,.%
\]
Also, it is easily seen that each $q$-dimensional subspace of $W[H]$ has equally many
pre-images under the mapping $Q\to (Q\xor H)\cap W[H]$.
\end{proof}
Let $F[\cdot]$ be the $(j,\delta)$-assignment to $G(V,\ell)$ respecting the side condition $H$ and $S$ be the set of its vertices that have been assigned.
 We ``move'' to the lower order Grassmann graph $G(W[H], \ell)$ and define a
$(j,\delta)$-assignment $\tilde{F}[\cdot]$  to it as (denoting the set of its vertices assigned as $\tilde{S}$)
\[
\tilde{S} = \sett{L\in G(W[H],\ell)}{L\xor H\in S}\,.%
\]
\[
\tilde{F}[L] = F[L\xor H]|_{L}\,.%
\]
By \expref{Hypothesis}{conj:main}, with probability at least $\alpha(\ell)$ over the choice of $Q$, there exists
$g_{Q}\colon W[H]\to\power{}$ such that
\[
  \Prob{L: Q\subseteq L \subseteq W[H]}{g_{Q}|_{L}\in \tilde{F}[L]}\geqslant C\,.%
\]
Define $g'_Q$ to be the unique extension of $g_{Q}$ to $V$ respecting the side condition. Since spaces $L \subseteq W[H]$ can be
pulled back to $L \xor H$, $g'_Q$ satisfies \expeqref{Equation}{eqn:global-fn-side} of \expref{Hypothesis}{conj:main} as required.

\subsection{Proof of \expref{Theorem}{thm:list_dec_upper}}\label{appendix:list_dec_upper}
Denote by $N$ the size of ${\cal L} \defeq \sett{L\in G(V,\ell)}{Q\subseteq L}$
and let $f_1,...,f_m$ be all functions
agreeing with $F[\cdot]$ on at least $C$ fraction of $L \in {\cal L}$. We construct a bipartite graph, where the left side
consists of $f_1,...,f_m$ and the right side consists of pairs $\sett{(L,\sigma)}{L\in {\cal L}, \sigma\in F[L]}$. We connect $f_i$
and $(L,\sigma)$ by an edge if ${f_i}|_L\equiv \sigma$. Then the degree of each $f_i$ is at least $C\cdot N$ and the number
of vertices on the right side is at most $j N$. Let us remove edges if necessary so that the degree of each
$f_i$ is exactly $C\cdot N$.

Denote by $d(L,\sigma)$ the degree of $(L,\sigma)$
and let us count the %
%
%
%
triples  $\{f_i,f_j, (L,\sigma)\}$
where $i\neq j$ and $(f_i, (L,\sigma))$, $(f_j, (L,\sigma))$ are
both edges in the bipartite graph. Using %
Cauchy--Schwarz %
and noting that the number of vertices on the right side is at most $j N$ and
$\sum_{L \in {\cal L}, \sigma \in F[L]}{d(L,\sigma)} = CmN$,
the number of such triples is %
at least  %
$$ \sum\limits_{L\in {\cal L}, \atop{\sigma\in F[L]}}{{d(L,\sigma) \choose 2}} = \sum\limits_{L\in {\cal L}, \atop{\sigma\in F[L]}} \frac{d(L,\sigma)^2}{2} - \frac{d(L,\sigma)}{2}
 \ \geq \ j N \cdot \frac{\left(\frac{CmN}{jN}\right)^2}{2} - \frac{CmN}{2} \  = \frac{C^2 m^2 N}{2j} - \frac{CmN}{2}\,.$$%

On the other hand, since any distinct pair of functions $f_i,f_j$ agree on at most $2^{q-\ell}$ fraction of $L \in {\cal L}$,
the number of such triples is at most $
{m \choose 2} 2^{q-\ell} N \leq \frac{ m^2 2^{q-\ell} N}{2}.  $ Combining the two bounds gives $m \leq \frac{jC}{C^2 - j 2^{q-\ell}}$.


\subsection{Auxiliary lemmas and facts}

%
%
\begin{claim}\label{claim:gen_claim}
Let $X$ be a real random variable such that $\card{X}\leq 1$ always, and let $B$ be an event.
Then $\left|\Expect{}{X}-\cExpect{}{B}{X}\right|\leq 2\Prob{}{\bar{B}}$.
\end{claim}
\begin{proof}
Let $Z = 1_B$ be the random variable indicating if $B$ occurred. Then by the law of 
total expectation we have
\[
\Expect{}{X} = \Expect{Z}{\cExpect{}{Z}{X}} = \Prob{}{B} \cExpect{}{B}{X} + \Prob{}{\bar{B}}\cExpect{}{\bar{B}}{X}.
\]
Thus $\Expect{}{X}-\cExpect{}{B}{X} = \Prob{}{\bar{B}}(\cExpect{}{\bar{B}}{X}-\cExpect{}{B}{X})$,
and the claim follows.
\end{proof}

\begin{lemma}\label{lemma:basis_exp}
Let $s_1, \ldots, s_b \in {\mathbb F}_2^b$ such that at least one of them is non-zero. Let $v_1, \ldots, v_b \in {\mathbb F}_2^b$ be chosen
at random. Then the following conditional expectation is bounded as:
$$ \left| \Expect{v_1,\ldots,v_b}{\prod_{i=1}^b (-1)^{\left<s_i, v_i \right> }
\mid {\sf Rank}(v_1,\ldots,v_b)=b }\right|  \leq
  2^{-b+2}\,.$$%
\end{lemma}
\begin{proof} Note that without the conditioning, the expectation is clearly zero. The point is to prove the upper bound conditional on the event that
$v_1, \ldots, v_b$ are linearly independent
(and hence form a basis of ${\mathbb F}_2^b\,$). %
Assume w.l.o.g. that $s_1$ is non-zero. Let
 $$ {\cal A}  = \{ A =  (v_1, \ldots, v_b) \mid {\sf Rank} (A) = b \}\,,$$%
so that we are interested in the expectation
 $$   \Expect{A}{\prod_{i=1}^b (-1)^{\left<s_i, v_i \right> } \mid A \in {\cal A} }\,.  $$%
 Let
 $$ {\cal A'}  = \{ A =  (v_1, \ldots, v_b) \mid {\sf Rank} (A) = b, ~\forall ~2 \leq i \leq b, ~\left<  s_1, v_i \right> = 0 \}\,.$$%
It is easily seen that $|{\cal A}'| \leq 2^{-b+1} \cdot |{\cal A}|$.
Indeed, %
imagine choosing $v_2, v_3, \ldots, v_b$ so that every $v_i$ is outside
the span of the previously chosen ones. If we require (in addition) that every $v_i$ also lies in
the hyperplane defined by the equation
$\left< s_1, x \right> = 0$, then at each step,
this happens with probability at most $\frac{1}{2}$, showing the desired upper bound on
$|{\cal A}'|$.
It follows from~\expref{Claim}{claim:gen_claim} that
%
$$ \Expect{A}{\prod_{i=1}^b (-1)^{\left<s_i, v_i \right> } \mid A \in {\cal A} }, \quad  \quad
\Expect{A}{\prod_{i=1}^b (-1)^{\left<s_i, v_i \right> } \mid A \in {\cal A}\setminus {\cal A}' } $$
differ by at most $2^{-b+2}$.  %
%
We show that the latter is zero. For fixed $\alpha_2, \alpha_3,\ldots, \alpha_b \in {\mathbb F}_2$,
consider the following bijection on $ {\cal A}\setminus {\cal A}' $ (that adds to the first vector, a linear combination of others):
  $$  (v_1, v_2, v_3, \ldots, v_b)  \to   (v_1 + \sum_{i=2}^b \alpha_i v_i, ~v_2, v_3, \ldots, v_b)\,. $$%
 The quantity of interest changes as follows:
  $$ \prod_{i=1}^b (-1)^{\left<s_i, v_i \right> }   \to   (-1)^{ \sum_{i=2}^b \alpha_i \left< s_1, v_i\right>} \cdot   \prod_{i=1}^b (-1)^{\left<s_i, v_i \right> }\,. $$%
Now take expectation of L.H.S. over the choice of $A = (v_1,\ldots,v_b)  \in {\cal A}\setminus {\cal A}'$ and expectation of R.H.S. over
the choice of $A  \in {\cal A}\setminus {\cal A}'$
as well as over a random choice of  $\alpha_2, \ldots, \alpha_b$. The two expectations are equal (due to bijectivity)
and the expectation of the L.H.S. is what we are interested in.
Since $ \left< s_1, v_i\right> \neq 0$ for some $2 \leq i \leq b$, the expectation over the R.H.S. is zero and we are done.
\end{proof}

%
\begin{fact}\label{fact:aux_1}
Let $V$ be an $n$-dimensional vector space over $\mathbb{F}_2$,
and $1 \leq \ell \leq  n-1$. Let $x_1, \ldots, x_\ell \in V$ be chosen randomly and
independently. Then $x_1, \ldots, x_\ell$ are
linearly independent with probability $\geq 1-2^{\ell -n}$.
\end{fact}
\begin{proof} If $x_1, \ldots, x_\ell$ are linearly dependent, then for some
$1 \leq i \leq \ell$, $x_i$ is in the span of $x_1, \ldots, x_{i-1}$. Hence
the probability that these $\ell$ vectors are linearly dependent
%
is at most  %
%
\[   %
\sum\limits_{i=1}^{\ell}{\Prob{x_1,...,x_{i}\in V}
{x_i\in {\sf Span}{\set{x_1,...,x_{i-1}}}}} \leq
\sum\limits_{i=1}^{\ell}{\frac{2^{i-1}}{2^n}}\leq
2^{\ell-n}\,.     \qedhere
%
\]   %
\end{proof}

\begin{fact}\label{fact:aux_2}
Let $V$ be an $n$-dimensional vector space over $\mathbb{F}_2$,
$H\subseteq U$ be a subspace of dimension $r$, and $1 \leq \ell \leq  n-\ell$.
Let $x_1, \ldots, x_\ell \in V$ be chosen randomly and independently. Then
\[
\Prob{x_1,...,x_\ell \in V }{\spnn{\set{x_1,...x_\ell}}\cap H =\set{0}}\geq 1-2^{r+\ell-n}\,.%
\]
\end{fact}
\begin{proof} If ${\sf Span}(\{ x_1, \ldots, x_\ell\}) \cap H \neq  \{0\}$, then for some $1 \leq i \leq \ell$, $x_i$ is in the span of $H \cup \{x_1, \ldots, x_{i-1}\}$. Hence
the probability that ${\sf Span}(\{ x_1, \ldots, x_\ell\}) \cap H \neq  \{0\}$
%
is at most  %
%
\[     %
\sum\limits_{i=1}^{\ell}{\Prob{x_1,...,x_\ell \in V}
{{x_i\in H\xor\spnn{\set{x_1,...,x_{i-1}}}}}}
\leqslant\sum\limits_{i=1}^{\ell}{\frac{2^{r+i-1}}{2^{n}}}
\leqslant 2^{r+\ell-n}\,.   \qedhere
%
\]         %
\end{proof}

\begin{fact}\label{fact:aux_3}  %
%
Let  $A,A',B$ be subspaces of a vector space $V$ over ${\mathbb F}_2$ such that
$A\xor B = A'\xor B$ and $(A\xor A')\cap B = \set{0}$. Then $A=A'$.
\end{fact}
\begin{proof}
By symmetry, it suffices to show that $A\subseteq A'$. Let $a\in A$. Then
$a\in A\xor B = A'\xor B$ and so there are $a'\in A'$, $b\in B$ such
that $a = a' \xor b$. Hence $b=a\xor a'\in A\xor A'$, and $b$ must be $0$.
\end{proof}


\section{Acknowledgements}   %
We thank  Prahladh Harsha,  Johan H{\aa}stad, and Sushant Sachdeva  for a counterexample
to an earlier (rather different) construction. We thank
Dana Moshkovitz and Ran Raz for related discussions over the years.  A big thanks to
Irit Dinur and Guy Kindler for discussions that clarified several of
the  %
authors' ideas and led to a better presentation.
We thank
the  %
anonymous referees for helpful comments.
%
%
%
%
%
%
The authors acknowledge the hospitality of the French Roast Bistro
(78 West 11th Street, New York) while working on the paper.
Unfortunately the place closed down since the conference publication of 
this paper~\cite{conf-version}.  %

\bibliographystyle{tocplain}   %

%
%
%
%
\bibliography{bibstrings,v021a010,bibtail}


%
%
%
%
%
\begin{tocauthors}
\begin{tocinfo}[khot]
 Subhash Khot\\
 Professor\\
 Courant Institute of Mathematical Sciences\\
 New York University\\
 NY, USA\\
 khot\tocat{}cs\tocdot{}nyu\tocdot{}edu \\   %
 \url{https://cs.nyu.edu/~khot/}      %
\end{tocinfo}
\begin{tocinfo}[minzer]
  Dor Minzer\\
  Associate Professor\\
  %
  Department of Mathematics\\
  Massachusetts Institute of Technology\\
  MA, USA\\
  dminzer\tocat{}mit\tocdot{}edu \\
  \url{https://sites.google.com/view/dorminzer/home}
\end{tocinfo}
\begin{tocinfo}[safra]
  Muli Safra\\
  Professor\\
  Department of Computer Science\\
  Tel-Aviv University\\
  Tel-Aviv, Israel\\
  safra\tocat{}mail\tocdot{}tau\tocdot{}ac\tocdot{}il\\
  \url{https://sites.google.com/site/mulisafra/home}
\end{tocinfo}
\end{tocauthors}

%
%
%
%
%
%
%
%
%
%
%
%
%
%
%
%
%
%
%
%
\begin{tocaboutauthors}
\begin{tocabout}[khot]  %
Subhash Khot is a Professor in the Computer Science Department at New York University. He has a \phd\ from Princeton University, completed in 2003 under the supervision of Sanjeev Arora.
His teacher and mentor at Vyankatrao Highschool, Mr. Vaman G. Gogate played a decisive role in directing his attention to mathematics. If not for Mr. Gogate’s guidance and some fortuitous turn of events, the chance of someone from the remote town of Ichalkaranji pursuing mathematical research was strictly nil. Subsequently, Mr. Gogate also mentored Amit Deshpande and Raghav Kulkarni (TCS), Abhijit Gadde (Physics), and Dinesh Bharadia (EE). He kept in touch with all his past students and continued to provide guidance on all aspects of life. He passed away in 2020 on Sept 5, which %
coincidentally  %
is celebrated as the %
Teachers' Day in India. %
\end{tocabout}
\begin{tocabout}[minzer]
Dor Minzer is an Associate Professor in the Mathematics Department at Massachusetts Institute of Technology. He has a \phd\ from Tel Aviv University, completed in 2018 under the
guidance of Muli Safra, after which he spent 2 years as a
%
postdoctoral  %
fellow at the Institute for Advanced Study, Princeton.
Dor is interested in Hardness of Approximation,
Analysis of Boolean functions and relations to adjacent areas such as extremal combinatorics and error correcting codes. He started his academic studies %
at  %
the Open University of Israel, %
where  %
he met Professor Vadim Greenstein %
who  %
has greatly influenced Dor and his passion for math.
\end{tocabout}
\begin{tocabout}[safra]
Muli Safra is a Professor of Computer Science at Tel Aviv University. He has a \phd\ from the Weizmann Institute, completed in 1990 under the supervision of Amir Pnueli. He is one of the recipients of the 2001 G\"{o}del Prize..
\end{tocabout}
\end{tocaboutauthors}

\end{document}

